% Lesson 13 — Listening: Texts on comparing prices or quality to determine the best buys for consumer economy and community resources/services — Anglais Tle A — Cours signé OnBuch+
% Programme officiel MINESEC. Compiler avec : tectonic EN13-listening-texts-on-comparing-prices-or-quality-to-deter.tex
\newcommand{\DOCMATIERE}{Anglais}
\newcommand{\DOCNIVEAU}{Tle A}
\newcommand{\DOCMODULE}{Chapter 2 : Economic Life and Occupations}
\newcommand{\DOCLECON}{EN13}
\newcommand{\DOCTITRE}{Lesson 13 — Listening: Texts on comparing prices or quality to determine the best buys for consumer economy and community resources/services}
% OnBuch+ — préambule « cours signés » ENRICHI (compilé avec tectonic / XeLaTeX)
% Système visuel « Pop » d'OnBuch+ : fond crème (jamais blanc), bordures encre
% épaisses, ombres « dures » décalées, titres Archivo Black en capitales.
% Version MATHÉMATIQUES : graphes (pgfplots/tikz), boîtes pédagogiques
% (définition, propriété, méthode, attention, le savais-tu, exercice résolu,
% exercice, TP, bilan), page de garde et signature OnBuch+ ancrée en bas.
%
% Macros attendues dans main.tex AVANT \input{preamble} :
%   \DOCMATIERE  \DOCNIVEAU  \DOCMODULE  \DOCLECON (numéro)  \DOCTITRE
\documentclass[11pt]{article}

\usepackage{fontspec}
\usepackage[english,french]{babel} % français = langue principale ; l'anglais via \eng{..} / textebox
\usepackage[a4paper,margin=2cm,top=3.4cm,bottom=2.8cm,headheight=1.4cm,footskip=1.3cm]{geometry}
\usepackage{amsmath,amssymb,amsfonts}
\usepackage{mathtools} % \xmapsto, \coloneqq... souvent utilisés par les modèles
\usepackage{cancel} % simplifications barrées (bilans de forces)
% \abs{z} (module, valeur absolue) et \norm{u} (norme), utilisés par les modèles
\providecommand{\abs}{}\let\abs\relax\DeclarePairedDelimiter{\abs}{\lvert}{\rvert}
\providecommand{\norm}{}\let\norm\relax\DeclarePairedDelimiter{\norm}{\lVert}{\rVert}
\usepackage[table]{xcolor} % [table] : fournit aussi \rowcolors (alternance de lignes)
\usepackage{pagecolor}
\usepackage{tikz}
\usetikzlibrary{arrows.meta,positioning,calc,shapes.geometric,shapes.misc,shapes.arrows,decorations.pathmorphing,decorations.pathreplacing,decorations.markings,patterns,shadows,mindmap,trees,fit,backgrounds,matrix,intersections,angles,quotes}
\usepackage{pgfplots}
\usepackage[version=4]{mhchem} % équations nucléaires \ce{^{235}_{92}U}
\usepackage[european,siunitx]{circuitikz} % schémas électriques
\pgfplotsset{compat=1.18}
\usepgfplotslibrary{fillbetween,groupplots} % aires entre courbes (calcul intégral)
\usepackage{enumitem}
\usepackage{fancyhdr}
% [most] charge toutes les bibliothèques tcolorbox (listings, minted, raster,
% documentation...) et gonfle inutilement la mémoire TeX sur un document long
% (~300 boîtes) : on ne charge que ce qui est réellement utilisé.
\usepackage[skins,breakable]{tcolorbox}
\usepackage{graphicx}
\usepackage{colortbl}
\usepackage{booktabs}
\usepackage{tabularx}
\usepackage{diagbox} % case d'en-tête diagonale des tableaux à double entrée
% Colonnes extensibles centrée / gauche / droite (C, L, R), souvent utilisées par les modèles
\newcolumntype{C}{>{\centering\arraybackslash}X}
\newcolumntype{L}{>{\raggedright\arraybackslash}X}
\newcolumntype{R}{>{\raggedleft\arraybackslash}X}
\usepackage{array}
\usepackage{multicol}
\usepackage{multirow}
\usepackage{siunitx}
% parse-numbers=false : accepte aussi \SI{2,5 \times 10^{-3}}{...}
\sisetup{locale=FR,output-decimal-marker={,},group-minimum-digits=4,parse-numbers=false}
\DeclareSIUnit\ohm{\ensuremath{\symup{\Omega}}} % U+2126 absent de la police
\DeclareSIUnit\division{div} % réglages d'oscilloscope (ms/div, V/div)
\DeclareSIUnit\tr{tr} % tours (vitesse de rotation, ex. \SI{1500}{\tr\per\minute})
\DeclareSIUnit\molar{M} % concentration molaire (ex. \SI{3}{\molar}, \SI{1}{\micro\molar})
\DeclareSIUnit\an{an} % année (le modèle confond parfois avec \year, un compteur TeX) — ex. \SI{5}{\kilo\meter\per\an}
\DeclareSIUnit\FCFA{FCFA} % franc CFA (ex. \SI{500}{\FCFA\per\kilo\gram})
\DeclareSIUnit\degreeBrix{\textdegree Brix} % teneur en sucre d'un sirop/jus (réfractométrie)
\DeclareSIUnit\month{mois} % le modèle utilise parfois \month, un compteur TeX, comme unité de durée
\usepackage{qrcode}
% \degree : commande fréquemment utilisée par les modèles pour un angle ou une
% température en dehors de \SI{}{} (non fournie par les paquets chargés).
\providecommand{\degree}{^{\circ}}
% \qed (fin de démonstration), utilisé par les modèles sans amsthm
\providecommand{\qed}{\hfill$\square$}
% \barre : double barre des valeurs interdites dans un tableau de variations (tkz-tab)
\providecommand{\barre}{\ensuremath{\|}}
% environnement proof (amsthm non chargé) utilisé par les modèles
\ifdefined\proof\else\newenvironment{proof}[1][Démonstration]{\par\noindent\textit{#1.}\ }{\qed\par}\fi
\usepackage{lastpage}
\usepackage{hyperref}
\hypersetup{hidelinks}

% ============================================================
% Palette « Pop » OnBuch+ — mêmes hex que la classe P (plus_ui.dart)
% ============================================================
\definecolor{poporange}{HTML}{F59321}
\definecolor{popdark}{HTML}{E07A0C}
\definecolor{popcream}{HTML}{FFF6E8}
\definecolor{popink}{HTML}{1C1714}
\definecolor{poppurple}{HTML}{7A5AE0}
\definecolor{popgreen}{HTML}{2E9E6A}
\definecolor{poppink}{HTML}{E0457B}
\definecolor{popblue}{HTML}{2F7DE1}
\definecolor{popgold}{HTML}{C98A12}
\definecolor{popmuted}{HTML}{8A7F76}
\definecolor{popline}{HTML}{EADFD2}
\definecolor{popgray}{HTML}{8A7F76}
\colorlet{popgrey}{popgray}
% Alias tolérants (le modèle nomme parfois les couleurs différemment)
\colorlet{popwhite}{white}
\colorlet{popblack}{popink}
\colorlet{popred}{poppink}
\colorlet{popyellow}{popgold}
% Teintes claires (fonds de tableaux/encadrés) — le modèle nomme parfois
% les couleurs avec un suffixe L (ex. popblueL) sans les avoir définies.
\colorlet{poporangeL}{poporange!20}
\colorlet{popdarkL}{popdark!20}
\colorlet{poppurpleL}{poppurple!20}
\colorlet{popgreenL}{popgreen!20}
\colorlet{poppinkL}{poppink!20}
\colorlet{popblueL}{popblue!20}
\colorlet{popgoldL}{popgold!20}
\colorlet{popredL}{popred!20}
\colorlet{popgrayL}{popgray!20}
% Teintes claires pour fonds de tableaux / zones
\colorlet{popblueL}{popblue!10!white}
\colorlet{poporangeL}{poporange!14!white}
\colorlet{popgreenL}{popgreen!12!white}
\colorlet{poppurpleL}{poppurple!12!white}
\colorlet{poppinkL}{poppink!12!white}
\colorlet{popgoldL}{popgold!14!white}

\pagecolor{popcream}

% ============================================================
% Typographie
% ============================================================
\defaultfontfeatures{Ligatures=TeX}
\setmainfont{PlusJakartaSans-Medium.ttf}[Path=../../../fonts/,BoldFont=PlusJakartaSans-Bold.ttf]
% Symboles, API et emojis absents de la police principale : repli sur DejaVu Sans (caractères actifs)
\newfontface\symfont{DejaVuSans.ttf}[Path=../../../fonts/]
\makeatletter
\newcommand\symactive[1]{\catcode#1=13 \begingroup\lccode`\~=#1 \lowercase{\endgroup\def~}{{\symfont\char#1}}}
\newcommand\symblank[1]{\catcode#1=13 \begingroup\lccode`\~=#1 \lowercase{\endgroup\def~}{}}
\newcommand\symhyph[1]{\catcode#1=13 \begingroup\lccode`\~=#1 \lowercase{\endgroup\def~}{\mbox{-}}}
\newcount\symc
\newcommand\symrange[2]{\symc=#1 \@whilenum\symc<#2 \do{\expandafter\symactive\expandafter{\the\symc}\advance\symc by 1 }}
\newcommand\symblankrange[2]{\symc=#1 \@whilenum\symc<#2 \do{\expandafter\symblank\expandafter{\the\symc}\advance\symc by 1 }}
\makeatother
\symrange{"0250}{"02FF}\symactive{"2713}\symactive{"2714}\symactive{"2717}\symactive{"2718}\symactive{"2610}\symactive{"2611}\symactive{"2612}
\symhyph{"2011}\symhyph{"2E1A}\symblankrange{"1F300}{"1FAFF}\symblank{"FE0F}
% Maths en sans-serif (Fira Math) avec chiffres et lettres droites de Plus
% Jakarta, pour que les formules se fondent dans le texte.
\usepackage[math-style=ISO,bold-style=ISO]{unicode-math}
\setmathfont{FiraMath-Regular.otf}[Path=../../../fonts/]
\setmathfont{PlusJakartaSans-Medium.ttf}[Path=../../../fonts/,range={up/num,up/{Latin,latin},\mathup}]
\usepackage{icomma} % virgule décimale sans espace en mode math
% Caractères mathématiques Unicode absents des polices de texte (Plus Jakarta,
% Archivo Black) : rendus par leur équivalent mathématique, en texte comme en
% formule — sinon ils s'affichent en carré vide (ex. « Arithmétique dans ℤ »).
\usepackage{newunicodechar}
\newunicodechar{ℝ}{\ensuremath{\mathbb{R}}}
\newunicodechar{ℤ}{\ensuremath{\mathbb{Z}}}
\newunicodechar{ℂ}{\ensuremath{\mathbb{C}}}
\newunicodechar{ℕ}{\ensuremath{\mathbb{N}}}
\newunicodechar{ℚ}{\ensuremath{\mathbb{Q}}}
\newunicodechar{𝔻}{\ensuremath{\mathbb{D}}}
\newunicodechar{α}{\ensuremath{\alpha}}
\newunicodechar{β}{\ensuremath{\beta}}
\newunicodechar{γ}{\ensuremath{\gamma}}
\newunicodechar{δ}{\ensuremath{\delta}}
\newunicodechar{ε}{\ensuremath{\varepsilon}}
\newunicodechar{θ}{\ensuremath{\theta}}
\newunicodechar{λ}{\ensuremath{\lambda}}
\newunicodechar{μ}{\ensuremath{\mu}}
\newunicodechar{σ}{\ensuremath{\sigma}}
\newunicodechar{τ}{\ensuremath{\tau}}
\newunicodechar{φ}{\ensuremath{\varphi}}
\newunicodechar{ψ}{\ensuremath{\psi}}
\newunicodechar{ω}{\ensuremath{\omega}}
\newunicodechar{Σ}{\ensuremath{\Sigma}}
% majuscules produites par \MakeUppercase dans les titres (α-aminés -> Α)
\newunicodechar{Α}{\ensuremath{\alpha}}
\newunicodechar{Β}{\ensuremath{\beta}}
\newunicodechar{Γ}{\ensuremath{\Gamma}}
\newunicodechar{Δ}{\ensuremath{\Delta}}
\newunicodechar{Ε}{\ensuremath{\varepsilon}}
\newunicodechar{Θ}{\ensuremath{\Theta}}
\newunicodechar{Λ}{\ensuremath{\Lambda}}
\newunicodechar{Μ}{\ensuremath{\mu}}
\newunicodechar{Τ}{\ensuremath{\tau}}
\newunicodechar{Φ}{\ensuremath{\Phi}}
\newunicodechar{Ψ}{\ensuremath{\Psi}}
\newunicodechar{Ω}{\ensuremath{\Omega}}
\newunicodechar{↦}{\ensuremath{\mapsto}}
\newunicodechar{∈}{\ensuremath{\in}}
\newunicodechar{∉}{\ensuremath{\notin}}
\newunicodechar{∧}{\ensuremath{\wedge}}
\newunicodechar{⇔}{\ensuremath{\Leftrightarrow}}
\newunicodechar{⟺}{\ensuremath{\Longleftrightarrow}}
\newunicodechar{⇒}{\ensuremath{\Rightarrow}}
\newunicodechar{⟹}{\ensuremath{\Longrightarrow}}
\newunicodechar{∘}{\ensuremath{\circ}}
\newunicodechar{∪}{\ensuremath{\cup}}
\newunicodechar{∩}{\ensuremath{\cap}}
\newunicodechar{≡}{\ensuremath{\equiv}}
\newunicodechar{⊂}{\ensuremath{\subset}}
\newunicodechar{⊥}{\ensuremath{\perp}}
\newunicodechar{∃}{\ensuremath{\exists}}
\newunicodechar{∀}{\ensuremath{\forall}}
\newunicodechar{∛}{\ensuremath{\sqrt[3]{\,}}}
\newunicodechar{∜}{\ensuremath{\sqrt[4]{\,}}}
\newunicodechar{⃗}{\ensuremath{{}^{\to}}}
\newunicodechar{ⁿ}{\ifmmode^{n}\else\textsuperscript{\ensuremath{n}}\fi}
\newunicodechar{ˣ}{\ifmmode^{x}\else\textsuperscript{\ensuremath{x}}\fi}
\newunicodechar{ᵇ}{\ifmmode^{b}\else\textsuperscript{\ensuremath{b}}\fi}
\newunicodechar{ʳ}{\ifmmode^{r}\else\textsuperscript{\ensuremath{r}}\fi}
\newunicodechar{ᵘ}{\ifmmode^{u}\else\textsuperscript{\ensuremath{u}}\fi}
\newunicodechar{ʸ}{\ifmmode^{y}\else\textsuperscript{\ensuremath{y}}\fi}
\newunicodechar{ᵃ}{\ifmmode^{a}\else\textsuperscript{\ensuremath{a}}\fi}
\newunicodechar{ᵉ}{\ifmmode^{e}\else\textsuperscript{\ensuremath{e}}\fi}
\newunicodechar{ᵏ}{\ifmmode^{k}\else\textsuperscript{\ensuremath{k}}\fi}
\newunicodechar{ᵖ}{\ifmmode^{p}\else\textsuperscript{\ensuremath{p}}\fi}
\newunicodechar{ᴺ}{\ifmmode^{N}\else\textsuperscript{\ensuremath{N}}\fi}
\newunicodechar{ᶜ}{\ifmmode^{c}\else\textsuperscript{\ensuremath{c}}\fi}
\newunicodechar{ʰ}{\ifmmode^{h}\else\textsuperscript{\ensuremath{h}}\fi}
\newunicodechar{ᵠ}{\ifmmode^{q}\else\textsuperscript{\ensuremath{q}}\fi}
\newunicodechar{ₙ}{\ifmmode_{n}\else\textsubscript{\ensuremath{n}}\fi}
\newunicodechar{ₐ}{\ifmmode_{a}\else\textsubscript{\ensuremath{a}}\fi}
\newunicodechar{ᵢ}{\ifmmode_{i}\else\textsubscript{\ensuremath{i}}\fi}
\newunicodechar{ᵣ}{\ifmmode_{r}\else\textsubscript{\ensuremath{r}}\fi}
\newunicodechar{ₖ}{\ifmmode_{k}\else\textsubscript{\ensuremath{k}}\fi}
\newunicodechar{ᵦ}{\ifmmode_{\beta}\else\textsubscript{\ensuremath{\beta}}\fi}
\newunicodechar{ₓ}{\ifmmode_{x}\else\textsubscript{\ensuremath{x}}\fi}
\newfontfamily\popdisplay{ArchivoBlack-Regular.ttf}[Path=../../../fonts/]
\newfontfamily\poplogo{SpaceGrotesk-Bold.ttf}[Path=../../../fonts/]
\newfontfamily\popsemi{PlusJakartaSans-SemiBold.ttf}[Path=../../../fonts/]
\newfontfamily\popxbold{PlusJakartaSans-ExtraBold.ttf}[Path=../../../fonts/]
\newfontfamily\popxboldls{PlusJakartaSans-ExtraBold.ttf}[Path=../../../fonts/,LetterSpace=3.0]

\setlist[enumerate]{leftmargin=1.7em,itemsep=5pt,topsep=4pt}
\setlist[itemize]{leftmargin=1.7em,itemsep=4pt,topsep=4pt,label={\color{poporange}\textbullet}}
\setlength{\parindent}{0pt}
\setlength{\parskip}{5pt}
\renewcommand{\arraystretch}{1.25}

% pgfplots : style Pop par défaut
\pgfplotsset{
  every axis/.append style={axis line style={popink,line width=1pt},tick style={popink},
    grid style={popline,line width=0.5pt},label style={font=\small},tick label style={font=\footnotesize}},
  popcurve/.style={poporange,line width=1.8pt,no markers,smooth},
}
% Même style disponible dans un \draw TikZ simple (hors pgfplots)
\tikzset{popcurve/.style={poporange,line width=1.8pt}}
\tikzset{popgrid/.style={popline,very thin}}
\colorlet{popcurve}{poporange} % le nom du style est parfois utilisé comme couleur

% ============================================================
% Pill / squiggle
% ============================================================
\newcommand{\poppill}[3]{%
  \tikz[baseline=-0.55ex]\node[draw=popink,line width=1.2pt,rounded corners=2.6pt,%
    fill=#2,inner xsep=6.5pt,inner ysep=3pt]{%
    \popxboldls\fontsize{7}{7}\selectfont\color{#3}\MakeUppercase{#1}};%
}
\newcommand{\popsquiggle}[1]{%
  \tikz[baseline=-0.4ex]{
    \draw[#1,line width=2.6pt,line cap=round]
      plot [smooth] coordinates {(0,0.09) (0.34,0.01) (0.68,0.21) (1.02,0.01) (1.36,0.10)};
  }%
}
% Mot-clé surligné (marqueur orange sous le texte)
\newcommand{\cle}[1]{{\popxbold\color{popdark}#1}}

% ============================================================
% En-tête / pied
% ============================================================
\pagestyle{fancy}
\fancyhf{}
\renewcommand{\headrulewidth}{0pt}
\renewcommand{\footrulewidth}{0pt}
\lhead{\raisebox{-0.15ex}{\poplogo\fontsize{13}{13}\selectfont\color{popink}OnBuch\color{poporange}+}}
\rhead{\poppill{\DOCMATIERE\ \textbullet\ \DOCNIVEAU\ \textbullet\ Leçon \DOCLECON}{white}{popink}}
\lfoot{\popxbold\fontsize{7.6}{7.6}\selectfont\color{popmuted}\MakeUppercase{Cours signé OnBuch+}\ \textbullet\ {\popsemi\DOCTITRE}}
\rfoot{\tikz[baseline=-0.6ex]\node[rounded corners=8.5pt,fill=popink,minimum height=17pt,minimum width=17pt,inner xsep=5pt,inner sep=0]{\popxbold\fontsize{8}{8}\selectfont\color{popcream}\thepage\,/\,\pageref*{LastPage}};}

% ============================================================
% Titres
% ============================================================
\newcommand{\chaptitre}[2]{%
  \begin{tcolorbox}[enhanced,colback=poporange,colframe=popink,boxrule=2.6pt,arc=11pt,
      drop shadow=popink,left=15pt,right=15pt,top=12pt,bottom=11pt,before skip=3pt,after skip=18pt]
    \poppill{#2}{popink}{popcream}\par\vspace{9pt}
    {\raggedright\hyphenpenalty=10000\exhyphenpenalty=10000\popdisplay\fontsize{22}{24}\selectfont\color{popcream}\MakeUppercase{#1}\par}
  \end{tcolorbox}
}
% Section numérotée : pastille orange avec le numéro + titre Archivo
\newcounter{popsec}
\newcommand{\coursec}[1]{%
  \stepcounter{popsec}\setcounter{popsub}{0}%
  \par\vspace{16pt}\noindent
  \begin{minipage}{\linewidth}
  \tikz[baseline=-0.7ex]\node[draw=popink,line width=1.6pt,fill=poporange,rounded corners=4pt,minimum size=22pt,inner sep=0]{\popdisplay\fontsize{12}{12}\selectfont\color{popcream}\thepopsec};%
  \hspace{8pt}{\popdisplay\fontsize{15}{17}\selectfont\color{popink}\MakeUppercase{#1}}\par
  \vspace{4pt}\noindent{\color{poporange}\rule{36pt}{3.6pt}}
  \end{minipage}\par\nopagebreak\vspace{8pt}
}
\newcounter{popsub}
\newcommand{\courssub}[1]{%
  \stepcounter{popsub}%
  \par\vspace{9pt}\noindent{\popxbold\fontsize{11.5}{13}\selectfont\color{popink}{\color{poporange}\thepopsec.\thepopsub}\hspace{6pt}#1}\par\nopagebreak\vspace{4pt}
}

% ============================================================
% Boîtes pédagogiques — même recette : fond blanc, bordure épaisse colorée,
% ombre dure encre, badge plein. Titre optionnel [#1].
% ============================================================
\tcbset{popbase/.style={breakable,enhanced,colback=white,boxrule=2.3pt,arc=9pt,
  shadow={3pt}{-3pt}{0pt}{popink},left=14pt,right=14pt,top=11pt,bottom=11pt,before skip=11pt,after skip=15pt,
  fontupper=\popsemi\fontsize{10.6}{14.5}\selectfont\color{popink},
  fontlower=\popsemi\fontsize{10.6}{14.5}\selectfont\color{popink}}}
% Coche dessinée (le glyphe ✓ n'existe pas dans les polices du document)
\newcommand{\popcheck}[1]{\tikz[baseline=-0.5ex]\draw[#1,line width=1.6pt,line cap=round,line join=round](-0.13,0.0)--(-0.04,-0.09)--(0.13,0.1);}
\providecommand{\coche}{\popcheck{popgreen}}
\providecommand{\resultat}[1]{\par\smallskip\noindent\fcolorbox{popgreen}{popgreenL}{\parbox{\dimexpr\linewidth-2\fboxsep-2\fboxrule\relax}{\textbf{Résultat.} #1}}\par\smallskip}
\newcommand{\popboxhead}[3]{\poppill{#1}{#2}{white}\ifx\relax#3\relax\else\hspace{6pt}{\popxbold\fontsize{10.6}{12}\selectfont\color{popink}#3}\fi\par\vspace{7pt}}

\newtcolorbox{aretenir}[1][]{popbase,colframe=popgreen,before upper={\popboxhead{À retenir}{popgreen}{#1}}}
% Texte en anglais (document de lecture, dialogue, extrait) : cadre
% d'étude. Un titre optionnel (source, auteur) s'affiche dans l'en-tête.
\newtcolorbox{textebox}[1][]{popbase,colframe=popblue,before upper={\popboxhead{Text}{popblue}{#1}\begin{otherlanguage}{english}},after upper={\end{otherlanguage}}}
% Marques ✓ / ✗ (forme correcte / fautive) souvent utilisées par les modèles
\providecommand{\cross}{\textcolor{poppink}{\ensuremath{\times}}}
\providecommand{\xmark}{\cross}\providecommand{\crossmark}{\cross}\providecommand{\cmark}{\ensuremath{\checkmark}}
% Ligne de réponse / justification à compléter (inventée par les modèles)
\providecommand{\justification}[1]{\par\noindent\textit{Justification~:} \dotfill\par}
% \textipa (package tipa non chargé) : la police ne couvre pas l'API -> simple passage
\providecommand{\textipa}[1]{#1}
% Soulignement (ulem non chargé) : \uline, \ul
\providecommand{\uline}[1]{\underline{#1}}\providecommand{\ul}[1]{\underline{#1}}
% \glossaire{\item ...} inventée par les modèles : liste de glossaire
\providecommand{\glossaire}[1]{\begin{itemize}#1\end{itemize}}
% \alert (beamer) inventée par les modèles : texte mis en évidence
\providecommand{\alert}[1]{\textbf{\textcolor{poppink}{#1}}}
% variantes ASCII d'ordinaux écrites en macro
\providecommand{\iere}{\textsuperscript{re}}\providecommand{\ieme}{\textsuperscript{e}}\providecommand{\ere}{\textsuperscript{re}}\providecommand{\eme}{\textsuperscript{e}}\providecommand{\er}{\textsuperscript{er}}\providecommand{\ie}{\textsuperscript{e}}
% Ordinaux français écrits en macro par les modèles : 1\ière, 3\ième
\newcommand{\ière}{\textsuperscript{re}}\newcommand{\ième}{\textsuperscript{e}}
% \exemple (inventée par les modèles) : étiquette « Exemple : »
\providecommand{\exemple}{\textbf{Exemple~:}\ }
% \square utilisable aussi hors mode math (cases à cocher)
\AtBeginDocument{\let\squareMath\square\renewcommand{\square}{\ensuremath{\squareMath}}}
% Mot ou phrase en anglais à l'intérieur d'un paragraphe français
\newcommand{\eng}[1]{\ifmmode\text{\textit{#1}}\else\begingroup\selectlanguage{english}\itshape #1\endgroup\fi}
\let\esp\eng % alias toléré
% flèches d'intonation inventées par les modèles
\providecommand{\textsearrow}{}\renewcommand{\textsearrow}{\ensuremath{\searrow}}\providecommand{\textnearrow}{}\renewcommand{\textnearrow}{\ensuremath{\nearrow}}
% \fr{..} : mot français (inventé par les modèles)
\providecommand{\fr}[1]{\textit{#1}}
\newtcolorbox{exemplebox}[1][]{popbase,colframe=poporange,before upper={\popboxhead{Exemple}{poporange}{#1}}}
\newtcolorbox{definition}[1][]{popbase,colframe=popblue,before upper={\popboxhead{Définition}{popblue}{#1}}}
\newtcolorbox{propriete}[1][]{popbase,colframe=poppurple,before upper={\popboxhead{Propriété}{poppurple}{#1}}}
\newtcolorbox{methode}[1][]{popbase,colframe=popgold,before upper={\popboxhead{Méthode}{popgold}{#1}}}
\newtcolorbox{attention}[1][]{popbase,colframe=poppink,before upper={\popboxhead{Attention}{poppink}{#1}}}
\newtcolorbox{experience}[1][]{popbase,colframe=popblue,colback=popblueL,before upper={\popboxhead{Expérience}{popblue}{#1}}}
% « Le savais-tu ? » : carte violette pleine (contexte, culture, Cameroun)
\newtcolorbox{savaistu}[1][]{popbase,colback=poppurple,colframe=popink,
  fontupper=\popsemi\fontsize{10.4}{14.2}\selectfont\color{white},
  before upper={\poppill{Le savais-tu\,?}{white}{poppurple}\ifx\relax#1\relax\else\hspace{6pt}{\popxbold\color{white}#1}\fi\par\vspace{7pt}}}
% Exercice résolu : énoncé en haut, \tcblower puis la solution sur fond crème
\newcounter{popexo}\newcounter{popexr}
\newtcolorbox{exoresolu}[1][]{popbase,colframe=poporange,colbacklower=popcream,
  segmentation style={popink,line width=1.2pt,dashed},
  before upper={\stepcounter{popexr}\popboxhead{Exercice résolu \thepopexr}{poporange}{#1}},
  before lower={\poppill{Solution}{popink}{popcream}\par\vspace{6pt}}}
% Exercice (non résolu en place) — la correction est dans la partie Corrigés
\newtcolorbox{exercice}[1][]{popbase,colframe=popink,
  before upper={\stepcounter{popexo}\popboxhead{Exercice \thepopexo}{popink}{#1}}}
% Corrigé
\newtcolorbox{corrige}[1][]{popbase,colframe=popgreen,colback=popgreenL,
  before upper={\popboxhead{Corrigé}{popgreen}{#1}}}
% Objectifs / prérequis / bilan
\newtcolorbox{objectifs}{popbase,colframe=popink,colback=poporangeL,
  before upper={\popboxhead{Objectifs de la leçon}{popink}{}}}
\newtcolorbox{prerequis}{popbase,colframe=popink,colback=popblueL,
  before upper={\popboxhead{Prérequis}{popblue}{}}}
\newtcolorbox{bilan}{popbase,colframe=popink,colback=white,boxrule=2.8pt,
  before upper={\popboxhead{Fiche bilan}{poporange}{L'essentiel à connaître pour l'examen}}}
% Figure encadrée avec légende
\newtcolorbox{popfigure}[1][]{enhanced,breakable=false,colback=white,colframe=popline,boxrule=1.4pt,arc=8pt,
  left=10pt,right=10pt,top=10pt,bottom=8pt,before skip=10pt,after skip=12pt,halign=center,
  lower separated=false,halign lower=center,
  fontlower=\popsemi\fontsize{9}{11.5}\selectfont\color{popmuted},
  #1}
\newcommand{\legende}[1]{\par\vspace{6pt}{\popsemi\fontsize{9}{11.5}\selectfont\color{popmuted}#1}}

% ============================================================
% Signature OnBuch+
% ============================================================
\newcommand{\poplogoinline}[1]{{\poplogo\fontsize{#1}{#1}\selectfont\color{popink}OnBuch\color{poporange}+}}
\newcommand{\signatureonbuch}{%
  \begin{tcolorbox}[enhanced,colback=popink,colframe=popink,boxrule=2.2pt,arc=10pt,
      drop shadow=poporange,left=14pt,right=14pt,top=10pt,bottom=10pt,before skip=0pt,after skip=0pt]
    \tikz[baseline=-0.7ex]\node[circle,fill=popgreen,draw=popcream,line width=1.3pt,minimum size=22pt,inner sep=0]{\popcheck{white}};%
    \hspace{9pt}\begin{minipage}[c]{0.62\linewidth}
      {\popxbold\fontsize{12}{13}\selectfont\color{popcream}Signé\ \textbullet\ {\poplogo OnBuch\color{poporange}+}}\par\vspace{2pt}
      {\raggedright\popsemi\fontsize{8}{10.5}\selectfont\color{popline}\MakeUppercase{Équipe pédagogique OnBuch+}\\\MakeUppercase{Conforme au programme officiel MINESEC}\par}
    \end{minipage}\hfill
    {\poplogo\fontsize{22}{22}\selectfont\color{popcream}OnBuch\color{poporange}+}
  \end{tcolorbox}%
}

% ============================================================
% Page de garde
% #1 titre · #2 matière · #3 niveau · #4 module · #5 n° leçon · #6 durée ·
% #7 description / objectifs (texte ou itemize)
% ============================================================
\newcommand{\pagegarde}[7]{%
  \thispagestyle{empty}
  \newgeometry{a4paper,margin=1.7cm,top=1.6cm,bottom=1.4cm}%
  \begin{tikzpicture}[remember picture,overlay]
    \fill[poporange!18] ([xshift=-3.2cm,yshift=-3.6cm]current page.north east) circle (4.6cm);
    \fill[poppurple!14] ([xshift=2.4cm,yshift=7.6cm]current page.south west) circle (3.8cm);
  \end{tikzpicture}%
  {\poplogo\fontsize{30}{30}\selectfont\color{popink}OnBuch\color{poporange}+}\hfill
  \raisebox{0.6ex}{\poppill{Cours signé \textbullet\ Premium}{popink}{popcream}}\par
  \vspace{1pt}\popsquiggle{poppurple}\par
  \vspace{8pt}
  \begin{tcolorbox}[enhanced,colback=poporange,colframe=popink,boxrule=2.8pt,arc=13pt,
      drop shadow=popink,left=18pt,right=18pt,top=12pt,bottom=13pt,after skip=10pt]
    \poppill{#2 \textbullet\ #3}{popink}{popcream}\hspace{5pt}\poppill{#4}{white}{popink}\par\vspace{8pt}
    {\popdisplay\fontsize{14}{14}\selectfont\color{popink}LEÇON #5}\par\vspace{4pt}
    {\raggedright\hyphenpenalty=10000\exhyphenpenalty=10000\popdisplay\fontsize{25}{27}\selectfont\color{popcream}\MakeUppercase{#1}\par}
    \vspace{8pt}
    {\popxbold\fontsize{9.5}{9.5}\selectfont\color{popink}\MakeUppercase{Durée conseillée : #6}}
  \end{tcolorbox}
  \begin{tcolorbox}[enhanced,colback=white,colframe=popink,boxrule=2.2pt,arc=10pt,
      drop shadow=popink,left=16pt,right=16pt,top=10pt,bottom=10pt,after skip=8pt]
    \poppill{Dans cette leçon}{poppurple}{white}\par\vspace{5pt}
    \popsemi\fontsize{9.3}{12.6}\selectfont\color{popink}#7
  \end{tcolorbox}
  \noindent\begin{minipage}[t]{0.487\linewidth}
    \begin{tcolorbox}[enhanced,colback=popgreen,colframe=popink,boxrule=2.2pt,arc=10pt,
        drop shadow=popink,left=11pt,right=11pt,top=10pt,bottom=10pt,height=3.1cm,valign=center]
      \tcbox[colback=white,colframe=popink,boxrule=1.3pt,arc=5pt,boxsep=0pt,nobeforeafter,
        left=3pt,right=3pt,top=3pt,bottom=3pt,valign=center]{\qrcode[height=1.9cm]{https://play.google.com/store/apps/details?id=com.luvvix.onbuch&hl=fr}}%
      \hspace{8pt}\begin{minipage}[c]{0.5\linewidth}
        {\popdisplay\fontsize{11}{12.5}\selectfont\color{white}\MakeUppercase{Télécharge OnBuch}}\par\vspace{4pt}
        {\raggedright\popsemi\fontsize{8.4}{11}\selectfont\color{white}Gratuit \textbullet\ Google Play\\Tuteur IA, annales, concours, résultats\par}
      \end{minipage}
    \end{tcolorbox}
  \end{minipage}\hfill
  \begin{minipage}[t]{0.487\linewidth}
    \begin{tcolorbox}[enhanced,colback=poppurple,colframe=popink,boxrule=2.2pt,arc=10pt,
        drop shadow=popink,left=11pt,right=11pt,top=10pt,bottom=10pt,height=3.1cm,valign=center]
      \tikz[baseline=-0.7ex]\node[circle,fill=white,draw=popink,line width=1.3pt,minimum size=1.6cm,inner sep=0]{\popdisplay\fontsize{24}{24}\selectfont\color{poppurple}+};%
      \hspace{8pt}\begin{minipage}[c]{0.6\linewidth}
        {\popdisplay\fontsize{11}{12.5}\selectfont\color{white}\MakeUppercase{Passe à OnBuch+}}\par\vspace{4pt}
        {\raggedright\popsemi\fontsize{8.4}{11}\selectfont\color{white}Tous les cours signés, Tuteur IA illimité, fiches PDF — onglet OnBuch+ de l'app\par}
      \end{minipage}
    \end{tcolorbox}
  \end{minipage}
  \vspace{8pt}
  \popcontenu
  \vfill
  \signatureonbuch
  \restoregeometry
  \newpage
}

% Rangée « Ce cours contient » (page de garde)
\newcommand{\poptile}[3]{% #1 couleur #2 gros texte #3 légende
  \begin{tcolorbox}[enhanced,colback=#1,colframe=popink,boxrule=2pt,arc=9pt,shadow={2.5pt}{-2.5pt}{0pt}{popink},
      left=6pt,right=6pt,top=9pt,bottom=9pt,halign=center,valign=center,height=1.75cm,width=\linewidth,nobeforeafter]
    {\popdisplay\fontsize{12.5}{13}\selectfont\color{white}\MakeUppercase{#2}}\par\vspace{3pt}
    {\centering\popsemi\fontsize{7.8}{9.5}\selectfont\color{white}#3\par}
  \end{tcolorbox}}
\newcommand{\popcontenu}{%
  \par\noindent{\popxboldls\fontsize{7.5}{7.5}\selectfont\color{popmuted}CE COURS SIGNÉ CONTIENT}\par\vspace{7pt}
  \noindent\begin{minipage}[t]{0.235\linewidth}\poptile{popblue}{Cours}{complet, illustré, pas à pas}\end{minipage}\hfill
  \begin{minipage}[t]{0.235\linewidth}\poptile{poporange}{Méthodes}{et exercices résolus}\end{minipage}\hfill
  \begin{minipage}[t]{0.235\linewidth}\poptile{poppink}{Exercices}{type examen + corrigés}\end{minipage}\hfill
  \begin{minipage}[t]{0.235\linewidth}\poptile{popgreen}{Fiche bilan}{l'essentiel à retenir}\end{minipage}\par}

% Sommaire
\newtcolorbox{sommaire}{enhanced,colback=white,colframe=popink,boxrule=2.2pt,arc=10pt,
  drop shadow=popink,left=18pt,right=18pt,top=13pt,bottom=13pt,before skip=6pt,after skip=20pt,
  before upper={\poppill{Sommaire}{poppurple}{white}\par\vspace{9pt}\popsemi\fontsize{10.6}{14.5}\selectfont\color{popink}}}

% Signature de fin de cours — poussée tout en bas de la dernière page
\newcommand{\findecours}{%
  \par\vfill\nopagebreak
  \begin{center}\popsquiggle{poporange}\end{center}\vspace{4pt}
  \signatureonbuch
}
\AtBeginDocument{\def\llbracket{\mathopen{[\mkern-3.5mu[}}\def\rrbracket{\mathclose{]\mkern-3.5mu]}}} % intervalles d'entiers (glyphe ⟦ absent de la police)
% Glyphes absents de Fira Math : redéfinis à partir de glyphes disponibles
\AtBeginDocument{%
  \def\setminus{\mathbin{\backslash}}\let\smallsetminus\setminus
  \def\vdots{\mathord{\vbox{\baselineskip3.2pt\lineskiplimit0pt\kern4pt\hbox{.}\hbox{.}\hbox{.}}}}%
  \def\ddots{\mathinner{\mkern1mu\raise6.5pt\hbox{.}\mkern2mu\raise3.6pt\hbox{.}\mkern2mu\raise0.7pt\hbox{.}\mkern1mu}}%
  \def\triangle{\mathord{\scalebox{0.9}{$\symup{\Delta}$}}}%
  \def\nabla{\mathord{\raisebox{\depth}{\scalebox{1}[-1]{$\symup{\Delta}$}}}}%
  \def\checkmark{\popcheck{popdark}}%
  \def\star{\mathbin{\ast}}\def\bigstar{\mathord{\ast}}%
  \def\diamond{\mathbin{\scalebox{0.6}{$\Diamond$}}}%
  \def\bigcup{\mathop{\mathchoice{\raisebox{-0.3em}{\scalebox{1.7}{$\cup$}}}{\scalebox{1.25}{$\cup$}}{\cup}{\cup}}\displaylimits}%
  \def\bigcap{\mathop{\mathchoice{\raisebox{-0.3em}{\scalebox{1.7}{$\cap$}}}{\scalebox{1.25}{$\cap$}}{\cap}{\cap}}\displaylimits}%
  \def\lesseqqgtr{\mathrel{\lessgtr}}\def\bigoplus{\mathop{\oplus}\displaylimits}%
}
\providecommand{\ssi}{si et seulement si} % abréviation fréquente
\AtBeginDocument{\providecommand{\wideparen}{\overparen}} % arc de cercle

\begin{document}

\pagegarde{Lesson 13 — Listening: Texts on comparing prices or quality to determine the best buys for consumer economy and community resources/services}{Anglais}{Tle A}{Chapter 2 : Economic Life and Occupations}{EN13}{2 h}{Cette leçon de compréhension orale entraîne les élèves à écouter des textes comparant prix et qualité afin d'identifier les meilleurs achats (best buys). Elle développe les stratégies d'écoute (repérage du thème, des mots-clés, des chiffres et des noms propres) et enrichit le lexique de la vie économique et des services communautaires.
\vspace{4pt}
\begin{multicols}{2}\small
\begin{itemize}
\item À la fin de cette leçon, je sais identifier le thème général d'un enregistrement sur la comparaison des prix et de la qualité.
\item Je sais repérer les mots-clés, les chiffres (prix, pourcentages, quantités) et les noms propres dans un texte écouté.
\item Je sais relever les informations essentielles et les détails pour déterminer le meilleur achat (best buy).
\item Je sais répondre à des questions de compréhension de types variés : vrai/faux, QCM et réponses courtes.
\item Je sais utiliser le lexique de la consommation (price, quality, bargain, value for money...) dans de courtes phrases.
\item Je sais comparer des produits ou des services en anglais à l'aide des comparatifs et superlatifs.
\end{itemize}\end{multicols}}

\begin{sommaire}
\begin{multicols}{2}
\begin{enumerate}
\item Warm-up et découverte du thème
\item Lexique thématique : consommation et comparaison
\item Stratégies et méthode d'écoute
\item Texte support : script d'écoute 'The Best Buy'
\item Compréhension : questions de types variés
\item Production guidée : conseiller le meilleur achat
\item Activité d'intégration
\item Méthodes et exercices résolus
\item Exercices
\item Corrigés des exercices
\item Fiche bilan
\end{enumerate}
\end{multicols}
\end{sommaire}

\begin{objectifs}
\begin{itemize}
\item À la fin de cette leçon, je sais identifier le thème général d'un enregistrement sur la comparaison des prix et de la qualité.
\item Je sais repérer les mots-clés, les chiffres (prix, pourcentages, quantités) et les noms propres dans un texte écouté.
\item Je sais relever les informations essentielles et les détails pour déterminer le meilleur achat (best buy).
\item Je sais répondre à des questions de compréhension de types variés : vrai/faux, QCM et réponses courtes.
\item Je sais utiliser le lexique de la consommation (price, quality, bargain, value for money...) dans de courtes phrases.
\item Je sais comparer des produits ou des services en anglais à l'aide des comparatifs et superlatifs.
\item Je sais prendre des notes efficaces pendant l'écoute d'un texte lu ou enregistré.
\item Je sais produire un court paragraphe ou dialogue conseillant le meilleur achat à un consommateur.
\end{itemize}
\end{objectifs}

\begin{prerequis}
\begin{itemize}
\item Les comparatifs et superlatifs (cheaper than, the best value) vus en classe de Première
\item Les nombres, prix et monnaies en anglais (pounds, dollars, francs CFA)
\item Le vocabulaire de base des achats et du marché (to buy, to sell, shop, market)
\item Les question words (what, where, how much, which) pour la compréhension
\item La notion de consumer goods et services vue au chapitre 2 (Economic Life and Occupations)
\end{itemize}
\end{prerequis}

\newpage

\coursec{Warm-up et découverte du thème}

\courssub{Situation d'accroche : le dilemme de Mama Béa}

\begin{textebox}[At the market — a consumer's dilemma]
At Mokolo Market in Yaoundé, Mama Béa hesitates between two bags of rice: one costs 12,000 FCFA, the other 15,000 FCFA but is of better quality. Which one is the \cle{best buy}? Every day, consumers in Cameroon, Nigeria, the UK or the USA compare prices and quality before spending their money. Radio programmes and consumer reports help them make wise choices. In this lesson, you will listen to such a programme and learn how to catch the essential information to determine the best buys.
\end{textebox}

\textbf{Problématique de la leçon.} Comment un consommateur peut-il comparer les \emph{prix} et la \emph{qualité} pour déterminer le \cle{meilleur achat} (\eng{best buy}) ? Et comment, à l'écoute d'un texte en anglais, repérer rapidement les chiffres, les noms de produits et les arguments qui permettent de trancher ? C'est toute la compétence que cette leçon va développer : \emph{écouter pour comparer, comparer pour choisir}.

\begin{definition}[Best buy]
A \cle{best buy} is the product or service that offers the best \cle{value for money} — in other words, the best balance between price and quality. En français : le \emph{meilleur rapport qualité-prix}. Attention : le meilleur achat n'est pas forcément le produit le moins cher, mais celui qui donne le plus de satisfaction pour l'argent dépensé.
\par\eng{The 15,000 FCFA bag is more expensive, but it is the best buy because the rice lasts longer and tastes better.} « Le sac à 15 000 FCFA est plus cher, mais c'est le meilleur achat car le riz dure plus longtemps et a meilleur goût. »
\end{definition}

\courssub{Brainstorming : où achète-t-on moins cher ?}

Avant toute écoute, un bon élève \emph{active ses connaissances} : il fait remonter à la surface tout ce qu'il sait déjà sur le sujet. Répondez mentalement à ces questions, en anglais :

\begin{itemize}
\item \eng{Where do people in Yaoundé, Douala or Bamenda buy food more cheaply: at the market, at the supermarket or at roadside stalls?} « Où les habitants de Yaoundé, Douala ou Bamenda achètent-ils la nourriture moins cher : au marché, au supermarché ou aux étalages au bord de la route ? »
\item \eng{Where can you bargain, and where are prices fixed?} « Où peut-on marchander, et où les prix sont-ils fixes ? »
\item \eng{Is the cheapest product always the best buy? Why or why not?} « Le produit le moins cher est-il toujours le meilleur achat ? Pourquoi ? »
\end{itemize}

Voici le vocabulaire des \emph{lieux d'achat} qui reviendra dans l'enregistrement :

\begin{center}
\begin{tabularx}{\linewidth}{|X|X|X|}\hline
\rowcolor{popblueL} English & Français & Remarque \\ \hline
a market (an open-air market) & un marché (en plein air) & \eng{Mokolo Market} \\ \hline
a supermarket & un supermarché & prix fixes, \eng{fixed prices} \\ \hline
a roadside stall & un étalage au bord de la route & souvent le moins cher \\ \hline
a shop / a store (US) & un magasin & \eng{store} est américain \\ \hline
a hawker / a street vendor & un vendeur ambulant & prononcer « ho-keur » \\ \hline
a wholesaler & un grossiste & vend en gros \\ \hline
a retailer & un détaillant & vend au détail \\ \hline
\end{tabularx}
\end{center}

\begin{exemplebox}[Parler des lieux d'achat]
\eng{Tomatoes are cheaper at the market than at the supermarket.} « Les tomates sont moins chères au marché qu'au supermarché. »
\par\eng{At roadside stalls, you can often bargain over the price.} « Aux étalages au bord de la route, on peut souvent marchander. »
\par\eng{Supermarkets are more expensive, but the products have a guarantee.} « Les supermarchés sont plus chers, mais les produits ont une garantie. »
\end{exemplebox}

\courssub{Les critères du bon choix : price, quality, durability, brand, guarantee}

Quand Mama Béa hésite entre les deux sacs de riz, elle compare en réalité plusieurs \cle{critères}. Les voici, avec leur vocabulaire anglais :

\begin{center}
\begin{tabularx}{\linewidth}{|X|X|X|}\hline
\rowcolor{popblueL} Critère (English) & Français & Question du consommateur \\ \hline
price & le prix & \eng{How much does it cost?} \\ \hline
quality & la qualité & \eng{Is it well made?} \\ \hline
durability & la durabilité & \eng{How long will it last?} \\ \hline
brand & la marque & \eng{Is it a reliable brand?} \\ \hline
guarantee / warranty & la garantie & \eng{Is there a guarantee?} \\ \hline
after-sales service & le service après-vente & \eng{Can they repair it if it breaks?} \\ \hline
\end{tabularx}
\end{center}

\begin{attention}[Cheap : un mot à manier avec prudence]
\eng{Cheap} signifie « bon marché », mais il peut aussi suggérer « de mauvaise qualité », « qui ne vaut rien ». \eng{This radio is cheap} peut vouloir dire « cette radio est bas de gamme ». Pour rester neutre et poli, utilisez \eng{inexpensive} (« peu coûteux ») ou \eng{affordable} (« abordable »).
\par\eng{This phone is inexpensive and very reliable.} « Ce téléphone est peu coûteux et très fiable. » (neutre, positif)
\par\eng{I don't want a cheap phone that will break in two weeks.} « Je ne veux pas d'un téléphone bas de gamme qui cassera en deux semaines. » (péjoratif)
\par De même, ne confondez pas \eng{economic} (« économique », relatif à l'économie) et \eng{economical} (« économe », qui fait faire des économies) : \eng{an economical car} = « une voiture économe ».
\end{attention}

La carte mentale suivante rassemble les cinq critères du \eng{best buy}. Apprenez-la : ce sont exactement les mots-clés à guetter pendant l'écoute.

\begin{popfigure}
\begin{center}
\begin{tikzpicture}[mindmap, grow cyclic, every node/.style={concept}, concept color=popblue!60, text=white,
level 1/.append style={level distance=3.2cm, sibling angle=72},
level 1 concept/.append style={font=\small}]
\node[concept, font=\large\bfseries] {BEST BUY}
child[concept color=poporange!80] { node[align=center]{price\\cost} }
child[concept color=popgreen!70!black] { node {quality} }
child[concept color=poppurple!80] { node {durability} }
child[concept color=poppink!80] { node {guarantee} }
child[concept color=popgold!90!black] { node[align=center] {after-sales\\service} };
\end{tikzpicture}
\end{center}
\legende{Carte mentale : les cinq critères du \eng{best buy} (le meilleur rapport qualité-prix).}
\end{popfigure}

\courssub{Prédiction : de quoi va parler l'enregistrement ?}

Une stratégie d'écoute essentielle consiste à \emph{prédire} le contenu d'un enregistrement avant de l'entendre, à partir de son titre, de son contexte et de ce que l'on sait du thème. L'enregistrement de cette leçon s'intitule \eng{The Best Buy} et sera diffusé dans le cadre d'une émission de radio de conseil aux consommateurs.

\begin{methode}[Prédire avant d'écouter]
\begin{enumerate}
\item Lis le titre (\eng{The Best Buy}) et demande-toi : de quoi va-t-on parler ? (des achats, de la comparaison de produits).
\item Note les mots que tu t'attends à entendre : \eng{price, cheaper, quality, guarantee, FCFA, market, supermarket}...
\item Anticipe les chiffres : des prix, des pourcentages, des quantités. Prépare-toi à les noter vite.
\item Formule deux ou trois questions auxquelles l'enregistrement devrait répondre : \eng{Which product is recommended? How much does it cost? Why is it the best buy?}
\end{enumerate}
\end{methode}

\begin{experience}[Classroom warm-up — Predicting]
En binômes, regardez le titre \eng{The Best Buy} et la carte mentale des critères. Chaque binôme écrit trois prédictions en anglais, par exemple : \eng{The programme will compare two products.} / \eng{The speaker will give prices in FCFA.} / \eng{The cheaper product will not be the best buy.} Après l'écoute (section 4), vous vérifierez quelles prédictions étaient correctes.
\end{experience}

\courssub{Texte de mise en train : Consumers' Corner}

Pour vous plonger dans le thème et préparer votre oreille au vocabulaire de l'écoute, lisez ce court article de presse original. Il utilise exactement le type de langue que vous entendrez dans l'enregistrement.

\begin{textebox}[Consumers' Corner — a newspaper article]
\textbf{Smart shopping: how Douala families find the best buys}
\par Every Saturday morning, thousands of shoppers walk through the markets of Douala with the same question in mind: where can I find the best buy? For Mrs Etoundi, a mother of four, the answer is simple. “I compare prices in at least three shops before I buy anything,” she explains. “Last week, a 25-kilo bag of rice cost 14,500 FCFA at one stall and 13,000 FCFA at another. Same brand, same quality!”
\par But price is not everything. Consumer experts warn that the cheapest product is not always the most economical one. A cheap pair of shoes that lasts two months is more expensive, in the end, than a good pair that lasts two years. Durability matters. So does the guarantee: a radio with a one-year guarantee is often a better buy than a cheaper one with no guarantee at all.
\par In Western countries, consumer magazines and radio programmes test products and publish comparisons. In Cameroon, shoppers rely mostly on word of mouth: they ask friends, neighbours and market women for advice. “My sister told me this brand of oil is the best,” says Mrs Etoundi. “She has used it for years.”
\par Whatever the method, the golden rule remains the same everywhere: never buy the first product you see. Compare the price, check the quality, ask about the guarantee — and then decide. That is how smart consumers get real value for money.
\end{textebox}

\textbf{Glossaire.} \emph{a shopper} (n.) : un acheteur ; \emph{a stall} (n.) : un étalage ; \emph{to warn} (v.) : avertir ; \emph{to last} (v.) : durer ; \emph{word of mouth} (expr.) : le bouche-à-oreille ; \emph{to rely on} (v.) : compter sur, se fier à ; \emph{the golden rule} (expr.) : la règle d'or ; \emph{value for money} (expr.) : le rapport qualité-prix.

\textbf{Questions de compréhension rapide.}
\begin{enumerate}
\item \eng{How many shops does Mrs Etoundi visit before buying?} (three)
\item \eng{What two prices does she mention for the bag of rice?} (14,500 FCFA and 13,000 FCFA)
\item \eng{According to experts, why is the cheapest product not always the best buy?} (because a cheap product may not last long)
\item \eng{How do Cameroonian shoppers usually get information about products?} (by word of mouth)
\end{enumerate}

\begin{savaistu}[Le magazine Which? et le bouche-à-oreille camerounais]
Au Royaume-Uni, le magazine des consommateurs \eng{Which?}, fondé en 1957, teste des centaines de produits (téléphones, voitures, machines à laver...) et publie des comparatifs pour aider les Britanniques à trouver le \eng{best buy}. Aux États-Unis, \eng{Consumer Reports} joue le même rôle. Au Cameroun, il n'existe pas d'équivalent aussi connu : les consommateurs se fient surtout au \eng{word of mouth} (bouche-à-oreille), à l'expérience des vendeuses du marché et à la comparaison directe des prix. C'est pourquoi les émissions de radio de conseil aux consommateurs, comme celle que vous allez écouter, sont si utiles.
\end{savaistu}

\courssub{Les objectifs d'écoute de la leçon}

Écouter en anglais n'est pas « tout comprendre mot à mot » : c'est \emph{écouter avec un objectif précis}. Pour cette leçon, vous vous entraînerez à trois niveaux d'écoute, du plus global au plus fin :

\begin{center}
\begin{tabularx}{\linewidth}{|X|X|X|}\hline
\rowcolor{popblueL} Niveau d'écoute & Objectif & Ce qu'il faut noter \\ \hline
\eng{Listening for the gist} & saisir le thème général & le sujet, le type d'émission \\ \hline
\eng{Listening for specific information} & repérer les chiffres et les noms propres & prix, marques, lieux, durées \\ \hline
\eng{Listening for detail} & comprendre les arguments & raisons du choix, avantages, défauts \\ \hline
\end{tabularx}
\end{center}

\begin{aretenir}[Les mots-clés à guetter dans un texte de comparaison]
\begin{itemize}
\item Les \textbf{chiffres} : \eng{12,000 FCFA, 15 per cent off, a two-year guarantee}.
\item Les \textbf{comparatifs} : \eng{cheaper than, more expensive than, better than, the best, the most economical}.
\item Les \textbf{marques et noms propres} : noms de produits, de magasins, de quartiers.
\item Les \textbf{expressions de conseil} : \eng{we recommend, you should, the best buy is, value for money}.
\end{itemize}
\end{aretenir}

\begin{exoresolu}[Warm-up : classer les critères]
Énoncé — Voici cinq phrases entendues à la radio. Pour chacune, indiquez le critère du \eng{best buy} qu'elle illustre : \eng{price, quality, durability, guarantee} ou \eng{after-sales service}.
\begin{enumerate}
\item \eng{This blender costs only 9,500 FCFA, the lowest price in town.}
\item \eng{The shop will repair your television free of charge if it breaks down.}
\item \eng{These batteries come with a six-month guarantee.}
\item \eng{This fabric is made of pure cotton; it feels soft and strong.}
\item \eng{A good pair of leather shoes can last five years or more.}
\end{enumerate}
\tcblower
Solution détaillée —
\begin{enumerate}
\item \eng{price} : la phrase insiste sur le montant (9,500 FCFA) et le superlatif \eng{the lowest price}.
\item \eng{after-sales service} : \eng{will repair... free of charge} décrit un service après l'achat.
\item \eng{guarantee} : le mot \eng{guarantee} est cité explicitement (\eng{a six-month guarantee}).
\item \eng{quality} : \eng{pure cotton, soft and strong} décrit la qualité du produit.
\item \eng{durability} : \eng{can last five years} parle de la durée de vie du produit.
\end{enumerate}
\end{exoresolu}

\begin{exercice}[À vous de jouer — préparer l'écoute]
\begin{enumerate}
\item Traduisez en anglais : « Le produit le moins cher n'est pas toujours le meilleur achat. »
\item Complétez avec le bon mot (\eng{cheap, inexpensive, guarantee, durability, brand}) : \eng{This radio is not \dots\ but it is very \dots\ — only 8,000 FCFA — and it has a one-year \dots .}
\item Écrivez deux phrases pour comparer le marché et le supermarché de votre ville, en utilisant \eng{cheaper than} et \eng{more expensive than}.
\item Formulez, en anglais, trois questions que l'enregistrement \eng{The Best Buy} devrait selon vous permettre de résoudre.
\end{enumerate}
\end{exercice}

\begin{corrige}[Exercice « À vous de jouer »]
\begin{enumerate}
\item \eng{The cheapest product is not always the best buy.} (Attention à l'ordre : superlatif + nom ; ne dites pas « the most cheap ».)
\item \eng{This radio is not \textbf{cheap} but it is very \textbf{inexpensive} — only 8,000 FCFA — and it has a one-year \textbf{guarantee}.} (On accepte aussi \eng{expensive} dans le premier blanc : \eng{not expensive but very inexpensive} est redondant mais correct ; la meilleure réponse reste \eng{cheap}.)
\item Réponses possibles : \eng{Vegetables are cheaper at the market than at the supermarket.} / \eng{Rice is more expensive at the supermarket, but the quality is better.}
\item Exemples : \eng{Which products will be compared?} / \eng{How much does each product cost?} / \eng{Which one is the best buy, and why?}
\end{enumerate}
\end{corrige}

\coursec{Lexique thématique : consommation et comparaison}

\courssub{Verbes d'action du consommateur}

Voici les verbes essentiels pour décrire l'acte d'acheter, de comparer et de décider. Maîtrisez leur sens et leur construction (transitivité, prépositions).

\begin{center}
\begin{tabularx}{\linewidth}{|X|X|X|}\hline
\rowcolor{popblueL} Verbe (infinitif) & Sens & Construction / Exemple \\ \hline
to compare & comparer & \eng{compare A with/to B} — \eng{Compare the prices with those of last year.} \\ \hline
to cost & coûter & \eng{cost + montant} (pas de to) — \eng{It costs 5,000 FCFA.} \\ \hline
to last & durer & \eng{last + durée} — \eng{This battery lasts two years.} \\ \hline
to guarantee & garantir & \eng{guarantee sth for + durée} — \eng{They guarantee the engine for five years.} \\ \hline
to recommend & recommander & \eng{recommend sth to sb} / \eng{recommend doing sth} — \eng{I recommend this brand.} \\ \hline
to afford & pouvoir s'offrir (financièrement) & \eng{can/could afford + nom / to + V} — \eng{We can't afford a new car.} \\ \hline
to bargain (UK) / to haggle & marchander & \eng{bargain over/for sth} — \eng{She bargained over the price.} \\ \hline
to save & économiser / sauver & \eng{save money / save up for sth} — \eng{He saves 10\% of his salary.} \\ \hline
to spend & dépenser & \eng{spend money on sth} — \eng{Don't spend all your money on clothes.} \\ \hline
to waste & gaspiller & \eng{waste money on sth} — \eng{Buying cheap tools is a waste of money.} \\ \hline
\end{tabularx}
\end{center}

\begin{attention}[Cost : pas de "to" après le verbe]
Erreur fréquente : \eng{*It costs to 5,000 FCFA.}
Correct : \eng{It costs 5,000 FCFA.} (verbe transitif direct).
De même : \eng{How much does it cost?} (pas \eng{cost to}).
\end{attention}

\courssub{Noms et expressions utiles pour l'écoute}

\begin{center}
\begin{tabularx}{\linewidth}{|X|X|X|}\hline
\rowcolor{popblueL} Terme anglais & Français & Contexte / Collocation \\ \hline
a bargain & une bonne affaire & \eng{to pick up a bargain, a real bargain} \\ \hline
a receipt & un reçu / un ticket de caisse & \eng{Keep your receipt for the guarantee.} \\ \hline
a refund & un remboursement & \eng{ask for a refund, get a full refund} \\ \hline
an exchange & un échange & \eng{exchange policy, exchange an item} \\ \hline
a price tag & une étiquette de prix & \eng{The price tag says 10,000 FCFA.} \\ \hline
a label & une étiquette (composition) & \eng{Read the label carefully.} \\ \hline
expiry date / best before date & date de péremption / DLUO & \eng{Check the expiry date.} \\ \hline
value for money & rapport qualité-prix & \eng{good/excellent/poor value for money} \\ \hline
a discount / a reduction & une remise / une réduction & \eng{10\% discount, a price reduction} \\ \hline
a special offer & une offre spéciale & \eng{special offer: buy one get one free} \\ \hline
\end{tabularx}
\end{center}

\begin{exemplebox}[Lexique en contexte]
\eng{Always keep the \textbf{receipt} in case you need a \textbf{refund} or an \textbf{exchange}.} « Gardez toujours le reçu au cas où vous auriez besoin d'un remboursement ou d'un échange. »
\par\eng{This laptop is excellent \textbf{value for money}; it has a two-year \textbf{guarantee}.} « Cet ordinateur a un excellent rapport qualité-prix ; il a une garantie de deux ans. »
\par\eng{There is a 20\% \textbf{discount} on all shoes this week.} « Il y a 20\% de remise sur toutes les chaussures cette semaine. »
\end{exemplebox}

\courssub{Comparatif et superlatif : rappels indispensables}

L'écoute de comparaison repose massivement sur ces formes. Révisez-les avant l'exercice d'écoute.

\begin{propriete}[Comparatif de supériorité / infériorité / égalité]
\begin{itemize}
\item \textbf{Court (1 syllabe)} : \eng{cheap $\rightarrow$ cheaper than} / \eng{less cheap than} / \eng{as cheap as}.
\item \textbf{Long ($\ge$ 2 syllabes)} : \eng{expensive $\rightarrow$ more expensive than} / \eng{less expensive than} / \eng{as expensive as}.
\item \textbf{Irréguliers} : \eng{good $\rightarrow$ better than} ; \eng{bad $\rightarrow$ worse than} ; \eng{far $\rightarrow$ further/farther than} ; \eng{much/many $\rightarrow$ more than} ; \eng{little $\rightarrow$ less than}.
\end{itemize}
\end{propriete}

\begin{propriete}[Superlatif]
\begin{itemize}
\item \textbf{Court} : \eng{the cheapest} / \eng{the least cheap}.
\item \textbf{Long} : \eng{the most expensive} / \eng{the least expensive}.
\item \textbf{Irréguliers} : \eng{the best} ; \eng{the worst} ; \eng{the most} ; \eng{the least}.
\end{itemize}
\par Attention : \eng{the best buy} (nom composé), mais \eng{the best product}.
\end{propriete}

\begin{exemplebox}[Comparer pour l'écoute]
\eng{Brand A is \textbf{cheaper than} Brand B, but Brand B is \textbf{better quality}.} « La marque A est moins chère que la marque B, mais la marque B est de meilleure qualité. »
\par\eng{This is \textbf{the most economical} choice in the long run.} « C'est le choix le plus économique sur le long terme. »
\par\eng{Product X is \textbf{as durable as} Product Y, but \textbf{less expensive}.} « Le produit X est aussi durable que le produit Y, mais moins cher. »
\end{exemplebox}

\begin{exercice}[Entraînement : formes comparatives]
Mettez l'adjectif entre parenthèses à la forme correcte (comparatif ou superlatif).
\begin{enumerate}
\item This market is \dots\ (cheap) than the supermarket.
\item That was \dots\ (bad) purchase I ever made.
\item Brand X is \dots\ (reliable) than Brand Y.
\item This fabric is not \dots\ (durable) as the other one.
\item Which one is \dots\ (good) value for money?
\end{enumerate}
\end{exercice}

\begin{corrige}[Exercice « Formes comparatives »]
\begin{enumerate}
\item \textbf{cheaper} (court : -er than)
\item \textbf{the worst} (superlatif irrégulier de bad)
\item \textbf{more reliable} (long : more + adj + than)
\item \textbf{as durable} (égalité : as + adj + as)
\item \textbf{the best} (superlatif irrégulier de good)
\end{enumerate}
\end{corrige}

\coursec{Stratégies et méthode d'écoute}

\courssub{La méthode en 4 étapes (avant / pendant / après)}

\begin{methode}[Réussir une compréhension orale de type "comparatif"]
\begin{enumerate}
\item \textbf{Avant l'écoute (Préparation — 2 min)} : Lisez le titre, les questions, les items de QCM/Vrai-Faux. Soulignez les \emph{mots-clés} (noms produits, chiffres, critères). Prédisez le vocabulaire (voir § Prédiction).
\item \textbf{Première écoute (Globalité — Gist)} : Ne prenez pas de notes détaillées. Répondez : \eng{What is the topic? How many products? What is the final recommendation?} Cochez l'idée générale.
\item \textbf{Deuxième écoute (Détails — Specific info)} : Notez les \textbf{chiffres} (prix, quantités, durées), les \textbf{noms propres} (marques, magasins), les \textbf{arguments} (avantages/inconvénients). Utilisez un tableau de prise de notes (voir ci-dessous).
\item \textbf{Troisième écoute (Vérification — Detail)} : Vérifiez vos réponses, complétez les trous, confirmez les Vrai/Faux. Attention aux \emph{distracteurs} (infos données puis corrigées : \eng{At first I thought... but actually...}).
\end{enumerate}
\end{methode}

\courssub{Grille de prise de notes type "Best Buy"}

Préparez ce tableau \emph{avant} la deuxième écoute. Remplissez-le au fur et à mesure.

\begin{center}
\begin{tabularx}{\linewidth}{|X|X|X|X|X|}\hline
\rowcolor{popblueL} Critère & \textbf{Produit A : \dots} & \textbf{Produit B : \dots} & \textbf{Produit C (si applicable)} & \textbf{Verdict / Notes} \\ \hline
Price (Prix) & & & & \\ \hline
Quality (Qualité) & & & & \\ \hline
Durability (Durabilité) & & & & \\ \hline
Guarantee (Garantie) & & & & \\ \hline
After-sales service (SAV) & & & & \\ \hline
\cle{Best Buy?} & & & & \\ \hline
\end{tabularx}
\end{center}

\begin{attention}[Pièges sonores fréquents]
\begin{itemize}
\item \textbf{Chiffres} : \eng{15,000} (quinze mille) vs \eng{50,000} (cinquante mille) — écoutez le \eng{teen} / \eng{ty}. \eng{13,000} vs \eng{30,000}.
\item \textbf{Négations} : \eng{not cheaper}, \eng{no guarantee}, \eng{doesn't last long}. Le \eng{not} peut être "avalé" à l'oral (\eng{isn't} $\rightarrow$ [\eng{iznt}]).
\item \textbf{Connecteurs de concession} : \eng{however, but, although, even though, despite} signalent souvent l'argument décisif (le \eng{best buy} n'est pas le premier présenté).
\item \textbf{Hésitations / Corrections} : \eng{It costs 12,000... sorry, 15,000 FCFA.} Notez la correction finale.
\end{itemize}
\end{attention}

\coursec{Texte support : script d'écoute « The Best Buy »}

\courssub{Consigne}
Vous allez écouter (ou lire) l'extrait suivant d'une émission de radio intitulée \eng{Consumers' Corner}. L'animateur compare deux marques d'huile végétale très courantes au Cameroun. Utilisez la grille de prise de notes ci-dessus.

\begin{textebox}[Listening Script — Consumers' Corner: "The Best Buy Edition"]
\textbf{Host:} Good morning, listeners, and welcome to \eng{Consumers' Corner}, your weekly guide to smart shopping. I'm your host, Peter Ngwa. Today, we tackle a question many housewives in Yaoundé, Douala and Buea ask themselves every market day: which cooking oil gives the best value for money? We have tested two leading brands on sale in our major supermarkets and open-air markets: \eng{"Mama's Choice"} and \eng{"Golden Drop"}.

\par Let's start with \eng{Mama's Choice}. A 5-litre jerrycan sells for an average of 12,000 FCFA. That is currently the lowest price on the market. The label says "100\% pure vegetable oil". However, our laboratory tests show a slightly higher acidity level than the national standard. Several consumers we interviewed complained that the oil foams excessively when heated and leaves a dark residue after frying plantains or fish. It does not carry the \eng{NAFDAC} certification number on the pack, though the manufacturer claims it is pending. There is no money-back guarantee offered by the distributor.

\par Now, \eng{Golden Drop}. The same 5-litre jerrycan costs 15,000 FCFA — that is 3,000 FCFA more expensive, a 25\% price difference. But the label displays a valid \eng{NAFDAC} registration number and the \eng{ISO 9001} quality mark. Our tests confirm low acidity, clear colour, and high smoke point — ideal for deep frying. Most importantly, the manufacturer offers a \textbf{one-year satisfaction guarantee}: if you are not satisfied with the quality, you can return the empty container to any authorised retailer for a full refund. Their customer service hotline is open six days a week.

\par So, the verdict. \eng{Mama's Choice} is cheaper at the checkout, no doubt. But \eng{Golden Drop} lasts longer — you use less oil per frying session because it doesn't foam away — it is safer for your health thanks to certified standards, and it comes with a risk-free guarantee. For the smart consumer who thinks \emph{total cost} and \emph{peace of mind}, \eng{Golden Drop} is clearly \textbf{the best buy}.

\par Remember, listeners: the cheapest price tag is not always the lowest cost. Until next week, shop smart!
\end{textebox}

\textbf{Glossaire de l'écoute.} \emph{a jerrycan} (n.) : un bidon ; \emph{acidity level} (n.) : taux d'acidité ; \emph{foam} (v.) : mousser ; \emph{residue} (n.) : résidu ; \emph{smoke point} (n.) : point de fumée (température où l'huile fume) ; \emph{pending} (adj.) : en attente ; \emph{hotline} (n.) : numéro vert / ligne d'assistance ; \emph{peace of mind} (expr.) : tranquillité d'esprit ; \emph{checkout} (n.) : caisse (supermarché).

\coursec{Compréhension : questions de types variés}

\courssub{Exercice 1 : Vrai ou Faux ? Justifiez par une phrase du texte.}
Écoutez (ou relisez) le script et décidez si les affirmations sont vraies (\eng{True}) ou fausses (\eng{False}). Citez la phrase (ou le passage) qui justifie votre réponse.

\begin{enumerate}
\item \eng{Mama's Choice is more expensive than Golden Drop.}
\item \eng{Both oils have the NAFDAC certification number on the pack.}
\item \eng{Golden Drop has a higher smoke point, making it better for deep frying.}
\item \eng{The manufacturer of Mama's Choice offers a one-year money-back guarantee.}
\item \eng{The host recommends Golden Drop because it is cheaper per litre.}
\end{enumerate}

\courssub{Exercice 2 : Questionnaire à choix multiples (QCM)}
Choisissez la bonne réponse (a, b ou c).

\begin{enumerate}
\item What is the price difference between the two 5-litre jerrycans?
\begin{itemize}
\item[a)] 2,000 FCFA
\item[b)] 3,000 FCFA
\item[c)] 5,000 FCFA
\end{itemize}
\item Why do some consumers complain about Mama's Choice?
\begin{itemize}
\item[a)] It smells bad.
\item[b)] It foams excessively and leaves a dark residue.
\item[c)] The jerrycan leaks.
\end{itemize}
\item What does the "one-year satisfaction guarantee" for Golden Drop allow?
\begin{itemize}
\item[a)] A free refill after one year.
\item[b)] A full refund if the empty container is returned.
\item[c)] A free repair of the stove.
\end{itemize}
\item According to the host, why does Golden Drop represent the "best buy"?
\begin{itemize}
\item[a)] It has the lowest price tag.
\item[b)] It lasts longer, is certified, and has a guarantee.
\item[c)] It is the most popular brand in Yaoundé.
\end{itemize}
\end{enumerate}

\courssub{Exercice 3 : Réponses courtes (Short answers)}
Répondez en anglais, par des phrases courtes ou des groupes nominaux.

\begin{enumerate}
\item \eng{What is the price of a 5-litre jerrycan of Golden Drop?}
\item \eng{Which certification does Golden Drop have that Mama's Choice lacks?}
\item \eng{How many days a week is the Golden Drop customer service hotline open?}
\item \eng{What does the host mean by "the cheapest price tag is not always the lowest cost"?}
\end{enumerate}

\begin{corrige}[Compréhension — The Best Buy]
\textbf{Exercice 1 (Vrai/Faux)}
\begin{enumerate}
\item \textbf{False.} — \eng{"A 5-litre jerrycan sells for an average of 12,000 FCFA... Golden Drop... costs 15,000 FCFA."} (Mama's Choice is cheaper).
\item \textbf{False.} — \eng{"It does not carry the NAFDAC certification number on the pack... Golden Drop... displays a valid NAFDAC registration number."}
\item \textbf{True.} — \eng{"...high smoke point — ideal for deep frying."}
\item \textbf{False.} — \eng{"There is no money-back guarantee offered by the distributor."} (C'est Golden Drop qui l'offre).
\item \textbf{False.} — \eng{"Golden Drop is clearly the best buy... lasts longer... safer... risk-free guarantee."} (Il le recommande pour la valeur globale, pas le prix le plus bas).
\end{enumerate}

\textbf{Exercice 2 (QCM)}
\begin{enumerate}
\item \textbf{b)} 3,000 FCFA (\eng{15,000 - 12,000 = 3,000}).
\item \textbf{b)} \eng{foams excessively and leaves a dark residue}.
\item \textbf{b)} \eng{return the empty container... for a full refund}.
\item \textbf{b)} \eng{lasts longer... certified standards... risk-free guarantee}.
\end{enumerate}

\textbf{Exercice 3 (Réponses courtes)}
\begin{enumerate}
\item \eng{15,000 FCFA.} (ou \eng{It costs 15,000 FCFA.})
\item \eng{NAFDAC registration number (and ISO 9001 quality mark).}
\item \eng{Six days a week.}
\item \eng{The total cost over time (durability, health, guarantee) matters more than the initial price.}
\end{enumerate}
\end{corrige}

\coursec{Production guidée : conseiller le meilleur achat}

\courssub{Méthode : structurer un avis de consommateur (oral ou écrit)}

Pour dire quel produit est le \eng{best buy}, suivez ce canevas en 4 étapes. Utilisez les connecteurs logiques et le vocabulaire de la leçon.

\begin{methode}[Comment présenter un "Best Buy" — Modèle en 4 étapes]
\begin{enumerate}
\item \textbf{Présenter les options} : \eng{I compared [Product A] and [Product B].} / \eng{There are two main options on the market: ...}
\item \textbf{Comparer sur les critères clés} (prix, qualité, durabilité, garantie) : \eng{Product A is cheaper than Product B, but...} / \eng{Although Product A costs less, Product B is more durable.}
\item \textbf{Donner le verdict} : \eng{In my opinion, Product B is the best buy.} / \eng{I would recommend Product B because...} / \eng{Product B offers better value for money.}
\item \textbf{Justifier avec des détails chiffrés} : \eng{It costs X FCFA and comes with a [durée] guarantee.} / \eng{It lasts [durée] longer than the competitor.}
\end{enumerate}
\end{methode}

\begin{exemplebox}[Modèle oral / écrit]
\eng{"I compared two brands of phone chargers: \textbf{PowerFast} and \textbf{QuickCharge}. \textbf{PowerFast} is cheaper — 3,000 FCFA against 5,000 FCFA. \textbf{However}, \textbf{QuickCharge} is more durable: it has a reinforced cable and a two-year guarantee, whereas PowerFast has no guarantee at all. \textbf{In the long run}, QuickCharge is the best buy because you don't have to replace it every three months. \textbf{That's why I recommend QuickCharge}."}
\end{exemplebox}

\courssub{Expressions utiles pour l'argumentation}

\begin{center}
\begin{tabularx}{\linewidth}{|X|X|}\hline
\rowcolor{popblueL} Fonction & Expressions anglaises \\ \hline
Introduire la comparaison & \eng{I compared... / Let's compare... / There are two options: ...} \\ \hline
Concéder (avantage adverse) & \eng{Admittedly, ... is cheaper / ... does have a lower price.} \\ \hline
Pivoter (argument décisif) & \eng{However, / But / On the other hand, / Nevertheless,} \\ \hline
Insister sur le critère clé & \eng{The key point is... / What matters most is... / Crucially,...} \\ \hline
Conclure / Recommander & \eng{Therefore, / So, / For these reasons,} \eng{... is the best buy.} \\ \hline
\end{tabularx}
\end{center}

\begin{experience}[Jeu de rôle : "Le conseiller conso"]
\textbf{Contexte :} Une émission de radio locale. \textbf{Rôles :} Animateur (Host) / Auditeur appelant (Caller) / Expert conso (Expert).
\par \textbf{Scénario :} L'auditeur appelle pour demander conseil entre deux produits (ex. : deux marques de ciment, deux forfaits internet, deux modèles de moto). L'expert compare à l'aide de la grille (prix, qualité, garantie, SAV) et donne son verdict \eng{best buy}. L'animateur gère le temps de parole.
\par \textbf{Contraintes :} Utiliser au moins 3 comparatifs, 1 superlatif, 1 expression de concession (\eng{although}), le mot \eng{value for money}.
\par \textbf{Déroulement :} 5 min préparation en groupes de 3 $\rightarrow$ 2 min jeu par groupe $\rightarrow$ feedback classe (vocabulaire, justesse, prononciation des chiffres).
\end{experience}

\begin{exercice}[Production écrite : Avis consommateur (80-100 mots)]
Vous avez testé deux marques de savon lessive (\eng{"WhiteClean"} et \eng{"BrightWash"}). Rédigez un court article pour le journal de l'école \eng{The Student Voice} afin de conseiller vos camarades. Utilisez la méthode en 4 étapes et le vocabulaire de la leçon.
\par \textbf{Données :}
\begin{itemize}
\item \eng{WhiteClean} : 1 kg = 1,200 FCFA. Pas de parfum. Tâche les mains. Pas de garantie.
\item \eng{BrightWash} : 1 kg = 1,500 FCFA. Parfum frais. Doux pour les mains. Garantie "satisfait ou remboursé" 30 jours.
\end{itemize}
\end{exercice}

\begin{corrige}[Production écrite — Modèle de réponse]
\eng{\textbf{Best Buy Laundry Soap: BrightWash Wins}}
\par \eng{I compared two popular laundry soaps: WhiteClean and BrightWash. WhiteClean is cheaper at 1,200 FCFA per kilo, while BrightWash costs 1,500 FCFA. However, WhiteClean has no fragrance and irritates the skin. BrightWash, on the other hand, has a fresh scent and is gentle on hands. Crucially, BrightWash offers a 30-day money-back guarantee, which WhiteClean does not. For these reasons, BrightWash is the best buy. It offers better value for money because it protects your clothes and your skin, and you take no risk thanks to the guarantee.}
\par (92 mots)
\end{corrige}

\begin{savaistu}[Le "Total Cost of Ownership" (Coût total de possession)]
En économie, le concept de \eng{Total Cost of Ownership (TCO)} explique pourquoi le \eng{best buy} n'est pas le moins cher à l'achat. Le TCO inclut : prix d'achat + coûts d'entretien + coûts de réparation + consommation (électricité, carburant) + durée de vie + valeur de revente. Une moto à 500,000 FCFA qui consomme peu et ne tombe jamais en panne a un TCO plus bas qu'une moto à 350,000 FCFA qui casse tous les mois. C'est exactement ce que l'émission \eng{The Best Buy} vous apprend à calculer intuitivement !
\end{savaistu}

\coursec{Lexique thématique : consommation et comparaison}

Dans la section précédente, nous avons découvert le thème de la leçon : comparer les prix et la qualité pour faire les \emph{meilleurs achats} (\eng{best buys}), que ce soit au marché de Mokolo, dans une boutique de Douala ou lorsqu'on choisit un service communautaire (eau, électricité, transport). Avant d'écouter le script \eng{The Best Buy}, il faut maîtriser le vocabulaire qui reviendra sans cesse dans l'enregistrement. Cette section construit ce lexique de manière organisée : quatre champs lexicaux, des expressions idiomatiques, et les faux amis qui piègent les francophones.

\courssub{Observer d'abord : un mini-dialogue au marché}

\begin{textebox}[At the market — “How much is this bag of rice?”]
\textbf{A:} Good morning, madam. How much is this bag of rice?

\textbf{B:} This one is 12,000 francs. It is a genuine local brand, very good quality.

\textbf{A:} That's expensive! The one next door is cheaper and just as good.

\textbf{B:} Be careful, my dear. The cheaper one is not durable — the grains break easily. You get what you pay for!

\textbf{A:} Hmm. Can you reduce the price a little? 11,000 francs?

\textbf{B:} All right, because you are a regular customer, I will give you a discount. Take it for 11,500 francs.

\textbf{A:} That's a real bargain! It is good value for money. Here is the money — and please give me a receipt.

\textbf{B:} Thank you! Come again!
\end{textebox}

\textbf{Glossaire.} \emph{bag} (n.m.) : sac ; \emph{genuine} (adj.) : authentique ; \emph{brand} (n.f.) : marque ; \emph{durable} (adj.) : durable, résistant ; \emph{grains} (n.m.pl.) : grains ; \emph{to reduce} (v.) : réduire, baisser ; \emph{discount} (n.m.) : remise ; \emph{bargain} (n.m.) : bonne affaire ; \emph{receipt} (n.m.) : reçu.

\textbf{Questions rapides.} 1) How much does the vendor first charge for the rice? 2) Why does she say the cheaper rice is not a good buy? 3) What discount does the customer finally get?

Observez : presque chaque réplique contient un mot de notre champ lexical. Nous allons maintenant classer ces mots.

\courssub{Le champ lexical des prix}

\begin{definition}[Les mots du prix]
\cle{price} (n.m.) : le prix (le montant affiché). \cle{cost} (n.m. et v.) : le coût / coûter. \cle{cheap} (adj.) : bon marché, pas cher. \cle{expensive} (adj.) : cher, coûteux. \cle{affordable} (adj.) : abordable (que l'on peut se permettre d'acheter). \cle{discount} (n.m.) : remise, réduction. \cle{bargain} (n.m.) : bonne affaire. \cle{to reduce} (v.) : réduire, baisser. \cle{to charge} (v.) : demander, facturer (un prix).
\end{definition}

\begin{exemplebox}[Les mots du prix en contexte]
\eng{What is the price of this shirt?} « Quel est le prix de cette chemise ? »

\eng{Petrol costs 750 francs a litre.} « L'essence coûte 750 francs le litre. »

\eng{Second-hand clothes are cheap at the market.} « Les fripes sont bon marché au marché. »

\eng{Imported rice is more expensive than local rice.} « Le riz importé est plus cher que le riz local. »

\eng{These school shoes are affordable for most families.} « Ces chaussures scolaires sont abordables pour la plupart des familles. »

\eng{The shopkeeper reduced the price by ten per cent.} « Le commerçant a baissé le prix de dix pour cent. »

\eng{How much do they charge for a bendskin ride?} « Combien demandent-ils pour une course en moto-taxi ? »
\end{exemplebox}

\begin{attention}[Price ou cost ?]
\eng{price} désigne le montant demandé par le vendeur ; \eng{cost} désigne ce que quelque chose coûte à l'acheteur. Comme verbe, on dit \eng{It costs 5,000 francs} — jamais \eng{It prices}. Attention aussi : l'adjectif français « cher » se traduit \eng{expensive}, pas \eng{dear} (qui signifie « cher » au sens affectif : \eng{my dear friend}).
\end{attention}

\courssub{Le champ lexical de la qualité}

\begin{definition}[Les mots de la qualité]
\cle{quality} (n.f.) : la qualité. \cle{durable} (adj.) : durable, qui dure longtemps. \cle{reliable} (adj.) : fiable, sur lequel on peut compter. \cle{genuine} (adj.) : authentique, véritable. \cle{fake} / \cle{counterfeit} (adj. et n.) : faux, contrefait / contrefaçon. \cle{brand} (n.f.) : marque. \cle{guarantee} / \cle{warranty} (n.f.) : garantie.
\end{definition}

\begin{exemplebox}[La qualité en contexte]
\eng{This radio is very reliable; I have used it for ten years.} « Cette radio est très fiable ; je l'utilise depuis dix ans. »

\eng{Beware of counterfeit phones in the streets of Douala.} « Méfiez-vous des téléphones contrefaits dans les rues de Douala. »

\eng{Is this a genuine product or a fake?} « Est-ce un produit authentique ou un faux ? »

\eng{The television has a two-year warranty.} « La télévision a une garantie de deux ans. »
\end{exemplebox}

\begin{attention}[Prononciation et orthographe]
\eng{guarantee} se prononce « ga-ran-ti » (le \emph{u} ne se prononce pas) ; \eng{receipt} se prononce « ri-ssit » (le \emph{p} est muet). Ne confondez pas \eng{quality} (la qualité) et \eng{quantity} (la quantité) : \eng{Quality is more important than quantity.}
\end{attention}

\courssub{Le champ lexical de l'achat}

\begin{definition}[Les acteurs et les gestes de l'achat]
\cle{customer} (n.) : client (d'un magasin). \cle{consumer} (n.) : consommateur (dans l'économie en général). \cle{seller} / \cle{vendor} (n.) : vendeur. \cle{to purchase} (v.) : acheter (registre soutenu). \cle{to compare} (v.) : comparer. \cle{value for money} (expr.) : rapport qualité-prix. \cle{receipt} (n.m.) : reçu.
\end{definition}

\begin{exemplebox}[Exemples traduits]
\eng{This phone is good value for money.} « Ce téléphone a un bon rapport qualité-prix. »

\eng{The market vendor gave me a discount.} « Le vendeur du marché m'a fait une remise. »

\eng{Consumers should compare prices before they purchase anything.} « Les consommateurs devraient comparer les prix avant d'acheter quoi que ce soit. »

\eng{Always keep your receipt in case the product is faulty.} « Gardez toujours votre reçu au cas où le produit serait défectueux. »
\end{exemplebox}

\courssub{Le champ lexical des services communautaires}

\begin{definition}[Community resources and services]
\cle{water supply} (n.f.) : l'approvisionnement en eau. \cle{electricity} (n.f.) : l'électricité. \cle{health centre} (n.m.) : centre de santé. \cle{public transport} (n.m.) : transport en commun. \cle{community radio} (n.f.) : radio communautaire.
\end{definition}

\begin{exemplebox}[Les services en contexte]
\eng{The water supply in our quarter is not reliable.} « L'approvisionnement en eau de notre quartier n'est pas fiable. »

\eng{The new health centre charges affordable fees.} « Le nouveau centre de santé pratique des tarifs abordables. »

\eng{Public transport is cheaper than taxis for students.} « Le transport en commun est moins cher que les taxis pour les élèves. »

\eng{Our community radio announces market prices every morning.} « Notre radio communautaire annonce les prix du marché chaque matin. »
\end{exemplebox}

\courssub{Expressions utiles pour comparer et conseiller}

\begin{aretenir}[Quatre expressions à réemployer à l'oral]
\begin{itemize}
\item \eng{It's worth the price.} « Ça vaut son prix. »
\item \eng{You get what you pay for.} « On en a pour son argent » (un produit très bon marché est souvent de mauvaise qualité).
\item \eng{It's a real bargain!} « C'est une vraie bonne affaire ! »
\item \eng{It's good value for money.} « C'est d'un bon rapport qualité-prix. »
\end{itemize}
\end{aretenir}

\begin{center}
\begin{tabularx}{\linewidth}{|X|X|X|}
\hline
\rowcolor{popblueL} English & Français & Remarque \\
\hline
price & prix & \eng{What is the price?} \\
\hline
bargain & bonne affaire & \eng{It's a real bargain!} \\
\hline
quality & qualité & ne pas confondre avec \eng{quantity} \\
\hline
guarantee / warranty & garantie & \eng{guarantee} : \emph{u} muet ; \eng{warranty} : sans \emph{u} ni \emph{p} \\
\hline
discount & remise & \eng{to give a discount} \\
\hline
receipt & reçu & prononcé « ri-ssit » (\emph{p} muet) \\
\hline
value for money & rapport qualité-prix & expression invariable \\
\hline
\end{tabularx}
\end{center}

\begin{popfigure}
\begin{center}
\begin{tikzpicture}
\node[draw=popblue, fill=popblueL, rounded corners, align=center, font=\bfseries] (root) at (0,0) {CONSUMER\\VOCABULARY};
\node[draw=popgreen, fill=popgreenL, rounded corners, align=center] (price) at (-4.2,2.2) {PRICES\\price, cost, cheap,\\discount, to charge};
\node[draw=poporange, fill=poporangeL, rounded corners, align=center] (quality) at (4.2,2.2) {QUALITY\\durable, reliable,\\genuine, fake};
\node[draw=poppurple, fill=poppurpleL, rounded corners, align=center] (buy) at (-4.2,-2.2) {BUYING\\customer, vendor,\\receipt, bargain};
\node[draw=poppink, fill=poppinkL, rounded corners, align=center] (serv) at (4.2,-2.2) {SERVICES\\water supply,\\health centre,\\public transport};
\draw[-{Stealth}, thick, popblue] (root) -- (price);
\draw[-{Stealth}, thick, popblue] (root) -- (quality);
\draw[-{Stealth}, thick, popblue] (root) -- (buy);
\draw[-{Stealth}, thick, popblue] (root) -- (serv);
\end{tikzpicture}
\end{center}
\legende{Carte mentale : les quatre familles du lexique de la consommation.}
\end{popfigure}

\courssub{Faux amis et pièges des francophones}

\begin{attention}[« Économique » n'est pas « economic »]
\eng{economic} signifie « relatif à l'économie » : \eng{the economic situation of the country} « la situation économique du pays ». Pour dire « qui fait économiser de l'argent », utilisez \eng{economical} ou \eng{money-saving} : \eng{This stove is very economical.} « Ce réchaud est très économique. » De même, « magasin » se dit \eng{shop} (Grande-Bretagne) ou \eng{store} (États-Unis) — jamais \eng{magazine}, qui signifie « revue, magazine ».
\end{attention}

\begin{center}
\begin{tabularx}{\linewidth}{|X|X|X|}
\hline
\rowcolor{poporangeL} Français & Piège (faux ami) & Traduction correcte \\
\hline
économique (qui épargne) & economic & economical / money-saving \\
\hline
économique (de l'économie) & --- & economic \\
\hline
magasin & magazine & shop / store \\
\hline
cher (coûteux) & dear & expensive \\
\hline
demander un prix & to ask a price & to charge \\
\hline
\end{tabularx}
\end{center}

\courssub{Méthode pour mémoriser durablement ce lexique}

\begin{methode}[Apprendre par familles de mots]
\begin{enumerate}
\item \textbf{Regroupez} les mots par famille : \eng{buy} (v.) $\rightarrow$ \eng{buyer} (acheteur) $\rightarrow$ \eng{buying} (achat) $\rightarrow$ \eng{best buy} (meilleur achat). De même : \eng{sell} $\rightarrow$ \eng{seller} $\rightarrow$ \eng{sale} ; \eng{consume} $\rightarrow$ \eng{consumer} $\rightarrow$ \eng{consumption}.
\item \textbf{Classez} chaque nouveau mot dans l'une des quatre familles de la carte mentale (prix, qualité, achat, services).
\item \textbf{Réemployez} chaque mot dans une phrase personnelle liée à votre vie : \eng{My mother always compares prices at Nkoulouloun market.}
\item \textbf{Testez-vous} en couvrant la colonne française du tableau et en retrouvant la traduction.
\end{enumerate}
\end{methode}

\begin{exoresolu}[Complétez avec le mot qui convient]
Énoncé. Complétez chaque phrase avec un mot de la liste : \emph{bargain, warranty, affordable, charge, counterfeit, receipt}.

1. This bag is not genuine; it is a \ldots\ product.

2. The school fees are \ldots\ for most parents in our village.

3. How much does the taxi driver \ldots\ from Buea to Mile 17?

4. Keep your \ldots\ ; you may need it if the radio breaks down.

5. This blender has a one-year \ldots\ .

6. Three dresses for 5,000 francs? What a \ldots\ !
\tcblower
Solution détaillée. 1. \textbf{counterfeit} — le mot \eng{genuine} (authentique) appelle son contraire : « contrefait ». 2. \textbf{affordable} — « abordable », que les parents peuvent payer. 3. \textbf{charge} — « combien demande le chauffeur de taxi » ; \eng{to charge} = demander un prix. 4. \textbf{receipt} — le reçu sert de preuve d'achat. 5. \textbf{warranty} — une garantie d'un an sur un appareil. 6. \textbf{bargain} — trois robes à ce prix, c'est une bonne affaire (\eng{What a bargain!}).
\end{exoresolu}

\begin{exercice}[À vous de jouer]
Traduisez en anglais : 1) Ce téléphone a un bon rapport qualité-prix. 2) Le vendeur m'a fait une remise de cinq pour cent. 3) Comparez toujours les prix avant d'acheter. 4) Ce four est très économique. Puis écrivez deux phrases personnelles avec \eng{bargain} et \eng{reliable}.
\end{exercice}

\begin{corrige}[Exercice « À vous de jouer »]
1) \eng{This phone is good value for money.} 2) \eng{The vendor gave me a five per cent discount.} 3) \eng{Always compare prices before you buy.} 4) \eng{This oven is very economical.} (et non \eng{economic} !). Phrases personnelles : réponses libres, par exemple \eng{These shoes were a real bargain.} / \eng{Our community radio is very reliable.}
\end{corrige}

Ce lexique est désormais votre boîte à outils : dans la prochaine section, nous verrons les \emph{stratégies d'écoute} qui vous permettront de repérer ces mots-clés, les chiffres et les noms propres dans l'enregistrement \eng{The Best Buy}.

\coursec{Stratégies et méthode d'écoute}

Dans la section précédente, nous avons enrichi notre vocabulaire de la consommation : \eng{price}, \eng{quality}, \eng{bargain}, \eng{cheaper than}, \eng{value for money}... Mais connaître des mots ne suffit pas : il faut encore être capable de les \cle{reconnaître à l'oral}, dans le flux rapide d'un enregistrement. C'est précisément l'objet de cette section : apprendre une \cle{méthode d'écoute} efficace, celle qui vous permettra de réussir l'épreuve de \emph{listening comprehension} au Baccalauréat. Suivez le guide, étape par étape.

\courssub{Pourquoi une méthode ? L'erreur de l'écoute « passive »}

Beaucoup d'élèves écoutent un enregistrement comme ils écoutent une chanson : ils attendent que le sens « vienne tout seul ». Résultat : à la première phrase incomprise, ils paniquent, se décrochent, et perdent le fil de tout le passage. L'écoute efficace est au contraire une activité \cle{active} et \cle{ciblée} : on sait \emph{ce qu'on cherche} avant même d'appuyer sur « play ».

\begin{methode}[La règle d'or de l'écoute]
Ne cherche \textbf{jamais} à comprendre chaque mot. Même un anglophone natif n'entend pas chaque mot : il capte les \cle{mots porteurs de sens} et reconstruit le reste. Concentre-toi sur :
\begin{enumerate}
\item les \textbf{noms} (produits, lieux, personnes) : \eng{rice}, \eng{Mokolo market}, \eng{Mrs Bih} ;
\item les \textbf{chiffres et prix} : \eng{fifteen thousand francs}, \eng{50 kilos}, \eng{20\% discount} ;
\item les \textbf{adjectifs et structures de comparaison} : \eng{cheaper}, \eng{better quality}, \eng{the best buy} ;
\item les \textbf{verbes d'action} : \eng{costs}, \eng{recommend}, \eng{choose}, \eng{save}.
\end{enumerate}
Les petits mots (\eng{the}, \eng{of}, \eng{and}, \eng{you know}...) peuvent être ignorés sans dommage.
\end{methode}

\courssub{Étape 1 — Avant l'écoute : lire, prédire, anticiper}

L'écoute commence \emph{avant} l'écoute. Pendant que l'examinateur ou le professeur distribue les questions, tu disposes d'un temps précieux : utilise-le.

\begin{methode}[Les trois actions à mener avant la première écoute]
\begin{enumerate}
\item \textbf{Lis toutes les questions} (vrai/faux, QCM, questions ouvertes). Elles te révèlent le thème probable et les informations à repérer. Si une question demande \eng{How much does the bag of rice cost at Shop B?}, tu sais qu'il faudra noter des \cle{prix} et distinguer \cle{deux magasins}.
\item \textbf{Prédis le contenu}. À partir du titre et des questions, formule une hypothèse : « Il s'agit probablement d'un dialogue entre un client et un vendeur qui comparent des produits. » La prédiction prépare ton cerveau : les mots attendus seront reconnus plus vite.
\item \textbf{Repère les mots attendus}. Souligne dans les questions les mots-clés : chiffres, noms propres, noms de produits. Prépare mentalement leur prononciation : \eng{Mokolo} se dit « mo-ko-lo », \eng{fifteen} se dit « fif-tine ».
\end{enumerate}
\end{methode}

\begin{exemplebox}[Prédire à partir d'une question]
Question donnée avant l'écoute : \eng{Which shop does the speaker finally recommend, and why?}

Prédictions possibles : le texte comparera \textbf{au moins deux} magasins ; il contiendra des adjectifs de comparaison (\eng{cheaper}, \eng{fresher}, \eng{better}) ; la réponse se trouvera probablement \textbf{vers la fin} du texte (la recommandation est une conclusion). Tu sais donc qu'il faut rester concentré jusqu'au bout.
\end{exemplebox}

\courssub{Étape 2 — Première écoute : saisir le thème général (global listening)}

La première écoute a un objectif unique et limité : comprendre \cle{de quoi parle le texte} et \cle{qui parle}. Ne t'occupe pas encore des détails.

\begin{propriete}[Les deux écoutes au Baccalauréat]
Au Bac, le texte de \emph{listening comprehension} est lu \textbf{deux fois}. Chaque écoute a sa fonction :
\begin{itemize}
\item \textbf{Première écoute} (\emph{global listening}) : identifier le \cle{thème général}, la \cle{situation} (marché, magasin, radio...) et les \cle{personnes} qui parlent (vendeuse, client, présentateur...).
\item \textbf{Deuxième écoute} (\emph{selective listening}) : relever les \cle{détails} (chiffres, prix, noms propres) et \cle{vérifier} les réponses déjà esquissées.
\end{itemize}
Ne réponds jamais à toutes les questions dès la première écoute : tu risquerais de rater la suite du texte en écrivant.
\end{propriete}

Pendant la première écoute, pose-toi trois questions mentales :

\begin{center}
\begin{tabularx}{\linewidth}{|X|X|X|}
\hline
\rowcolor{popblueL} Question mentale & English & Réponse attendue \\
\hline
De quoi parle-t-on ? & What is the topic? & \eng{Comparing prices of rice in two shops.} \\
\hline
Qui parle ? & Who is speaking? & \eng{A market woman and a customer.} \\
\hline
Où et dans quelle situation ? & Where / in what situation? & \eng{At Mokolo market, Yaoundé.} \\
\hline
\end{tabularx}
\end{center}

\courssub{Étape 3 — Deuxième écoute : noter les détails (selective listening)}

La deuxième écoute est une \cle{écoute sélective} : tu ne cherches que ce que les questions demandent. C'est le moment de la \cle{prise de notes}.

\begin{propriete}[Technique de prise de notes : abréviations et symboles]
Écris vite, sans phrases complètes. Utilise :
\begin{itemize}
\item les \textbf{abréviations} : \eng{kg} (kilos), \eng{F} ou \eng{FCFA} (francs), \eng{£} (pounds), \eng{\%} (per cent), \eng{min} (minutes), \eng{qty} (quantity) ;
\item les \textbf{flèches de comparaison} : \eng{A > B} (A est plus cher / meilleur que B), \eng{A < B}, \eng{A = B}, \eng{A $\rightarrow$ 15 000 F} (A coûte 15 000 F) ;
\item les \textbf{initiales} des noms propres : \eng{Mrs B.}, \eng{Shop A / Shop B}.
\end{itemize}
Exemple de notes : \eng{rice A: 50kg $\rightarrow$ 15 000F; B: 50kg $\rightarrow$ 13 500F but old stock; A > B quality}.
\end{propriete}

\begin{attention}[Piège majeur : les chiffres qui se ressemblent]
À l'oral, les nombres en \eng{-teen} et en \eng{-ty} se confondent facilement : \eng{thirteen} / \eng{thirty}, \eng{fifteen} / \eng{fifty}, \eng{sixteen} / \eng{sixty}. Deux indices pour les distinguer :
\begin{enumerate}
\item \textbf{L'accent tonique} : \eng{fifTEEN} (accent sur la 2\textsuperscript{e} syllabe, « fif-TINE ») contre \eng{FIFty} (accent sur la 1\textsuperscript{re}, « FIF-ti »).
\item \textbf{Le contexte du prix} : un sac de riz de 50 kg ne peut pas coûter 50 francs ; si tu hésites entre \eng{fifteen thousand} et \eng{fifty thousand}, demande-toi quel montant est \emph{plausible}.
\end{enumerate}
Vérifie toujours le chiffre à la deuxième écoute.
\end{attention}

\begin{exemplebox}[Exemples traduits : prix et quantités]
\eng{The 50-kilo bag costs fifteen thousand francs.} « Le sac de 50 kilos coûte quinze mille francs. » (\eng{fifteen} $\neq$ \eng{fifty} !)

\eng{Shop B sells it for thirteen thousand, five hundred francs.} « Le magasin B le vend treize mille cinq cents francs. »

\eng{You get a ten per cent discount if you buy two bags.} « Vous bénéficiez d'une remise de dix pour cent si vous achetez deux sacs. »

\eng{This one is cheaper, but the quality is lower.} « Celui-ci est moins cher, mais la qualité est inférieure. »
\end{exemplebox}

\courssub{Étape 4 — Vérification : répondre et justifier}

Après la deuxième écoute, il reste souvent une courte pause. C'est le moment de \cle{finaliser} tes réponses.

\begin{methode}[Vérifier ses réponses en quatre points]
\begin{enumerate}
\item \textbf{Réponds à toutes les questions}, même par une hypothèse : une case vide est un point perdu, une supposition logique peut être juste.
\item \textbf{Justifie avec tes notes} : chaque réponse doit pouvoir s'appuyer sur un élément noté (un chiffre, un mot-clé). Si tu as noté \eng{B: 13 500F}, la réponse à \eng{Which shop is cheaper?} est nécessairement \eng{Shop B}.
\item \textbf{Contrôle la cohérence} : si le texte recommande le magasin A et que ta réponse au vrai/faux dit le contraire, l'une des deux réponses est fausse.
\item \textbf{Réécoute mentalement les passages difficiles} : en classe, demande au professeur de repasser le passage ; seul, réécoute l'extrait en te concentrant uniquement sur le mot manquant.
\end{enumerate}
\end{methode}

\courssub{Schéma récapitulatif : la méthode en six étapes}

\begin{popfigure}
\begin{center}
\begin{tikzpicture}[
node distance=0.55cm,
etape/.style={rectangle, rounded corners=3pt, draw=popblue, fill=popblueL, text=popdark, align=center, minimum width=4.6cm, minimum height=0.85cm, font=\small},
fleche/.style={-{Stealth[length=2.5mm]}, thick, poporange}
]
\node[etape] (e1) {\textbf{1. Read} the questions};
\node[etape, below=of e1] (e2) {\textbf{2. Predict} the content};
\node[etape, below=of e2] (e3) {\textbf{3. Listen 1} $\rightarrow$ theme, speakers};
\node[etape, below=of e3, align=center] (e4){\textbf{4. Listen 2} $\rightarrow$ details:\\ figures, prices, names};
\node[etape, below=of e4] (e5) {\textbf{5. Answer} the questions};
\node[etape, below=of e5] (e6) {\textbf{6. Check} with your notes};
\draw[fleche] (e1) -- (e2);
\draw[fleche] (e2) -- (e3);
\draw[fleche] (e3) -- (e4);
\draw[fleche] (e4) -- (e5);
\draw[fleche] (e5) -- (e6);
\end{tikzpicture}
\end{center}
\legende{Organigramme de la méthode d'écoute : six étapes, de la lecture des questions à la vérification.}
\end{popfigure}

\courssub{Exemple traité en classe : repérage guidé}

Mettons la méthode en application sur un court extrait, comme nous le ferions en classe. Lis d'abord les questions (étape 1), prédis le contenu (étape 2), puis « écoute » le script ci-dessous lu par le professeur.

\textbf{Questions données avant l'écoute :}
\begin{enumerate}
\item What product are the speakers comparing?
\item How much does the bag cost at Central Store?
\item Which shop does the customer finally choose, and why?
\end{enumerate}

\begin{textebox}[Listening script — “At the provision store”]
\textbf{Customer:} Good morning, madam. I need a bag of rice for my restaurant. How much is the fifty-kilo bag?

\textbf{Shopkeeper:} Here at Central Store, it is fifteen thousand francs. It is good quality rice from Ndop.

\textbf{Customer:} Fifteen thousand? At New Deal Shop, they told me thirteen thousand five hundred.

\textbf{Shopkeeper:} Yes, but be careful. Their rice is old stock. It breaks when you cook it. Mine arrived last week. Look at the grains: they are long and clean.

\textbf{Customer:} Hmm, you are right. But fifteen thousand is expensive for me.

\textbf{Shopkeeper:} If you buy two bags, I give you a ten per cent discount. That is thirteen thousand five hundred per bag, for better quality.

\textbf{Customer:} That is a good bargain. I will take two bags, please.

\textbf{Shopkeeper:} A wise choice, sir. Your customers will taste the difference.
\end{textebox}

\textbf{Glossaire :} \emph{shopkeeper} (n., commerçant) ; \emph{old stock} (n., vieux stock) ; \emph{grain} (n., grain) ; \emph{discount} (n., remise) ; \emph{bargain} (n., bonne affaire) ; \emph{wise} (adj., judicieux).

\begin{exoresolu}[Repérage guidé des prix et du produit recommandé]
Énoncé. À l'aide du script ci-dessus, réponds aux trois questions posées avant l'écoute, en justifiant chaque réponse par une note (chiffre ou mot-clé).
\tcblower
Solution détaillée.
\begin{enumerate}
\item \eng{They are comparing the price and quality of a 50-kilo bag of rice.} Note : \eng{50kg bag, rice from Ndop}.
\item \eng{It costs fifteen thousand francs (15 000 F) at Central Store.} Note : \eng{Central Store $\rightarrow$ 15 000F}. Attention : c'est bien \eng{fifteen} (15 000), pas \eng{fifty} ; le contexte le confirme, car le client trouve le prix élevé par rapport à 13 500 F.
\item \eng{The customer chooses Central Store, because the rice is better quality (new stock, long clean grains) and the ten per cent discount brings the price down to 13 500 F per bag.} Notes : \eng{New Deal = old stock; Central + 10\% discount $\rightarrow$ 13 500F; quality A > B}.
\end{enumerate}
Remarque la stratégie du vendeur : le prix affiché est plus élevé, mais la remise et la qualité renversent la comparaison. À l'écoute, il faut donc suivre \emph{toute} la négociation avant de conclure.
\end{exoresolu}

\begin{attention}[Erreurs fréquentes des francophones à l'écoute]
\begin{itemize}
\item \textbf{S'arrêter sur un mot inconnu.} Pendant que tu réfléchis à \eng{wise}, tu rates trois phrases. Note le son (« ouaïze ») et continue ; le contexte expliquera le mot.
\item \textbf{Confondre les chiffres} : \eng{fifteen} / \eng{fifty}, \eng{thirteen} / \eng{thirty}. Vérifie l'accent tonique et la vraisemblance du prix.
\item \textbf{Traduire mentalement mot à mot.} La traduction prend trop de temps ; comprends directement en anglais à partir des mots-clés.
\item \textbf{Écrire des phrases complètes} pendant l'écoute : utilise les abréviations (\eng{kg}, \eng{F}, \eng{\%}) et les flèches (\eng{A > B}).
\item \textbf{Se fier à la première information entendue} : le premier prix cité n'est pas toujours le prix final (remises, négociation). Écoute jusqu'au bout.
\end{itemize}
\end{attention}

\begin{exercice}[À toi de jouer]
Relis le script “At the provision store” à voix haute avec un camarade (ou écoute-le lu par le professeur), puis :
\begin{enumerate}
\item Note uniquement avec des abréviations et des flèches : les deux prix, la remise et la comparaison de qualité.
\item Réponds par \eng{True} ou \eng{False} : \eng{The rice at New Deal Shop is new stock.} / \eng{With the discount, one bag costs 13 500 F.}
\item Rédige une phrase de recommandation avec \eng{I recommend... because...}.
\end{enumerate}
\end{exercice}

\begin{corrige}[Exercice « À toi de jouer »]
\begin{enumerate}
\item Notes possibles : \eng{Central: 15 000F; New Deal: 13 500F (old stock); 2 bags $\rightarrow$ -10\% $\rightarrow$ 13 500F/bag; quality Central > New Deal}.
\item \eng{The rice at New Deal Shop is new stock.} $\rightarrow$ \textbf{False} (it is old stock). \eng{With the discount, one bag costs 13 500 F.} $\rightarrow$ \textbf{True} (15 000 F moins 10 \% = 13 500 F).
\item Modèle : \eng{I recommend Central Store because the rice is better quality and, with the ten per cent discount, it costs the same price as at New Deal Shop.} « Je recommande Central Store parce que le riz est de meilleure qualité et, avec la remise de dix pour cent, il coûte le même prix qu'à New Deal Shop. »
\end{enumerate}
\end{corrige}

\begin{aretenir}[L'essentiel de la méthode d'écoute]
\begin{itemize}
\item \textbf{Avant} : lis les questions, prédis le contenu, repère les mots attendus.
\item \textbf{Écoute 1} : thème général et locuteurs uniquement.
\item \textbf{Écoute 2} : détails — chiffres, prix, noms propres — avec abréviations (\eng{kg}, \eng{F}, \eng{\%}) et flèches (\eng{A > B}).
\item \textbf{Après} : réponds à tout, justifie avec tes notes, vérifie la cohérence.
\item Méfie-toi des couples \eng{-teen} / \eng{-ty} et ne t'accroche jamais à un mot inconnu.
\end{itemize}
\end{aretenir}

\coursec{Texte support : script d'écoute « The Best Buy »}

Dans la section précédente, nous avons appris à préparer notre écoute : anticiper le thème grâce au titre, repérer les mots-clés, se concentrer sur les chiffres et les noms propres. Le moment est venu de mettre ces stratégies en pratique sur un enregistrement authentique. Le script ci-dessous est un \cle{reportage radio} : c'est le type de document sonore le plus fréquent dans les épreuves de compréhension orale. L'enseignant le lit deux fois (ou diffuse l'enregistrement) ; vous écoutez d'abord \emph{sans prendre de notes}, puis une seconde fois \emph{avec une grille de relevé}.

\courssub{Le script d'écoute : « Consumer Watch »}

\begin{textebox}[Consumer Watch — Community Radio, Buea]
Good morning, listeners! Welcome to \emph{Consumer Watch} on Community Radio. Today we compare three shops in Buea where you can buy a 50-kilo bag of rice.

At Mile 17 Market, the bag costs 12,000 francs, but some customers complain \textbf{that} the rice is mixed with stones. At the Central Supermarket, the same bag costs 15,000 francs; it is a well-known brand with a guarantee of quality. Finally, at Mama Ngo's store near the health centre, the bag costs 13,500 francs and the rice is clean and \textbf{holds up well in cooking}.

Our reporter interviewed ten customers: seven said Mama Ngo's store offers the best value for money.

Remember, listeners: the cheapest product is not always the best buy. Compare prices, check the quality, and always ask for a receipt!
\end{textebox}

\textbf{Glossaire du script}\par

\begin{center}
\begin{tabularx}{\linewidth}{|l|l|X|}
\hline
\rowcolor{popblueL} \textbf{Mot anglais} & \textbf{Nature} & \textbf{Traduction française} \\
\hline
listeners & nom pluriel & auditeurs, auditrices \\
\hline
to complain (about / that) & verbe & se plaindre (de / que) \\
\hline
stones & nom pluriel & cailloux, pierres \\
\hline
brand & nom & marque (d'un produit) \\
\hline
guarantee & nom & garantie \\
\hline
holds up well in cooking & locution & tient bien à la cuisson, ne se défait pas \\
\hline
value for money & expression & rapport qualité-prix \\
\hline
receipt & nom & reçu, ticket de caisse \\
\hline
\end{tabularx}
\end{center}

\begin{attention}[Prononciation et faux amis]
\begin{itemize}
\item \eng{receipt} se prononce « ri-site » (API : /rɪˈsiːt/) : le \eng{p} est \textbf{muet}. Ne le prononcez jamais.
\item \eng{durable} (non utilisé dans le script final, remplacé par \eng{holds up well in cooking}) est un \textbf{faux ami} partiel : en anglais, il signifie « qui dure longtemps, solide », et non « soutenable » au sens écologique (qui se dit \eng{sustainable}).
\item \eng{to complain} se construit avec \eng{about} (nom) ou \eng{that} (proposition) : \eng{Customers complain about the quality.} / \eng{Customers complain that the rice is mixed with stones.} — jamais \eng{complain the quality}.
\item \eng{stones} se prononce « steunz » (/stəʊnz/) ; attention au \eng{s} final sonore \eng{/z/}.
\item \eng{health centre} : orthographe britannique (\eng{centre}), standard au Cameroun.
\end{itemize}
\end{attention}

\courssub{Première écoute : compréhension globale}

Lors de la \cle{première écoute}, vous ne cherchez aucun détail. Votre seul objectif est de répondre à trois questions générales :

\begin{enumerate}
\item \textbf{What is the topic?} (Quel est le thème ?) — \emph{Comparing the price and quality of rice in three shops in Buea.}
\item \textbf{Who is speaking?} (Qui parle ?) — \emph{A radio presenter on a consumer programme.}
\item \textbf{Where does the scene take place?} (Où se passe la scène ?) — \emph{In Buea, on Community Radio; the shops mentioned are Mile 17 Market, the Central Supermarket and Mama Ngo's store.}
\end{enumerate}

\begin{methode}[Réussir l'écoute globale]
\begin{enumerate}
\item Écoutez le \textbf{titre} et la \textbf{première phrase} : ils annoncent presque toujours le thème (ici : \eng{Today we compare three shops...}).
\item Identifiez les \textbf{marqueurs de situation} : \eng{Good morning, listeners!}, \eng{Welcome to...} indiquent une émission de radio.
\item Repérez les \textbf{noms propres} (Buea, Mile 17, Mama Ngo's) : ils donnent le lieu et les personnages.
\item Ne paniquez pas sur les mots inconnus : le sens général suffit à ce stade.
\end{enumerate}
\end{methode}

\courssub{Deuxième écoute : écoute sélective et relevé d'informations}

À la \cle{deuxième écoute}, vous remplissez une grille. Concentrez-vous sur les \textbf{chiffres} (les prix) et les \textbf{adjectifs de qualité}. Voici le tableau à compléter pendant l'écoute :

\begin{popfigure}
\begin{center}
\begin{tikzpicture}
\node[draw=popblue, fill=popblueL, rounded corners, minimum width=3.2cm, minimum height=0.8cm, font=\bfseries] (h1) at (0,0) {Shop};
\node[draw=popblue, fill=popblueL, rounded corners, minimum width=3.2cm, minimum height=0.8cm, font=\bfseries] (h2) at (3.6,0) {Price};
\node[draw=popblue, fill=popblueL, rounded corners, minimum width=3.2cm, minimum height=0.8cm, font=\bfseries] (h3) at (7.2,0) {Quality};
\node[draw=popblue, fill=popblueL, rounded corners, minimum width=3.4cm, minimum height=0.8cm, font=\bfseries] (h4) at (10.9,0) {Verdict};
\node[draw=popline, fill=popcream, rounded corners, minimum width=3.2cm, minimum height=0.8cm, align=center] at (0,-1) {Mile 17 Market};
\node[draw=popline, fill=popcream, rounded corners, minimum width=3.2cm, minimum height=0.8cm, align=center] at (3.6,-1) {12,000 F};
\node[draw=popline, fill=popcream, rounded corners, minimum width=3.2cm, minimum height=0.8cm, align=center] at (7.2,-1) {poor (stones)};
\node[draw=popline, fill=popcream, rounded corners, minimum width=3.4cm, minimum height=0.8cm, align=center] at (10.9,-1) {risky};
\node[draw=popline, fill=popcream, rounded corners, minimum width=3.2cm, minimum height=0.8cm, align=center] at (0,-2) {Central Supermarket};
\node[draw=popline, fill=popcream, rounded corners, minimum width=3.2cm, minimum height=0.8cm, align=center] at (3.6,-2) {15,000 F};
\node[draw=popline, fill=popcream, rounded corners, minimum width=3.2cm, minimum height=0.8cm, align=center] at (7.2,-2) {high (brand)};
\node[draw=popline, fill=popcream, rounded corners, minimum width=3.4cm, minimum height=0.8cm, align=center] at (10.9,-2) {expensive};
\node[draw=popline, fill=popgreenL, rounded corners, minimum width=3.2cm, minimum height=0.8cm, align=center] at (0,-3) {Mama Ngo's store};
\node[draw=popline, fill=popgreenL, rounded corners, minimum width=3.2cm, minimum height=0.8cm, align=center] at (3.6,-3) {13,500 F};
\node[draw=popline, fill=popgreenL, rounded corners, minimum width=3.2cm, minimum height=0.8cm, align=center] at (7.2,-3) {good (clean)};
\node[draw=popline, fill=popgreenL, rounded corners, minimum width=3.4cm, minimum height=0.8cm, align=center, font=\bfseries] at (10.9,-3) {best buy};
\end{tikzpicture}
\end{center}
\legende{Grille de relevé complétée : les trois boutiques comparées}
\end{popfigure}

\begin{aretenir}[Les informations à relever en écoute sélective]
\begin{itemize}
\item \textbf{Les trois prix} : 12,000 F (Mile 17), 15,000 F (Central Supermarket), 13,500 F (Mama Ngo's).
\item \textbf{Les trois lieux} : Mile 17 Market, Central Supermarket, Mama Ngo's store near the health centre.
\item \textbf{L'avis des clients} : seven out of ten customers prefer Mama Ngo's store.
\end{itemize}
\end{aretenir}

\courssub{Le mot-clé du script : value for money}

\begin{definition}[Value for money]
\eng{Value for money} signifie \emph{what you get compared with what you pay} — ce que tu reçois par rapport à ce que tu paies. En français : le \textbf{rapport qualité-prix}. Un produit offre \eng{good value for money} quand sa qualité est élevée par rapport à son prix, même s'il n'est pas le moins cher.
\end{definition}

C'est exactement la conclusion du reportage : le riz de Mama Ngo n'est ni le moins cher ni le plus cher, mais c'est le \emph{meilleur achat} (\eng{the best buy}) parce que son rapport qualité-prix est le plus avantageux.

\begin{exemplebox}[Phrases-clés du script, traduites]
\eng{The cheapest product is not always the best buy.} « Le produit le moins cher n'est pas toujours le meilleur achat. »

\eng{Always ask for a receipt.} « Demandez toujours un reçu. »

\eng{Mama Ngo's store offers the best value for money.} « La boutique de Mama Ngo offre le meilleur rapport qualité-prix. »

\eng{Some customers complain that the rice is mixed with stones.} « Certains clients se plaignent que le riz est mélangé à des cailloux. »
\end{exemplebox}

\begin{attention}[Pièges de traduction]
\begin{itemize}
\item \eng{the best buy} ne se traduit pas mot à mot par « le meilleur achat » au sens d'acheter : c'est un \textbf{nom} qui désigne le produit le plus avantageux. De même, \eng{a good buy} = « une bonne affaire ».
\item \eng{ask for} = demander \emph{quelque chose}. Ne dites jamais \eng{ask a receipt} : la préposition \eng{for} est obligatoire.
\item Ordre des mots : \eng{a 50-kilo bag of rice} = « un sac de riz de 50 kilos ». En anglais, la mesure précède avec un trait d'union et \eng{kilo} reste au singulier.
\end{itemize}
\end{attention}

\courssub{Reformulation : résumer le message}

Après les deux écoutes, l'élève doit être capable de \cle{résumer} le script en deux ou trois phrases, avec ses propres mots, en réemployant le lexique du thème.

\begin{exoresolu}[Summarise the radio report in two or three sentences]
Énoncé : Using your notes, summarise the « Consumer Watch » report. Use the words \eng{compare}, \eng{value for money} and \eng{receipt}.
\tcblower
Solution détaillée : un bon résumé reprend le thème, les données essentielles (prix et lieux) et la conclusion du journaliste, sans recopier le texte.

Modèle de réponse : \eng{The radio programme compares the price and quality of a 50-kilo bag of rice in three shops in Buea. Mile 17 Market is the cheapest at 12,000 francs, but the quality is poor, while the Central Supermarket is the most expensive at 15,000 francs. Seven customers out of ten say Mama Ngo's store, at 13,500 francs, offers the best value for money, and the presenter advises listeners to compare prices, check quality and always ask for a receipt.}

« L'émission compare le prix et la qualité d'un sac de riz de 50 kilos dans trois boutiques de Buea. Mile 17 est le moins cher à 12 000 francs, mais la qualité est mauvaise, tandis que le supermarché est le plus cher à 15 000 francs. Sept clients sur dix disent que la boutique de Mama Ngo, à 13 500 francs, offre le meilleur rapport qualité-prix, et le présentateur conseille de comparer les prix, de vérifier la qualité et de toujours demander un reçu. »
\end{exoresolu}

\begin{methode}[Résumer un document sonore en trois étapes]
\begin{enumerate}
\item \textbf{Phrase 1 — le thème} : qui parle et de quoi ? (\eng{The programme compares...})
\item \textbf{Phrase 2 — les faits} : les chiffres et lieux essentiels, reliés par \eng{but}, \eng{while}, \eng{finally}.
\item \textbf{Phrase 3 — la conclusion} : l'avis des clients ou le conseil du journaliste (\eng{...offers the best value for money}).
\end{enumerate}
\end{methode}

\begin{savaistu}[Les radios communautaires au Cameroun anglophone]
Dans les régions du Nord-Ouest et du Sud-Ouest du Cameroun, les radios communautaires jouent un rôle central dans la vie quotidienne. Elles diffusent régulièrement des émissions de conseils aux consommateurs — comparaisons de prix, alertes sur les produits de mauvaise qualité, droits des acheteurs — en anglais, mais aussi en \cle{pidgin} (Cameroon Pidgin English), la langue véhiculaire parlée par des millions de Camerounais. Écouter ces émissions est un excellent entraînement à la compréhension orale : les présentateurs y parlent clairement, répètent les chiffres et utilisent un vocabulaire concret du marché et de la vie quotidienne.
\end{savaistu}

\begin{exercice}[À vous de jouer]
Répondez en phrases complètes, en anglais :
\begin{enumerate}
\item How much does the bag of rice cost at the Central Supermarket?
\item Why do some customers complain about Mile 17 Market?
\item Which shop offers the best value for money, according to the customers?
\item What two pieces of advice does the presenter give at the end?
\end{enumerate}
\end{exercice}

\begin{corrige}[Exercice : compréhension du script]
\begin{enumerate}
\item \eng{It costs 15,000 francs.} (« Il coûte 15 000 francs. »)
\item \eng{They complain because the rice is mixed with stones.} (« Parce que le riz est mélangé à des cailloux. »)
\item \eng{Mama Ngo's store offers the best value for money.} (« La boutique de Mama Ngo offre le meilleur rapport qualité-prix. »)
\item \eng{He advises listeners to compare prices, check the quality, and always ask for a receipt.} (« Il conseille de comparer les prix, de vérifier la qualité et de toujours demander un reçu. »)
\end{enumerate}
\end{corrige}

Vous savez désormais exploiter un script d'écoute en deux passes : d'abord le sens global, puis les détails chiffrés. La section suivante vous entraînera à répondre à des questions de compréhension de types variés (vrai/faux, QCM, réponses courtes) sur ce même document.

\coursec{Compréhension : questions de types variés}

Après avoir écouté et lu le script \eng{The Best Buy} dans la section précédente, nous allons maintenant vérifier notre compréhension à l'aide des trois grands types de questions que l'on rencontre à l'examen : le \cle{vrai/faux justifié}, le \cle{QCM} (questionnaire à choix multiples) et les \cle{questions à réponse courte}. Chaque type exige une méthode précise : c'est ce que nous allons apprendre ici, pas à pas.

\courssub{Rappel du texte support}

Avant de répondre, relisons une version condensée du script d'écoute. Gardez en tête les chiffres, les noms de lieux et les opinions exprimées : ce sont eux qui seront testés.

\begin{textebox}[The Best Buy — script résumé (Buea)]
\eng{Good morning, listeners! Welcome to “Consumer Watch”, your weekly programme on Radio Buea. Today we ask a simple question: where can you find the best buy in town? Our reporter visited three shops. At the City Supermarket, a fifty-kilogram bag of rice costs 15,000 CFA francs. It is a well-known brand, but it is the most expensive. At Mile 17 Market, the same bag costs only 12,000 CFA francs. It is the cheapest, but be careful: the rice is mixed with small stones, so its quality is poor. Finally, at Mama Ngo's Store near the Post Office, the bag costs 13,500 CFA francs. The rice is clean, well-dried and of good quality. We asked ten customers in the street: seven of them prefer Mama Ngo's Store because it offers the best value for money. Our advice to consumers: compare prices, check the quality, and always ask for a receipt after paying. A receipt is your proof of purchase if there is a problem. Remember: the cheapest product is not always the best buy. A best buy means a good price and good quality together. Thank you for listening, and see you next week!}
\end{textebox}

\textbf{Glossaire}\par

\begin{center}
\begin{tabularx}{\linewidth}{|l|l|X|}
\hline
\rowcolor{popblueL} \textbf{Mot anglais} & \textbf{Nature} & \textbf{Traduction française} \\
\hline
best buy & nom & la meilleure affaire, le meilleur achat \\
\hline
brand & nom & la marque \\
\hline
to be mixed with & verbe & être mélangé à \\
\hline
value for money & expression & le rapport qualité-prix \\
\hline
receipt & nom (« ri-ssit », le \emph{p} ne se prononce pas) & le reçu, la quittance \\
\hline
proof of purchase & expression & la preuve d'achat \\
\hline
well-dried & adjectif & bien séché \\
\hline
CFA francs & nom & francs CFA (monnaie locale) \\
\hline
\end{tabularx}
\end{center}

\courssub{Type 1 : le vrai/faux justifié}

\begin{methode}[Comment répondre à un vrai/faux]
\begin{enumerate}
\item Lis d'abord l'affirmation et souligne les mots-clés (chiffres, noms propres, adjectifs).
\item Retrouve dans le script la phrase qui parle du même sujet.
\item Compare : l'affirmation dit-elle exactement la même chose que le texte ?
\item Réponds \eng{True} ou \eng{False}, puis \cle{cite toujours la phrase du texte} qui justifie ta réponse. Une réponse non justifiée perd des points au Bac, même si elle est correcte.
\end{enumerate}
\end{methode}

\begin{exoresolu}[Vrai/Faux — corrigé type]
Énoncé : réponds par \eng{True} ou \eng{False} et justifie par une citation du script.
\begin{enumerate}
\item \eng{The supermarket bag is the cheapest.}
\item \eng{Seven customers prefer Mama Ngo's Store.}
\item \eng{The rice at Mile 17 Market is of good quality.}
\end{enumerate}
\tcblower
Solution détaillée :
\begin{enumerate}
\item \eng{False.} Justification : \eng{“It is the most expensive.”} Le sac du supermarché coûte 15 000 F CFA, c'est le plus cher, pas le moins cher. L'affirmation inverse l'information du texte.
\item \eng{True.} Justification : \eng{“Seven of them prefer Mama Ngo's Store because it offers the best value for money.”} Le chiffre sept correspond exactement au texte.
\item \eng{False.} Justification : \eng{“The rice is mixed with small stones, so its quality is poor.”} Le riz de Mile 17 est de mauvaise qualité (\eng{poor quality}), et non de bonne qualité.
\end{enumerate}
\end{exoresolu}

\begin{attention}[Piège fréquent : l'inversion]
Au vrai/faux, l'examinateur adore inverser une information du texte : \eng{cheapest} à la place de \eng{most expensive}, \eng{good quality} à la place de \eng{poor quality}. Ne réponds jamais de mémoire : vérifie l'adjectif exact dans le script. Autre piège : ne traduis pas mot à mot. \eng{False} ne se justifie pas en disant « c'est faux », mais en citant ce que le texte dit réellement.
\end{attention}

\courssub{Type 2 : le QCM (multiple choice questions)}

Dans un QCM, une seule réponse est correcte ; les autres options, appelées \cle{distracteurs}, sont conçues pour te tromper.

\begin{attention}[Les distracteurs reprennent des chiffres du texte]
Dans les QCM, les distracteurs reprennent souvent des chiffres entendus dans le texte : ici 12,000 F CFA (Mile 17) et 15,000 F CFA (supermarché) apparaissent comme options alors que la bonne réponse est 13,500 F CFA. \cle{Écoute la phrase entière} avant de choisir : le chiffre seul ne suffit pas, il faut savoir à quel magasin il correspond.
\end{attention}

\begin{exoresolu}[QCM — corrigé type]
Énoncé : entoure la bonne réponse.
\begin{enumerate}
\item \eng{How much does the bag cost at Mama Ngo's Store?} a) 12,000 F CFA \quad b) 13,500 F CFA \quad c) 15,000 F CFA
\item \eng{What should consumers always ask for?} a) a discount \quad b) a receipt \quad c) a gift
\end{enumerate}
\tcblower
Solution détaillée :
\begin{enumerate}
\item Réponse \textbf{b) 13,500 F CFA}. Le script dit : \eng{“At Mama Ngo's Store near the Post Office, the bag costs 13,500 CFA francs.”} Les options a) et c) sont des distracteurs : 12,000 F CFA est le prix de Mile 17 Market et 15,000 F CFA celui du City Supermarket.
\item Réponse \textbf{b) a receipt}. Le script dit : \eng{“Always ask for a receipt after paying.”} Ni \eng{discount} (remise) ni \eng{gift} (cadeau) ne sont mentionnés.
\end{enumerate}
\end{exoresolu}

\courssub{Type 3 : les réponses courtes (short answers)}

Pour une question ouverte commençant par \eng{Where}, \eng{Why}, \eng{Which}, \eng{How much}, la réponse doit être \cle{courte mais complète} : une phrase simple qui reprend les mots de la question.

\begin{exemplebox}[Exemples traduits : question et réponse]
\eng{Why is the supermarket rice expensive?} « Pourquoi le riz du supermarché est-il cher ? »

\eng{Because it is a well-known brand.} « Parce que c'est une marque connue. »

\eng{Where does the programme take place?} « Où se déroule l'émission ? »

\eng{It takes place in Buea.} « Elle se déroule à Buéa. »
\end{exemplebox}

\begin{exoresolu}[Réponses courtes — corrigé type]
Énoncé : réponds par une phrase complète.
\begin{enumerate}
\item \eng{Where does the programme take place?}
\item \eng{Why is the cheapest bag not the best buy?}
\item \eng{Which shop offers the best value for money?}
\end{enumerate}
\tcblower
Solution détaillée :
\begin{enumerate}
\item \eng{The programme takes place in Buea.} (On reprend le verbe de la question, \eng{take place}, et on ajoute le lieu.)
\item \eng{Because its quality is poor.} (La question commence par \eng{Why}, donc la réponse commence par \eng{Because} : le riz à 12 000 F CFA est mélangé à des cailloux.)
\item \eng{Mama Ngo's Store offers the best value for money.} (La question \eng{Which shop...?} appelle un nom de magasin en tête de réponse.)
\end{enumerate}
\end{exoresolu}

\begin{propriete}[Correspondance question / réponse]
Le mot interrogatif annonce la nature de la réponse attendue :
\begin{itemize}
\item \eng{Where...?} → un lieu : \eng{In Buea. / At Mile 17 Market.}
\item \eng{Why...?} → une raison introduite par \eng{Because...} (jamais \eng{Because of} suivi d'une proposition conjuguée).
\item \eng{Which...?} → un choix parmi plusieurs éléments : \eng{Mama Ngo's Store.}
\item \eng{How much...?} → un prix : \eng{It costs 13,500 CFA francs.}
\item \eng{How many...?} → une quantité : \eng{Seven customers.}
\end{itemize}
\end{propriete}

\begin{attention}[Erreurs fréquentes des francophones]
\begin{itemize}
\item Ne dis pas \eng{*How much costs the bag?} ; dis \eng{How much does the bag cost?} (auxiliaire \eng{does} obligatoire).
\item Ne traduis pas « combien » toujours par \eng{how much} : devant un nom comptable, utilise \eng{how many} : \eng{How many customers prefer Mama Ngo's Store?}
\item \eng{Because} ne se combine jamais avec \eng{so} dans la même phrase : on dit \eng{Because the quality is poor, it is not the best buy} ou \eng{The quality is poor, so it is not the best buy}, mais pas les deux ensemble.
\end{itemize}
\end{attention}

\courssub{Question ouverte : donner son opinion}

La dernière question invite à une réponse personnelle, en deux phrases environ. Utilise les marqueurs d'opinion : \eng{In my opinion}, \eng{I think that}, \eng{For me}, \eng{I believe that}.

\begin{exemplebox}[Modèle de réponse personnelle]
\eng{In your opinion, what makes a product a “best buy”?}

\eng{In my opinion, a best buy is a product that combines a fair price and good quality. I also think that a good product must last a long time, so I do not waste my money.}

« À mon avis, une bonne affaire est un produit qui combine un prix juste et une bonne qualité. Je pense aussi qu'un bon produit doit durer longtemps, pour ne pas gaspiller mon argent. »
\end{exemplebox}

\courssub{Correction collective : justifier par le script}

En classe, la correction se fait à voix haute : chaque élève donne sa réponse \cle{puis lit la phrase du script} qui la prouve. Cette habitude entraîne à la fois la compréhension écrite et l'expression orale.

\begin{center}
\begin{tabularx}{\linewidth}{|X|X|X|}
\hline
\rowcolor{popblueL} \textbf{Question} & \textbf{Réponse} & \textbf{Justification (citation du script)} \\
\hline
\eng{The supermarket bag is the cheapest.} & \eng{False} & \eng{“It is the most expensive.”} \\
\hline
\eng{Seven customers prefer Mama Ngo's Store.} & \eng{True} & \eng{“Seven of them prefer Mama Ngo's Store.”} \\
\hline
\eng{The rice at Mile 17 Market is of good quality.} & \eng{False} & \eng{“The rice is mixed with small stones.”} \\
\hline
\eng{Price at Mama Ngo's Store?} & \eng{13,500 F CFA} & \eng{“The bag costs 13,500 CFA francs.”} \\
\hline
\eng{What should consumers ask for?} & \eng{A receipt} & \eng{“Always ask for a receipt after paying.”} \\
\hline
\eng{Best value for money?} & \eng{Mama Ngo's Store} & \eng{“It offers the best value for money.”} \\
\hline
\end{tabularx}
\end{center}

\begin{aretenir}[L'idée-clé de la leçon]
Le moins cher n'est pas toujours la meilleure affaire. Un \eng{best buy}, c'est un bon prix \emph{et} une bonne qualité réunis.
\end{aretenir}

\begin{popfigure}
\begin{tikzpicture}[node distance=0.6cm and 1.2cm]
\node[draw=poporange, fill=poporangeL, rounded corners, align=center, font=\small, text width=3.4cm] (cheap) {\textbf{Cheapest}\\ \eng{12,000 F CFA}\\ poor quality};
\node[draw=popmuted, fill=popcream, rounded corners, align=center, font=\small, text width=3.4cm, right=of cheap] (noteq) {\eng{Cheapest} $\neq$\\ \eng{Best buy}};
\node[draw=popgreen, fill=popgreenL, rounded corners, align=center, font=\small, text width=4.2cm, right=of noteq] (best) {\textbf{Best buy}\\ \eng{good price + good quality}\\ (Mama Ngo's : 13,500 F CFA)};
\draw[-{Stealth}, very thick, poporange] (cheap) -- (noteq);
\draw[-{Stealth}, very thick, popgreen] (noteq) -- (best);
\end{tikzpicture}
\legende{Du prix le plus bas au meilleur achat : le prix seul ne suffit pas.}
\end{popfigure}

\begin{exercice}[À toi de jouer]
Réponds aux questions suivantes sur le script \eng{The Best Buy}.
\begin{enumerate}
\item Vrai/Faux (justifie) : \eng{Mama Ngo's Store is near the Market.}
\item QCM : \eng{How many customers were asked?} a) three \quad b) seven \quad c) ten
\item Réponse courte : \eng{Why should consumers ask for a receipt?}
\item Opinion : \eng{Do you always compare prices before buying? Why or why not?} (deux phrases)
\end{enumerate}
\end{exercice}

\begin{corrige}[Exercice « À toi de jouer »]
\begin{enumerate}
\item \eng{False.} Justification : \eng{“At Mama Ngo's Store near the Post Office.”} Le magasin est près de la poste, pas du marché.
\item Réponse \textbf{c) ten} : \eng{“We asked ten customers in the street.”} Sept est le nombre de ceux qui préfèrent Mama Ngo's Store, pas le total interrogé : attention au distracteur.
\item \eng{Because a receipt is your proof of purchase if there is a problem.}
\item Réponse personnelle possible : \eng{Yes, I always compare prices before buying because I want to save money. I also check the quality, because the cheapest product is not always the best buy.}
\end{enumerate}
\end{corrige}

\coursec{Production guidée : conseiller le meilleur achat}

Après avoir écouté le script \eng{The Best Buy} et vérifié votre compréhension grâce aux questions de types variés (vrai/faux, QCM, réponses courtes), il est temps de passer à la \cle{production}. L'objectif de cette dernière étape est clair : réemployer le lexique de la consommation et les structures de comparaison pour \textbf{conseiller le meilleur achat}, d'abord à l'oral, puis à l'écrit. C'est la compétence qui vous servira dans la vie réelle : au marché Mokolo, dans une boutique d'électroménagers à Douala ou dans un supermarché de Bamenda, savoir comparer et décider en anglais est un atout précieux.

\courssub{Réemploi du lexique : les cinq mots-clés}

\begin{definition}[Le lexique du meilleur achat]
Voici les cinq mots essentiels de cette leçon. Retenez leur nature, leur prononciation et leur emploi exact :
\begin{itemize}
\item \eng{bargain} (nom) : une bonne affaire, un article acheté à très bon prix. \eng{At 5,000 francs, this bag is a real bargain!} = « À 5 000 francs, ce sac est une vraie bonne affaire ! »
\item \eng{quality} (nom) : la qualité. \eng{Quality is more important than price.} = « La qualité est plus importante que le prix. »
\item \eng{discount} (nom) : une remise, une réduction. \eng{The shopkeeper gave me a 10\% discount.} = « Le vendeur m'a accordé une remise de 10 \%. »
\item \eng{receipt} (nom) : un reçu (attention, le \emph{p} ne se prononce pas : « ri-site »). \eng{Keep your receipt in case of a problem.} = « Gardez votre reçu en cas de problème. »
\item \eng{value for money} (expression) : le rapport qualité-prix. \eng{This phone offers excellent value for money.} = « Ce téléphone offre un excellent rapport qualité-prix. »
\end{itemize}
\end{definition}

\begin{attention}[Faux amis et pièges lexicaux]
\begin{itemize}
\item \eng{receipt} ne veut pas dire « recette » de cuisine (qui se dit \eng{recipe}) : c'est un \textbf{reçu}, un justificatif d'achat.
\item \eng{bargain} (nom) ne veut pas dire « barguigner » : c'est une \textbf{bonne affaire}. Le verbe « marchander » se dit \eng{to bargain} ou \eng{to haggle}.
\item On dit \eng{good value for money}, jamais \eng{a good quality-price report} (calque du français).
\end{itemize}
\end{attention}

\begin{exoresolu}[Exercice de réemploi : phrases à trous]
Énoncé — Complétez les phrases avec : \eng{bargain}, \eng{quality}, \eng{discount}, \eng{receipt}, \eng{value for money}.
\begin{enumerate}
\item This market is famous for its low prices; every purchase is a real \ldots
\item Don't only look at the price; check the \ldots\ of the fabric.
\item Because I bought two radios, the seller gave me a \ldots\ of 2,000 francs.
\item Always ask for a \ldots\ when you pay, so you can exchange the product later.
\item This bicycle is expensive, but it will last ten years: it is excellent \ldots
\end{enumerate}
\tcblower
Solution détaillée :
\begin{enumerate}
\item \textbf{bargain} — « chaque achat est une vraie bonne affaire » (nom comptable précédé de \eng{a real}).
\item \textbf{quality} — « vérifiez la qualité du tissu » (on oppose prix et qualité).
\item \textbf{discount} — « une remise de 2 000 francs » (précédé de l'article \eng{a}).
\item \textbf{receipt} — « demandez toujours un reçu » (justificatif nécessaire pour un échange).
\item \textbf{value for money} — « un excellent rapport qualité-prix » (expression invariable, sans article).
\end{enumerate}
\end{exoresolu}

\courssub{Production orale guidée : le jeu de rôle vendeur / client}

\begin{experience}[Role-play : at the electronics shop]
Travaillez en binômes. L'élève A est le \textbf{vendeur} (\eng{shopkeeper}) d'une boutique d'électronique à Yaoundé ; l'élève B est le \textbf{client} (\eng{customer}) qui hésite entre deux radios ou deux téléphones. Le client doit comparer les deux produits selon trois critères : le \textbf{prix} (\eng{price}), la \textbf{qualité} (\eng{quality}) et la \textbf{garantie} (\eng{guarantee}). Le vendeur doit défendre ses produits et proposer éventuellement une remise. À la fin, le client décide, paie et demande un reçu. Puis échangez les rôles.
\end{experience}

La réussite du dialogue repose sur une progression logique en cinq étapes, qu'il faut mémoriser :

\begin{popfigure}
\begin{center}
\begin{tikzpicture}[
node distance=0.55cm,
etape/.style={rectangle, rounded corners=4pt, draw=popblue, fill=popblueL, thick, align=center, minimum width=3.1cm, minimum height=0.95cm, font=\small},
fleche/.style={-{Stealth[length=2.5mm]}, thick, popdark}
]
\node[etape, align=center] (n1){\textbf{1. \eng{Greet}}\\ \eng{Good morning!}};
\node[etape, right=of n1, align=center] (n2){\textbf{2. \eng{Ask prices}}\\ \eng{How much is this?}};
\node[etape, right=of n2, align=center] (n3){\textbf{3. \eng{Compare}}\\ \eng{cheaper / better than}};
\node[etape, right=of n3, align=center] (n4){\textbf{4. \eng{Decide}}\\ \eng{I'll take it.}};
\node[etape, right=of n4, align=center] (n5){\textbf{5. \eng{Pay \& receipt}}\\ \eng{Can I have a receipt?}};
\draw[fleche] (n1) -- (n2);
\draw[fleche] (n2) -- (n3);
\draw[fleche] (n3) -- (n4);
\draw[fleche] (n4) -- (n5);
\end{tikzpicture}
\end{center}
\legende{Organigramme du dialogue d'achat : saluer, demander les prix, comparer, décider, payer et demander le reçu.}
\end{popfigure}

\begin{textebox}[Model dialogue : “The best buy”]
\textbf{Customer:} Good morning! How much is this radio?\\
\textbf{Shopkeeper:} Good morning! That one is 25,000 francs.\\
\textbf{Customer:} And that one over there?\\
\textbf{Shopkeeper:} That one is 20,000 francs, but the first one is more durable and has a one-year guarantee.\\
\textbf{Customer:} I see. The second one is cheaper, but the first one seems to be of better quality. Can you give me a discount?\\
\textbf{Shopkeeper:} For you, I can reduce it to 23,000 francs. It's a real bargain!\\
\textbf{Customer:} Then I'll take the first one; it's the best buy. Can I have a receipt, please?\\
\textbf{Shopkeeper:} Of course. Here you are. Thank you and come again!

\emph{Glossaire :} \eng{durable} (adj.) = durable, solide ; \eng{one-year guarantee} = garantie d'un an ; \eng{to reduce} (v.) = réduire, baisser ; \eng{here you are} = voilà, je vous le donne.
\end{textebox}

\begin{exemplebox}[Phrases utiles traduites]
\begin{itemize}
\item \eng{This one is cheaper, but that one lasts longer.} = « Celui-ci est moins cher, mais celui-là dure plus longtemps. »
\item \eng{I'll take it.} = « Je le prends. » (formule de décision, très fréquente à l'oral)
\item \eng{How much is this radio?} = « Combien coûte cette radio ? »
\item \eng{It offers better value for money.} = « Il offre un meilleur rapport qualité-prix. »
\item \eng{Can I have a receipt, please?} = « Puis-je avoir un reçu, s'il vous plaît ? »
\end{itemize}
\end{exemplebox}

\begin{attention}[La question du prix : l'erreur n° 1 des francophones]
Ne dites jamais \eng{How much costs this?} — c'est un calque du français « combien coûte ceci ? ». En anglais, le verbe \eng{to cost} exige l'auxiliaire \eng{do} à la forme interrogative. Les deux formes correctes sont :
\begin{itemize}
\item \eng{How much does this cost?} (avec l'auxiliaire \eng{does} + base verbale \eng{cost}, sans \emph{s})
\item \eng{How much is this?} (plus simple, avec le verbe \eng{to be})
\end{itemize}
De même, ne dites pas \eng{It is more cheap} mais \eng{It is cheaper} : les adjectifs courts prennent \emph{-er} au comparatif.
\end{attention}

\begin{center}
\begin{tabularx}{\linewidth}{|X|X|X|}
\hline
\rowcolor{popblueL} Français & English & Remarque \\
\hline
Combien coûte ceci ? & How much is this? / How much does this cost? & jamais \eng{How much costs this?} \\
\hline
Celui-ci est moins cher que celui-là. & This one is cheaper than that one. & comparatif en \emph{-er} + \eng{than} \\
\hline
Celui-là est de meilleure qualité. & That one is of better quality. & \eng{good} devient \eng{better} (irrégulier) \\
\hline
Il dure plus longtemps. & It lasts longer. & \eng{long} $\rightarrow$ \eng{longer} \\
\hline
Je le prends. & I'll take it. & contraction de \eng{I will} \\
\hline
C'est une bonne affaire. & It's a bargain. & pas de lien avec « barguigner » \\
\hline
Puis-je avoir un reçu ? & Can I have a receipt? & \emph{p} muet : « ri-site » \\
\hline
\end{tabularx}
\end{center}

\courssub{Production écrite courte : “My best buy”}

\begin{methode}[Rédiger le paragraphe “My best buy” en quatre étapes]
Pour réussir votre paragraphe de 6 à 8 lignes, suivez ce plan strict :
\begin{enumerate}
\item \textbf{Phrase d'introduction} : présentez le produit acheté (quoi, où, quand). \eng{Last month, I bought a second-hand bicycle at the Buea market.} = « Le mois dernier, j'ai acheté un vélo d'occasion au marché de Buea. »
\item \textbf{2 à 3 phrases de comparaison} : opposez le prix et la qualité avec au moins deux comparatifs (\eng{cheaper than}, \eng{better than}, \eng{more durable than}, \eng{longer}). \eng{It was cheaper than the new one in the shop, and its frame was stronger.} = « Il était moins cher que le neuf du magasin, et son cadre était plus solide. »
\item \textbf{Phrase de conclusion} : expliquez pourquoi c'est votre meilleur achat. \eng{That is why it was the best buy of the year: real value for money.} = « C'est pourquoi c'était le meilleur achat de l'année : un vrai rapport qualité-prix. »
\item \textbf{Vérification finale} : comptez vos 5 mots du lexique (\eng{bargain}, \eng{quality}, \eng{discount}, \eng{receipt}, \eng{value for money}) et vos 2 comparatifs ; relisez les prix et les chiffres.
\end{enumerate}
\end{methode}

\begin{exemplebox}[Modèle de paragraphe rédigé]
\eng{Last Saturday, I bought a mobile phone at a shop in Douala. There were two models: the first one was 45,000 francs and the second one was cheaper, only 30,000 francs. However, the first phone was more durable and had a better camera, so its quality was higher. The shopkeeper gave me a small discount and a receipt with a one-year guarantee. In the end, I chose the first phone because it offered real value for money. It was not a bargain at first sight, but I am sure it was my best buy this year.}

\emph{Traduction :} Samedi dernier, j'ai acheté un téléphone portable dans une boutique de Douala. Il y avait deux modèles : le premier coûtait 45 000 francs et le second était moins cher, seulement 30 000 francs. Cependant, le premier téléphone était plus solide et avait un meilleur appareil photo, donc sa qualité était supérieure. Le vendeur m'a fait une petite remise et m'a donné un reçu avec une garantie d'un an. Finalement, j'ai choisi le premier téléphone parce qu'il offrait un vrai rapport qualité-prix. Ce n'était pas une bonne affaire à première vue, mais je suis sûr que c'était mon meilleur achat de l'année.
\end{exemplebox}

\begin{exercice}[À vous d'écrire]
Rédigez votre propre paragraphe \eng{My best buy} (6 à 8 lignes) sur un achat réel ou imaginaire (radio, chaussures, sac, vivres au marché). Contraintes obligatoires : au moins \textbf{5 mots du lexique} de la leçon et au moins \textbf{2 comparatifs}. Soulignez-les dans votre copie.
\end{exercice}

\courssub{Grille d'auto-évaluation}

Avant de rendre votre production (orale ou écrite), évaluez-vous honnêtement à l'aide de cette grille. Chaque critère vaut un point ; visez le score de 4 sur 4.

\begin{center}
\begin{tabularx}{\linewidth}{|X|l|l|}
\hline
\rowcolor{popgreenL} Critère & Oui & Non \\
\hline
J'ai utilisé correctement au moins 5 mots du lexique (\eng{bargain}, \eng{quality}, \eng{discount}, \eng{receipt}, \eng{value for money}). & & \\
\hline
J'ai formé au moins 2 comparatifs corrects (\eng{cheaper than}, \eng{better than}, \eng{more durable than}). & & \\
\hline
Les prix et les chiffres cités sont exacts et prononcés correctement (\eng{25,000 francs} = « twenty-five thousand francs »). & & \\
\hline
Ma conclusion est claire : j'explique pourquoi c'est le meilleur achat. & & \\
\hline
\end{tabularx}
\end{center}

\begin{aretenir}[L'essentiel de la section]
Pour conseiller le meilleur achat en anglais : suivez la trame \textbf{\eng{Greet} $\rightarrow$ \eng{Ask prices} $\rightarrow$ \eng{Compare} $\rightarrow$ \eng{Decide} $\rightarrow$ \eng{Pay \& receipt}} ; interrogez sur le prix avec \eng{How much is this?} ou \eng{How much does this cost?} ; comparez avec les comparatifs (\eng{cheaper than}, \eng{better than}) ; concluez avec \eng{I'll take it} et terminez toujours par \eng{Can I have a receipt?}. À l'écrit, structurez votre paragraphe : introduction, comparaison, conclusion justifiée.
\end{aretenir}

\newpage

\coursec{Activité d'intégration}

\coursec{Lesson 13 — Listening: Texts on comparing prices or quality to determine the best buys for consumer economy and community resources/services}

\courssub{Warm-up et découverte du thème}

\begin{experience}[Discussion d'entrée : nos habitudes d'achat]
L'enseignant affiche ou dicte ces questions. En binômes, les élèves discutent 3 minutes en anglais, puis quelques binômes rapportent à la classe.
\begin{enumerate}
\item \eng{Where do you usually buy household appliances in your area?} (Où achetez-vous habituellement les appareils ménagers dans votre quartier ?)
\item \eng{What do you check before buying: price, brand, guarantee, durability?} (Qu'est-ce que vous vérifiez avant d'acheter : le prix, la marque, la garantie, la durabilité ?)
\item \eng{Have you ever regretted a purchase because you didn't compare?} (Avez-vous déjà regretté un achat parce que vous n'aviez pas comparé ?)
\end{enumerate}
\emph{Objectif :} activer le schéma notionnel \eng{shopping / consumer choices} et le lexique de base (\eng{price, quality, guarantee, brand, market, shop}).
\end{experience}

\courssub{Lexique thématique : consommation et comparaison}

\begin{aretenir}[Mots et expressions clés du thème — \eng{Consumer Economy}]
\begin{center}
\begin{tabularx}{\linewidth}{|X|X|X|}
\hline
\rowcolor{popblueL} \textbf{English} & \textbf{Prononciation (approx.)} & \textbf{Français} \\
\hline
\eng{price} & « praïs » & prix \\
\eng{cheap / cheaper (than)} & « tchiip / tchiip-eu » & bon marché / moins cher (que) \\
\eng{expensive / more expensive (than)} & « ik-spen-siv / mor ik-spen-siv » & cher / plus cher (que) \\
\eng{quality} & « quo-li-ti » & qualité \\
\eng{durable / more durable} & « deu-rea-beul / mor deu-rea-beul » & durable / plus résistant \\
\eng{sturdy / sturdier} & « steur-di / steur-di-eu » & solide, robuste / plus solide \\
\eng{guarantee / warranty} & « ga-ran-ti / wo-ran-ti » & garantie \\
\eng{value for money} & « va-liou for ma-ni » & rapport qualité-prix \\
\eng{a bargain} & « a bar-gaine » & une bonne affaire \\
\eng{to afford} & « a forde » & avoir les moyens d'acheter \\
\eng{customer / buyer} & « cos-to-meu / bai-eu » & client / acheteur \\
\eng{seller / vendor} & « se-leu / ven-deu » & vendeur \\
\eng{best buy} & « beste bai » & le meilleur achat \\
\eng{to compare} & « tou com-pare » & comparer \\
\eng{pros and cons} & « proz and konz » & avantages et inconvénients \\
\hline
\end{tabularx}
\end{center}
\emph{Astuce :} \eng{warranty} s'emploie surtout pour les produits manufacturés (électronique, électroménager) ; \eng{guarantee} est plus général.
\end{aretenir}

\begin{attention}[Faux amis et pièges fréquents]
\begin{itemize}
\item \eng{Actually} = « en fait, en réalité » (pas « actuellement » $\rightarrow$ \eng{currently} / \eng{at the moment}).
\item \eng{Price} n'est pas un adjectif : on ne dit pas \eng{*The price is expensive}. On dit \eng{The price is high / low} ou \eng{The product is expensive / cheap}.
\item \eng{More cheaper} $\times$ $\rightarrow$ \eng{cheaper} (court) ou \eng{more expensive} (long).
\item \eng{Than} (comparaison) $\neq$ \eng{Then} (temps / conséquence). \eng{Cheaper \textbf{than}} (moins cher \textbf{que}).
\item \eng{Durable} (adjectif) $\neq$ \eng{Duration} (nom, durée). \eng{Durability} = durabilité.
\end{itemize}
\end{attention}

\courssub{Stratégies et méthode d'écoute}

\begin{methode}[Réussir une compréhension orale sur la comparaison de produits]
\begin{enumerate}
\item \textbf{Avant l'écoute (Prédiction)} : Lisez le titre, observez les images ou les mots-clés au tableau. Posez-vous : \eng{What products? What criteria (price, quality, guarantee)?} Préparez une grille de prise de notes (colonnes : Product, Price, Quality, Guarantee, Verdict).
\item \textbf{Première écoute (Globalité)} : Identifiez le \textbf{thème}, le \textbf{contexte} (radio, magasin, dialogue amis), les \textbf{intervenants} et le \textbf{résultat global} (quel produit gagne ?).
\item \textbf{Deuxième écoute (Détails chiffrés)} : Relevez \textbf{exactement} les nombres (prix, durée de garantie, puissance, poids), les noms propres (marques, vendeurs, lieux) et les adjectifs de qualité (\eng{sturdy, flimsy, reliable, faulty}).
\item \textbf{Troisième écoute (Vérification)} : Complétez les trous, corrigez les approximations, validez le verdict final.
\item \textbf{Après l'écoute (Réemploi)} : Reformulez oralement la comparaison en utilisant \eng{cheaper than, more durable than, the best value for money}.
\end{enumerate}
\end{methode}

\courssub{Texte support : script d'écoute « The Best Buy — Smartphone Edition »}

\begin{textebox}[Script audio — Teacher's reading text / Enregistrement OnBuch+]
\eng{Good morning, listeners! Welcome to “The Best Buy”, your weekly consumer guide on Buea Community Radio. I'm your host, Peter. Today, we help students prepare for the new academic year. We compared two popular smartphones sold at Molyko Market: the “Nova X” and the “TechPro 5”.}

\eng{First, the “Nova X”, sold by Madam Ngwa at Shop 12. It costs 125,000 FCFA. It has a 6.5-inch screen, 128 gigabytes of storage, and a 5000 mAh battery that lasts two days. The build quality is solid: aluminium frame, toughened glass. Madam Ngwa offers a two-year warranty and a free screen protector.}

\eng{Second, the “TechPro 5”, sold by Mr. Bello at the phone plaza. It is cheaper: 95,000 FCFA. It also has a 6.5-inch screen and 128 gigabytes. But the battery is smaller, 4000 mAh, lasting one day. The body is plastic, less sturdy. The warranty is only six months. No free accessories.}

\eng{So, listeners, the “Nova X” is more expensive, but it is more durable, has a longer battery life and a much longer warranty. The “TechPro 5” is a bargain if your budget is tight, but you risk changing the phone sooner. For us, the “Nova X” is the best buy for a student who needs reliability. Vote on our Facebook page!}
\end{textebox}

\begin{definition}[Glossaire du script]
\begin{itemize}
\item \eng{storage} (n.) : espace de stockage (mémoire)
\item \eng{battery life} (n.) : autonomie de la batterie
\item \eng{aluminium frame} (n.) : châssis en aluminium
\item \eng{toughened glass} (n.) : verre trempé
\item \eng{screen protector} (n.) : protection d'écran (film ou verre)
\item \eng{plastic} (n./adj.) : plastique
\item \eng{tight (budget)} (adj.) : serré (budget)
\item \eng{reliability} (n.) : fiabilité
\end{itemize}
\end{definition}

\courssub{Compréhension : questions de types variés}

\begin{exercice}[Questions sur le script « The Best Buy »]
\textbf{A. True or False? Justify with a quote from the text.}
\begin{enumerate}
\item The “Nova X” is sold by Mr. Bello.
\item Both phones have the same screen size and storage capacity.
\item The “TechPro 5” has a better battery life than the “Nova X”.
\item The “Nova X” comes with a two-year warranty.
\item The radio host recommends the “TechPro 5” for students.
\end{enumerate}

\textbf{B. Multiple Choice. Tick the correct answer.}
\begin{enumerate}
\item What is the price difference between the two phones?
      \begin{itemize}
      \item[a)] 20,000 FCFA
      \item[b)] 30,000 FCFA
      \item[c)] 35,000 FCFA
      \end{itemize}
\item Which material is used for the “Nova X” body?
      \begin{itemize}
      \item[a)] Plastic
      \item[b)] Aluminium
      \item[c)] Wood
      \end{itemize}
\item How long is the “TechPro 5” warranty?
      \begin{itemize}
      \item[a)] One year
      \item[b)] Six months
      \item[c)] Two years
      \end{itemize}
\end{enumerate}

\textbf{C. Short Answers. Answer in full sentences.}
\begin{enumerate}
\item Which shop sells the “Nova X” and where is it located?
\item Why is the “TechPro 5” considered a bargain?
\item According to the host, why is the “Nova X” the best buy for a student?
\end{enumerate}
\end{exercice}

\begin{corrige}[Exercice — Compréhension « The Best Buy »]
\textbf{A. True or False?}
\begin{enumerate}
\item \textbf{False.} \eng{“First, the 'Nova X', sold by Madam Ngwa at Shop 12.”}
\item \textbf{True.} \eng{“It also has a 6.5-inch screen and 128 gigabytes.”} (said about TechPro 5, matching Nova X specs).
\item \textbf{False.} \eng{“The battery is smaller, 4000 mAh, lasting one day.”} (Nova X lasts two days).
\item \textbf{True.} \eng{“Madam Ngwa offers a two-year warranty…”}
\item \textbf{False.} \eng{“For us, the 'Nova X' is the best buy for a student who needs reliability.”}
\end{enumerate}

\textbf{B. Multiple Choice.}
\begin{enumerate}
\item \textbf{b)} 125,000 - 95,000 = 30,000 FCFA.
\item \textbf{b)} \eng{“aluminium frame”}.
\item \textbf{b)} \eng{“The warranty is only six months.”}
\end{enumerate}

\textbf{C. Short Answers.}
\begin{enumerate}
\item The “Nova X” is sold by Madam Ngwa at Shop 12 in Molyko Market.
\item It is considered a bargain because it costs only 95,000 FCFA, which is 30,000 FCFA cheaper than the Nova X, for similar screen and storage specs.
\item Because it is more durable (aluminium frame), has a longer battery life (two days vs one day) and a much longer warranty (two years vs six months), offering better value for money.
\end{enumerate}
\end{corrige}

\courssub{Point de grammaire : Comparatifs et superlatifs (Rappel et approfondissement)}

\begin{propriete}[Formation et emploi des comparatifs et superlatifs]
\begin{center}
\begin{tabularx}{\linewidth}{|X|X|X|X|}
\hline
\rowcolor{popblueL} \textbf{Type} & \textbf{Adjectif} & \textbf{Comparatif de supériorité} & \textbf{Superlatif} \\
\hline
Court (1 syllabe) & \eng{cheap, strong, fast, long} & \eng{cheaper than, stronger than} & \eng{the cheapest, the strongest} \\
Court terminé par -e & \eng{large, safe} & \eng{larger than, safer than} & \eng{the largest, the safest} \\
Court CVC (double consonne) & \eng{big, hot, thin} & \eng{bigger than, hotter than, thinner than} & \eng{the biggest, the hottest, the thinnest} \\
Long ($\ge$ 2 syllabes) & \eng{expensive, durable, reliable, beautiful} & \eng{more expensive than, more durable than} & \eng{the most expensive, the most durable} \\
Terminé par -y & \eng{heavy, easy, busy} & \eng{heavier than, easier than, busier than} & \eng{the heaviest, the easiest, the busiest} \\
\hline
\end{tabularx}
\end{center}

\textbf{Irréguliers importants :}
\begin{center}
\begin{tabular}{|l|l|l|}
\hline
\rowcolor{popblueL} \textbf{Adjectif / Adverbe} & \textbf{Comparatif} & \textbf{Superlatif} \\
\hline
\eng{good / well} & \eng{better} & \eng{the best} \\
\eng{bad / badly} & \eng{worse} & \eng{the worst} \\
\eng{far} & \eng{farther / further} & \eng{the farthest / the furthest} \\
\eng{much / many} & \eng{more} & \eng{the most} \\
\eng{little} & \eng{less} & \eng{the least} \\
\hline
\end{tabular}
\end{center}

\textbf{Comparatif d'égalité :} \eng{as + adjectif + as} (aussi ... que). Ex. \eng{The TechPro 5 is as fast as the Nova X.}
\textbf{Comparatif d'infériorité :} \eng{less + adjectif + than} (moins ... que) ou \eng{not as ... as}. Ex. \eng{The TechPro 5 is less durable than the Nova X.} / \eng{The TechPro 5 is not as durable as the Nova X.}
\end{propriete}

\begin{exemplebox}[Exemples contextualisés — Comparaison de produits]
\begin{itemize}
\item \eng{The Nova X is \textbf{more expensive than} the TechPro 5.} (Le Nova X est plus cher que le TechPro 5.)
\item \eng{The TechPro 5 is \textbf{cheaper than} the Nova X.} (Le TechPro 5 est moins cher que le Nova X.)
\item \eng{The Nova X has a \textbf{longer} battery life.} (Le Nova X a une autonomie plus longue.)
\item \eng{The Nova X offers \textbf{the best value for money}.} (Le Nova X offre le meilleur rapport qualité-prix.)
\item \eng{The warranty of the Nova X is \textbf{twice as long as} the TechPro 5's.} (La garantie du Nova X est deux fois plus longue que celle du TechPro 5.)
\end{itemize}
\end{exemplebox}

\courssub{Activité d'intégration : Consumer Watch — Enregistrement en studio}

\courssub{Objectifs de l'activité}

À l'issue de cette activité d'intégration, l'élève sera capable de :
\begin{itemize}
\item \textbf{écouter} un mini-reportage oral sur la comparaison de produits et d'en relever les informations essentielles (nom du produit, prix, qualité, garantie) ;
\item \textbf{réemployer} le lexique de la consommation (\eng{price, quality, guarantee, durable, value for money, bargain}) et les structures de comparaison (\eng{cheaper than, more expensive than, the best buy, as ... as}) dans une production orale et écrite ;
\item \textbf{justifier} un choix de consommateur en anglais, à l'oral comme à l'écrit ;
\item \textbf{rédiger} un court compte rendu structuré de six lignes environ ;
\item travailler en \textbf{coopération} dans un cadre de communication authentique : une émission de radio communautaire.
\end{itemize}

\courssub{Situation}

La radio communautaire de ta ville, \eng{Buea Community Radio}, diffuse chaque samedi matin une émission très écoutée : \emph{Consumer Watch}. Le principe est simple : avant les grandes dépenses des familles (rentrée scolaire, fêtes de fin d'année, équipement de la maison), un reporter se rend au marché et compare deux produits similaires vendus par deux commerçants différents. À la fin de l'émission, les auditeurs votent pour le \eng{best buy}, le meilleur achat, en expliquant leur choix.

Cette semaine, l'émission porte sur les \textbf{réchauds à gaz de table} (\eng{table gas cookers}), très demandés avant la rentrée. Ta classe est invitée à enregistrer l'émission en studio. L'enseignant distribue à chaque groupe deux fiches produits. Voici un exemple de fiche utilisée par le groupe témoin :

\begin{textebox}[Product cards — Consumer Watch, Buea Community Radio]
\textbf{Product card A — “Flame King” table gas cooker}\\
Seller: Madam Eposi, Molyko Market, Buea\\
Price: 18,500 FCFA\\
Quality: strong stainless steel body; two burners; lights easily\\
Guarantee: one year; free repair during the guarantee period\\[0.4em]
\textbf{Product card B — “Quick Fire” table gas cooker}\\
Seller: Mr. Tabi, Buea Town Market\\
Price: 14,000 FCFA\\
Quality: lighter metal body; two burners; sometimes difficult to light\\
Guarantee: three months only
\end{textebox}

\begin{definition}[Glossaire utile — Fiches produits]
\begin{itemize}
\item \eng{stainless steel} (n.) « steïn-less steel » : acier inoxydable
\item \eng{burner} (n.) « beur-neu » : brûleur (d'un réchaud)
\item \eng{to light} (v.) « laït » : allumer (intransitif : \eng{it lights easily} = il s'allume facilement)
\item \eng{guarantee} (n.) « ga-ran-ti » : garantie
\item \eng{best buy} (n.) « beste bai » : le meilleur achat, la meilleure affaire
\item \eng{sturdy} (adj.) « steur-di » : solide, robuste (synonyme de \eng{strong} pour un objet)
\end{itemize}
\end{definition}

\courssub{Consignes}

Travail par \textbf{groupes de 4 élèves}. Durée totale : environ 50 minutes.

\textbf{Étape 1 — Préparation du reportage (10 min).} Dans chaque groupe, répartissez les rôles : un \eng{presenter} (présentateur), un \eng{reporter} et deux \eng{sellers} (vendeurs). Lisez vos deux fiches produits distribuées par l'enseignant. Les vendeurs préparent trois ou quatre phrases pour « vendre » leur produit ; le reporter prépare ses questions (\eng{How much is it? What is it made of? How long is the guarantee?}).

\textbf{Étape 2 — Enregistrement de l'émission (15 min).} Chaque groupe passe devant la classe et joue son mini-reportage de deux minutes : le présentateur introduit l'émission, le reporter interviewe les deux vendeurs, puis le présentateur conclut.

\textbf{Étape 3 — Écoute et prise de notes.} Pendant que les autres groupes passent, chaque élève remplit sa \textbf{grille d'écoute} :

\begin{center}
\begin{tabularx}{\linewidth}{|X|X|X|X|}
\hline
\rowcolor{popblueL} Product & Price & Quality / Guarantee & Verdict \\
\hline
Flame King & 18,500 FCFA & strong stainless steel; 1-year guarantee & \ldots \\
\hline
Quick Fire & 14,000 FCFA & lighter metal; 3-month guarantee & \ldots \\
\hline
\end{tabularx}
\end{center}

\textbf{Étape 4 — Le vote du « best buy » (10 min).} Après chaque reportage, la classe vote à main levée. Deux ou trois élèves justifient leur vote à l'oral avec la structure imposée : \eng{We think X is the best buy because it is cheaper and more durable.}

\textbf{Étape 5 — Compte rendu écrit (15 min).} Chaque groupe rédige un court compte rendu de \textbf{six lignes environ} présentant les deux produits comparés, les chiffres relevés et le verdict final justifié. Ce texte sera ramassé et évalué.

\begin{methode}[Réussir son mini-reportage en 5 étapes]
\begin{enumerate}
\item \textbf{Introduire} : nommer l'émission, le thème et les produits comparés. \eng{“Welcome to Consumer Watch. Today we compare two table gas cookers…”}
\item \textbf{Interviewer} : poser des questions courtes et claires sur le prix, la qualité et la garantie. \eng{“Madam Eposi, how much is the Flame King? What is the guarantee?”}
\item \textbf{Comparer} : utiliser au moins deux comparatifs (\eng{cheaper, stronger, longer, better}) et un superlatif (\eng{the best buy}).
\item \textbf{Conclure} : annoncer le vote et inviter les auditeurs à réagir. \eng{“The class voted for the Flame King. Listeners, call us and tell us your choice!”}
\item \textbf{Rédiger} : reprendre les faits dans l'ordre (produit A, produit B, verdict) avec les chiffres exacts.
\end{enumerate}
\end{methode}

\courssub{Ressources à mobiliser}

\begin{aretenir}[Lexique de la comparaison des achats — Récapitulatif]
\eng{price} (prix), \eng{cheap / cheaper (than)} (bon marché / moins cher que), \eng{expensive / more expensive (than)} (cher / plus cher que), \eng{quality} (qualité), \eng{durable / more durable} (durable / plus résistant), \eng{sturdy / sturdier} (solide / plus solide), \eng{guarantee / warranty} (garantie), \eng{value for money} (rapport qualité-prix), \eng{a bargain} (une bonne affaire), \eng{to afford} (avoir les moyens d'acheter), \eng{customer} (client), \eng{seller / vendor} (vendeur), \eng{best buy} (meilleur achat).
\end{aretenir}

\begin{propriete}[Rappel : les comparatifs et le superlatif — Points clés pour l'activité]
\begin{itemize}
\item Adjectifs courts (1 syllabe) : \eng{cheap} $\rightarrow$ \eng{cheaper than} $\rightarrow$ \eng{the cheapest}. \eng{strong} $\rightarrow$ \eng{stronger than}. Exemple : \eng{Quick Fire is cheaper than Flame King.}
\item Adjectifs longs ($\ge$ 2 syllabes) : \eng{durable} $\rightarrow$ \eng{more durable than} $\rightarrow$ \eng{the most durable}. \eng{expensive} $\rightarrow$ \eng{more expensive than}.
\item Irréguliers : \eng{good} $\rightarrow$ \eng{better} $\rightarrow$ \eng{the best} ; \eng{bad} $\rightarrow$ \eng{worse} $\rightarrow$ \eng{the worst}.
\item Égalité : \eng{as strong as} (aussi solide que).
\item Infériorité : \eng{less durable than} / \eng{not as durable as}.
\end{itemize}
\end{propriete}

\begin{attention}[Erreurs fréquentes des francophones — Vigilance pendant l'activité]
\begin{itemize}
\item \textbf{More cheaper} $\times$ $\rightarrow$ \eng{cheaper} (court) ou \eng{more expensive} (long). Jamais les deux.
\item \textbf{The price is expensive} $\times$ $\rightarrow$ Un prix est \eng{high} ou \eng{low}. Un produit est \eng{expensive} ou \eng{cheap}. Dis : \eng{The price is high.} / \eng{The cooker is expensive.}
\item \textbf{Than} vs \textbf{Then} : \eng{cheaper \textbf{than}} (que), pas \eng{cheaper \textbf{then}}.
\item \textbf{Actually} = « en fait » (pas « actuellement » $\rightarrow$ \eng{currently}).
\item \textbf{Guarantee} (nom) vs \eng{Guarantee} (verbe) : \eng{The seller guarantees the product for one year.} (Le vendeur garantit le produit un an.)
\item \textbf{Agreement} : \eng{The burners work well} (pluriel), pas \eng{*works}.
\end{itemize}
\end{attention}

\courssub{Proposition de production et corrigé commenté}

\begin{exoresolu}[Modèle de compte rendu écrit — Consumer Watch]
Rédige le compte rendu de six lignes du reportage sur les deux réchauds à gaz, en reprenant les chiffres exacts et en justifiant le verdict.
\tcblower
\textbf{Production modèle :}

\eng{On Consumer Watch this week, we compared two table gas cookers sold in Buea. Madam Eposi sells the “Flame King” for 18,500 FCFA. It is made of strong stainless steel and it has a one-year guarantee with free repairs. Mr. Tabi sells the “Quick Fire” for 14,000 FCFA, so it is cheaper, but its metal body is lighter and the guarantee is only three months. We think the “Flame King” is the best buy because it is more durable and has a longer guarantee, even though it is more expensive.}

\medskip
\textbf{Commentaire des choix de langue :}
\begin{itemize}
\item \textbf{Ordre logique} : le texte suit le plan produit A $\rightarrow$ produit B $\rightarrow$ verdict, ce qui facilite la lecture et montre la maîtrise de la prise de notes.
\item \textbf{Chiffres exacts} : \eng{18,500 FCFA}, \eng{14,000 FCFA}, \eng{one-year guarantee}, \eng{three months} — l'évaluation porte sur la justesse de ces relevés d'écoute.
\item \textbf{Comparatifs corrects} : \eng{cheaper} (adjectif court, sans \emph{more}), \eng{more durable} (adjectif long), \eng{longer} ; le superlatif \eng{the best buy} reprend le lexique de la leçon.
\item \textbf{Justification} : la structure \eng{We think ... is the best buy because ...} est exactement celle exigée à l'oral ; la concession \eng{even though it is more expensive} (« même s'il est plus cher ») enrichit l'argumentation au niveau Terminale.
\item \textbf{Connecteurs} : \eng{so} (donc), \eng{but} (mais), \eng{because} (parce que) articulent clairement la comparaison.
\item \textbf{Précision lexicale} : \eng{stainless steel}, \eng{free repairs}, \eng{lighter metal body} (corps en métal plus léger/fin) montrent un vocabulaire technique approprié.
\end{itemize}
\end{exoresolu}

\begin{exemplebox}[Modèle de justification orale du vote]
\eng{We think the Flame King is the best buy because it is stronger and more durable than the Quick Fire. The guarantee is also longer: one year instead of three months. It is more expensive, but it offers better value for money.}

« Nous pensons que le Flame King est le meilleur achat parce qu'il est plus solide et plus durable que le Quick Fire. La garantie est aussi plus longue : un an au lieu de trois mois. Il est plus cher, mais il offre un meilleur rapport qualité-prix. »
\end{exemplebox}

\begin{savaistu}[Le saviez-vous ? — Consommation au Cameroun]
Au Cameroun, l'expression \eng{“buyam-sellam”} (pidgin/argot) désigne le petit commerce informel. Sur les marchés comme Molyko (Buea), Marché Central (Yaoundé) ou Nkoulouloun (Douala), la négociation (\eng{bargaining / haggling}) est la norme. Un bon consommateur compare toujours : \eng{“Cut price!”} (Réduisez le prix !). Les garanties formelles (\eng{warranty}) sont rares dans l'informel ; on parle plutôt de \eng{“test before buy”} (essaie avant d'acheter). Les radios communautaires (comme \eng{Buea Community Radio}, \eng{Radio Environnement} à Yaoundé) jouent un rôle clé d'éducation du consommateur en langues locales et en anglais/français.
\end{savaistu}

\courssub{Bilan de l'activité}

Cette activité d'intégration a permis de réunir les quatre savoir-faire de la leçon dans une situation de communication authentique :
\begin{itemize}
\item \textbf{Compréhension orale (Listening)} : écouter les reportages des autres groupes et relever les chiffres exacts dans la grille d'écoute (prix, durée de garantie, matériaux).
\item \textbf{Lexique (Vocabulary)} : réemployer les mots de la consommation (\eng{price, quality, guarantee, durable, value for money, best buy, bargain, sturdy}) dans un contexte réel de marché camerounais.
\item \textbf{Grammaire (Grammar)} : mobiliser correctement comparatifs (\eng{cheaper than, more durable than, longer than}) et superlatifs (\eng{the best buy, the most reliable}) pour comparer et justifier.
\item \textbf{Production orale et écrite (Speaking / Writing)} : jouer un rôle à la radio (presenter, reporter, seller), voter en argumentant (\eng{We think... because...}), puis rédiger un compte rendu structuré de six lignes.
\end{itemize}

L'évaluation porte sur trois critères annoncés aux élèves :
\begin{enumerate}
\item La \textbf{justesse des chiffres relevés} (prix en FCFA, durées de garantie, nombre de brûleurs).
\item L'\textbf{usage correct du lexique} du thème (orthographe et sens).
\item La \textbf{correction des comparatifs et superlatifs} (\eng{cheaper than}, \eng{more durable than}, \eng{the best buy}, \eng{better value for money}).
\end{enumerate}

Au-delà de la langue, l'activité développe une compétence de vie (\eng{life skill}) : comparer avant d'acheter, vérifier la garantie et défendre un choix de consommateur averti — exactement l'objectif du chapitre \emph{Economic Life and Occupations}. Pour prolonger l'activité à la maison, chaque élève peut comparer deux produits réels dans un marché de sa localité (téléphones, chaussures, sacs, réchauds) et rédiger un court paragraphe de recommandation en anglais (6--8 lignes) à déposer sur la plateforme OnBuch+ ou à remettre en classe.

\newpage

\coursec{Méthodes et exercices résolus}

\courssub{Fiche méthode 1 — Comprendre un enregistrement sur la comparaison des prix}

\begin{methode}[Réussir l'écoute : repérer thème, mots-clés, chiffres et noms propres]
\begin{enumerate}
\item \textbf{Avant l'écoute} : lis d'abord les questions. Elles annoncent le thème (\eng{prices, quality, market, shop}) et te disent quoi écouter : un prix, un nom de lieu, une opinion.
\item \textbf{Première écoute} : ne note rien. Identifie le \cle{thème général} (qui parle ? où ? de quoi ?) et le ton (conseil, publicité, plainte).
\item \textbf{Deuxième écoute} : note uniquement les \cle{chiffres} (prix, quantités, pourcentages), les \cle{noms propres} (marchés, magasins, personnes) et les \cle{mots-clés} du thème (\eng{cheaper, discount, durable, value for money}).
\item \textbf{Attention aux distracteurs} : un enregistrement cite souvent plusieurs prix ou plusieurs produits. Vérifie \emph{à quoi} chaque chiffre se rapporte avant de répondre.
\item \textbf{Rédaction de la réponse} : réponds par une phrase complète en anglais, en reprenant les termes de la question. Exemple : \eng{Mrs Bello advises buying rice from the cooperative because it is cheaper and fresher.}
\end{enumerate}
\end{methode}

\begin{textebox}[Listening script — The Best Buy]
Mrs Amina Bello wants to buy a fifty-kilogram bag of rice for her family in Douala. At the supermarket near Bonanjo, the bag costs twenty-eight thousand francs CFA. At the Sandaga market, the same quality of rice costs only twenty-four thousand francs, but the bag is sometimes damaged during transport. Her neighbour, Mr Etoundi, tells her that a local farmers' cooperative in Mbanga sells the bag for twenty-two thousand francs, and the rice is fresh and clean. However, the cooperative is far away, and the journey by bus costs one thousand francs each way. After comparing the prices and the quality, Amina decides to buy from the cooperative. “It is the best buy,” she says, “because the rice is good and I still save four thousand francs, even with the transport.”
\end{textebox}

\begin{exoresolu}[Listening comprehension — “The Best Buy”]
\textbf{Énoncé.} Écoute le script ci-dessus lu par le professeur (ou lis-le si l'enregistrement n'est pas disponible), puis réponds aux questions.

\textbf{Questions.}
\begin{enumerate}
\item What does Mrs Bello want to buy?
\item How much does the bag of rice cost at the supermarket?
\item True or False: the rice at the cooperative is of poor quality.
\item Why does Amina choose the cooperative? Give two reasons.
\end{enumerate}
\tcblower
\textbf{Solution détaillée.}
\begin{enumerate}
\item La question porte sur l'objet de l'achat. Le script dit : \eng{“a fifty-kilogram bag of rice”}. Réponse rédigée : \eng{Mrs Bello wants to buy a fifty-kilogram bag of rice for her family.}
\item Question sur un chiffre. Attention au distracteur : trois prix sont cités (28 000, 24 000, 22 000). Le prix du supermarché est le premier. Réponse : \eng{At the supermarket, the bag costs twenty-eight thousand francs CFA.}
\item Vrai/faux. Le script dit : \eng{“the rice is fresh and clean”}. L'affirmation “poor quality” est donc fausse. Réponse : \eng{False. The rice at the cooperative is fresh and clean.}
\item Question à réponse courte avec justification. Deux raisons : le prix le plus bas et la bonne qualité, malgré le transport. Réponse : \eng{Amina chooses the cooperative because the rice is good quality and it is the cheapest option: she still saves four thousand francs even with the transport costs.}
\end{enumerate}
\textbf{Conclusion.} En écoute, chaque chiffre entendu doit être relié à un produit ou à un lieu ; c'est le piège classique de ce type d'épreuve.
\end{exoresolu}

\courssub{Fiche méthode 2 — Relever les informations essentielles et les détails}

\begin{methode}[Construire une grille de comparaison pendant l'écoute]
\begin{enumerate}
\item Trace mentalement (ou au brouillon) un tableau à trois colonnes : \cle{produit/lieu}, \cle{prix}, \cle{qualité}.
\item À chaque nom de lieu entendu (\eng{supermarket, market, cooperative}), place le prix et le commentaire sur la qualité dans la bonne ligne.
\item Repère les \cle{comparatifs} et \cle{superlatifs} : \eng{cheaper than}, \eng{the cheapest}, \eng{better value}, \eng{the best buy}. Ils signalent la conclusion du locuteur.
\item Repère les \cle{connecteurs de contraste} : \eng{but, however, although}. L'information importante (le défaut ou l'avantage décisif) vient souvent après.
\item Conclus en identifiant le \cle{choix final} du locuteur et sa justification : c'est presque toujours la réponse à la question principale.
\end{enumerate}
\end{methode}

\begin{exoresolu}[Relevé d'informations — tableau de comparaison]
\textbf{Énoncé.} À partir du script “The Best Buy”, complète le tableau puis réponds : \emph{Which option offers the best value for money? Justify.}

\begin{center}
\begin{tabularx}{\linewidth}{|X|X|X|}
\hline
\rowcolor{popblueL} Place & Price (FCFA) & Quality comment \\
\hline
Supermarket (Bonanjo) & \ldots & \ldots \\
\hline
Sandaga market & \ldots & \ldots \\
\hline
Cooperative (Mbanga) & \ldots & \ldots \\
\hline
\end{tabularx}
\end{center}
\tcblower
\textbf{Solution détaillée.}

\begin{center}
\begin{tabularx}{\linewidth}{|X|X|X|}
\hline
\rowcolor{popblueL} Place & Price (FCFA) & Quality comment \\
\hline
Supermarket (Bonanjo) & 28,000 & same quality, but most expensive \\
\hline
Sandaga market & 24,000 & same quality, but bag sometimes damaged \\
\hline
Cooperative (Mbanga) & 22,000 & fresh and clean; far away (2,000 FCFA return trip) \\
\hline
\end{tabularx}
\end{center}

Raisonnement : il faut comparer le \emph{coût total}, pas seulement le prix affiché. Coopérative : 22 000 + 2 000 de transport = 24 000 FCFA, soit le même total que Sandaga, mais avec une meilleure qualité (\eng{fresh and clean}) et sans risque de sac endommagé.

Réponse rédigée : \eng{The cooperative offers the best value for money. Even with the bus fare, the total cost is twenty-four thousand francs, like at Sandaga market, but the rice is fresh, clean and undamaged. That is why Amina calls it “the best buy”.}

\textbf{Conclusion.} Le “meilleur achat” n'est pas toujours le prix le plus bas : il faut croiser \cle{prix}, \cle{qualité} et \cle{coûts annexes} (transport, durabilité).
\end{exoresolu}

\courssub{Fiche méthode 3 — Maîtriser le lexique de la consommation et de la comparaison}

\begin{definition}[Lexique essentiel du thème]
\begin{itemize}
\item \eng{price} (nom) : le prix ; \eng{cheap} (adj.) : bon marché ; \eng{expensive} (adj.) : cher.
\item \eng{quality} (nom indénombrable) : la qualité ; \eng{value for money} : le rapport qualité-prix.
\item \eng{a bargain / a good deal} : une bonne affaire ; \eng{the best buy} : le meilleur achat.
\item \eng{to compare A with/to B} : comparer A à B ; \eng{to shop around} : comparer les prix dans plusieurs magasins.
\item \eng{to save money} : économiser de l'argent ; \eng{savings} (pluriel) : les économies.
\item \eng{a discount} : une remise, une réduction ; \eng{to buy in bulk} : acheter en gros.
\item \eng{durable / strong} : durable, solide ; \eng{to last} : durer ; \eng{to tear} (tore, torn) : se déchirer.
\end{itemize}
\end{definition}

\begin{methode}[Réemployer le vocabulaire des achats dans des phrases correctes]
\begin{enumerate}
\item Apprends les mots par \cle{familles} : \eng{price – cheap – cheaper – the cheapest} ; \eng{quality – good – better – the best} ; \eng{to compare – comparison}.
\item Mémorise les \cle{expressions figées} du thème : \eng{value for money} (rapport qualité-prix), \eng{a good deal / a bargain} (une bonne affaire), \eng{the best buy} (le meilleur achat), \eng{to shop around} (comparer les magasins), \eng{to save money} (économiser).
\item Respecte la structure du comparatif : \eng{cheaper than}, \eng{more expensive than}, jamais \eng{more cheap}.
\item Pour réemployer : prends un mot du texte, puis construis une phrase sur ta propre expérience (marché de Mokolo, supermarché de Douala, boutique du quartier).
\item Vérifie l'accord et l'article : \eng{rice} et \eng{money} sont indénombrables → \eng{rice is}, \eng{money is}, sans \eng{a}.
\end{enumerate}
\end{methode}

\begin{exoresolu}[Vocabulary in use — phrases personnelles]
\textbf{Énoncé.} Complete the sentences with the correct word from the box, then write two sentences of your own.

\begin{center}
\eng{bargain — value for money — cheaper — save — quality}
\end{center}

\begin{enumerate}
\item \eng{Tomatoes are \ldots\ at the market than at the supermarket.}
\item \eng{This phone is expensive, but it offers good \ldots\ because it lasts five years.}
\item \eng{I bought these shoes for 2,000 francs: what a \ldots !}
\item \eng{If you buy in bulk, you can \ldots\ a lot of money.}
\item \eng{The \ldots\ of this fabric is excellent: it does not fade.}
\end{enumerate}
\tcblower
\textbf{Solution détaillée.}
\begin{enumerate}
\item \eng{cheaper} — comparatif de \eng{cheap} (adjectif court : -er), suivi de \eng{than}. Phrase complète : \eng{Tomatoes are cheaper at the market than at the supermarket.}
\item \eng{value for money} — expression figée signifiant « rapport qualité-prix » : \eng{This phone is expensive, but it offers good value for money because it lasts five years.}
\item \eng{bargain} — « une bonne affaire » ; exclamation typique : \eng{What a bargain!}
\item \eng{save} — \eng{to save money} = économiser ; après \eng{can}, base verbale : \eng{If you buy in bulk, you can save a lot of money.}
\item \eng{quality} — nom indénombrable ici, avec article défini : \eng{The quality of this fabric is excellent: it does not fade.}
\end{enumerate}
Phrases personnelles attendues (modèles) : \eng{At Mokolo market, second-hand clothes are cheaper than in town, and the quality is often good.} / \eng{My mother always shops around before she buys plantains, so she saves money every week.}
\textbf{Conclusion.} Un mot n'est acquis que s'il est réemployé dans une phrase personnelle, grammaticalement correcte (comparatif, article, verbe).
\end{exoresolu}

\courssub{Fiche méthode 4 — Répondre aux questions de compréhension (vrai/faux, QCM, réponses courtes)}

\begin{methode}[Traiter chaque type de question comme au Bac]
\begin{enumerate}
\item \textbf{Vrai/Faux} : justifie toujours par une courte citation du texte (\eng{True: “the rice is fresh and clean”}). Une réponse sans justification perd des points.
\item \textbf{QCM} : élimine d'abord les options contredites par le texte, puis méfie-toi de l'option qui reprend les mots exacts du texte mais change le sens (distracteur classique).
\item \textbf{Réponse courte} : reprends les termes de la question pour former une phrase complète ; ne recopie jamais un bloc du texte sans l'adapter (personne, temps).
\item \textbf{Question “Why?”} : commence par \eng{because} + sujet + verbe, et donne autant de raisons que demandé (\eng{two reasons} = deux arguments distincts).
\item Relis-toi : verbe conjugué à la bonne personne, présent simple pour les faits du texte, passé pour les actions racontées (\eng{Amina decided}, \eng{she says}).
\end{enumerate}
\end{methode}

\begin{exoresolu}[Mixed questions — type Bac]
\textbf{Énoncé.} Answer according to the script “The Best Buy”.

\begin{enumerate}
\item True or False? Justify. \eng{The supermarket is the cheapest place to buy rice.}
\item Choose the correct answer. The journey to the cooperative costs: a) 1,000 FCFA in total; b) 2,000 FCFA in total; c) 4,000 FCFA in total.
\item How much money does Amina finally save compared with the supermarket?
\end{enumerate}
\tcblower
\textbf{Solution détaillée.}
\begin{enumerate}
\item \eng{False. The supermarket is the most expensive option: the bag costs twenty-eight thousand francs there, while the cooperative sells it for twenty-two thousand.} — Justification obligatoire par les chiffres du texte.
\item Piège du calcul : le script dit \eng{“one thousand francs each way”} (1 000 par trajet). Aller-retour = 2 000 FCFA. Réponse : \eng{b) 2,000 FCFA in total, because the bus costs one thousand francs each way.}
\item Calcul : supermarché 28 000 ; coopérative 22 000 + 2 000 de transport = 24 000. Économie : 28 000 − 24 000 = 4 000. Réponse rédigée : \eng{Amina finally saves four thousand francs compared with the supermarket, even after paying for the bus.}
\end{enumerate}
\textbf{Conclusion.} Au Bac, la réponse juste mais non justifiée, ou juste mais mal rédigée, ne rapporte pas tous les points : toujours \cle{réponse + preuve chiffrée + phrase complète}.
\end{exoresolu}

\courssub{Fiche méthode 5 — Production guidée : conseiller le meilleur achat}

\begin{methode}[Structurer un conseil d'achat à l'oral et à l'écrit]
\begin{enumerate}
\item \textbf{Présente le besoin} : \eng{You want to buy... / Your family needs...}
\item \textbf{Compare au moins deux options} avec comparatifs : \eng{Option A is cheaper, but option B is better quality.}
\item \textbf{Donne ton conseil} avec une structure de recommandation : \eng{I advise you to... / You should... / If I were you, I would buy...}
\item \textbf{Justifie} par deux arguments (prix + qualité, ou prix + durabilité) : \eng{because it is the best value for money and it will last longer.}
\item \textbf{Conclus} par la formule du thème : \eng{It is the best buy. / It is a real bargain.}
\end{enumerate}
\end{methode}

\begin{exoresolu}[Guided production — advising a friend]
\textbf{Énoncé.} Your friend has 15,000 FCFA to buy a school bag. Shop A sells a bag for 8,000 FCFA but it tears after three months. Shop B sells a strong leather bag for 14,000 FCFA that lasts three years. Write a short paragraph (5–6 lines) advising your friend on the best buy.
\tcblower
\textbf{Solution détaillée.}

Démarche : (1) besoin = un sac d'école ; (2) comparaison = prix contre durabilité ; (3) conseil avec \eng{I advise you to} ; (4) justification chiffrée : 8 000 × 4 sacs par an × 3 ans = 96 000 FCFA, ou plus simplement 8 000 × 4 = 32 000 FCFA la première année, contre 14 000 une seule fois ; (5) conclusion avec \eng{the best buy}.

Réponse rédigée (modèle) :

\eng{You need a new school bag, and you have two options. The bag from Shop A is much cheaper, but it tears after only three months, so you will have to buy another one again and again. The leather bag from Shop B is more expensive, but it is strong and it lasts three years. I advise you to buy the leather bag: in one year alone, the cheap bags would cost you thirty-two thousand francs, while the leather bag costs only fourteen thousand and lasts three years. It is better quality and better value for money. Believe me, it is the best buy.}

Points de grammaire vérifiés : comparatif \eng{cheaper / more expensive} ; \eng{advise you to buy} (advise + objet + to-infinitif) ; futur avec \eng{will} pour la conséquence ; indénombrable \eng{money} sans article.
\textbf{Conclusion.} Un bon conseil d'achat = comparaison + recommandation + justification chiffrée + formule de conclusion.
\end{exoresolu}

\begin{attention}[Les 5 erreurs qui coûtent des points au Bac]
\begin{enumerate}
\item \textbf{Faux ami} : \eng{économiser} se dit \eng{to save money}, jamais \eng{to economize money} ; et \eng{des économies} (somme épargnée) = \eng{savings}, pas \eng{economies} (qui désigne les économies des pays).
\item \textbf{Comparatif mal formé} : on écrit \eng{cheaper than} ou \eng{more expensive than}, jamais \eng{more cheaper} ni \eng{more cheap}. De même : \eng{good – better – the best}, jamais \eng{gooder} ou \eng{the most good}.
\item \textbf{Indénombrables} : \eng{rice, money, advice, information} n'ont pas de pluriel et prennent le verbe au singulier : \eng{The rice is cheap}, \eng{the advice is useful} — jamais \eng{rices are} ou \eng{advices}.
\item \textbf{Prépositions} : on dit \eng{to buy something \textbf{at} the market}, \eng{to compare A \textbf{with/to} B}, \eng{to listen \textbf{to} the recording}, \eng{to spend money \textbf{on} food}. Le calque du français (\eng{listen the radio}) est sanctionné.
\item \textbf{Ordre des mots et temps} : pas d'inversion dans la question indirecte — \eng{She asks how much the bag costs} (et non \eng{how much does the bag cost}) ; et dans le compte rendu d'un texte, respecte la concordance : \eng{Amina decided to buy from the cooperative because it was the cheapest.}
\end{enumerate}
\end{attention}

\newpage

\coursec{Exercices}

\coursec{Lesson 13 — Listening: Comparing Prices and Quality for Best Buys}

\courssub{Warm-up et découverte du thème}

\begin{experience}[Discussion de mise en route]
Le professeur affiche (ou dicte) trois offres pour un même produit (ex. sac de riz 50 kg) et demande aux élèves, par paires, de décider laquelle ils choisiraient et pourquoi. Vocabulaire attendu : \eng{price, quality, discount, guarantee, value for money, bargain, expensive, cheap, afford}.
\begin{itemize}
\item \textbf{Shop A :} 28,000 FCFA, marque connue, garantie 1 an.
\item \textbf{Shop B :} 25,500 FCFA, marque inconnue, pas de garantie.
\item \textbf{Shop C :} 27,000 FCFA, marque correcte, \eng{free delivery} (livraison gratuite) + 1 kg de sel offert.
\end{itemize}
Quelques paires présentent leur choix à la classe en anglais. Le professeur note au tableau les mots-clés : \cle{comparing}, \cle{price}, \cle{quality}, \cle{best buy}, \cle{value for money}.
\end{experience}

\courssub{Lexique thématique : consommation et comparaison}

\begin{definition}[Mots et expressions clés du thème]
Le tableau ci-dessous regroupe le vocabulaire essentiel pour comprendre, comparer et exprimer un choix d'achat. Apprenez la prononciation approximative (entre guillemets français) et la traduction.
\begin{center}
\begin{tabularx}{\linewidth}{|X|X|X|}
\hline
\rowcolor{popblueL} \textbf{English} & \textbf{Prononciation (approx.)} & \textbf{Français} \\
\hline
a bargain & « ba-gaïne » & une bonne affaire, un achat avantageux \\
to bargain (with sb) & « ba-gaïne » & marchander, négocier le prix \\
a discount & « dis-kaonte » & une remise, une réduction \\
a receipt & « ri-si-te » & un reçu, une facture \\
a guarantee / a warranty & « ga-ran-ti » / « wo-ran-ti » & une garantie \\
value for money & « va-lyu fo ma-ni » & rapport qualité-prix \\
affordable & « a-for-da-beul » & abordable, à prix raisonnable \\
expensive & « ik-spen-siv » & cher \\
cheap & « tchiip » & bon marché, peu cher \\
quality & « quo-li-ti » & qualité \\
to compare & « kom-païe » & comparer \\
a customer / a shopper & « kas-to-me » / « cho-pe » & un client, un acheteur \\
a seller / a shopkeeper & « se-le » / « cho-pe-ki-pe » & un vendeur \\
to afford & « a-ford » & avoir les moyens de s'offrir \\
to save money & « seiv ma-ni » & économiser de l'argent \\
a price tag & « prais tæg » & une étiquette de prix \\
on sale / on offer & « on seil » / « on o-fe » & en solde, en promotion \\
\hline
\end{tabularx}
\end{center}
\end{definition}

\begin{attention}[Faux amis et pièges fréquents]
\begin{itemize}
\item \textbf{Bargain} (nom) = une bonne affaire. \textbf{Bargain} (verbe) = marchander. Ne pas confondre avec \eng{to argue} (se disputer).
\item \textbf{Receipt} \eng{/rɪˈsiːt/} (reçu) : le \textbf{p} ne se prononce pas. Ne pas confondre avec \eng{recipe} (recette de cuisine).
\item \textbf{Afford} s'emploie presque toujours avec un modal (\eng{can/could afford}). \eng{*I afford it} est faux. Dire \eng{I can't afford it}.
\item \textbf{Cheap} a parfois une connotation négative (mauvaise qualité). Pour dire « prix raisonnable » de façon positive, préférez \eng{affordable} ou \eng{good value}.
\item \textbf{Quality} est un nom non comptable. On ne dit pas \eng{*a good quality} mais \eng{good quality} ou \eng{high quality}.
\end{itemize}
\end{attention}

\courssub{Stratégies et méthode d'écoute}

\begin{methode}[Réussir une épreuve de compréhension orale (Listening)]
\begin{enumerate}
\item \textbf{Avant l'écoute :} Lisez les questions et les options (QCM, Vrai/Faux). Soulignez les \cle{mots-clés} (noms propres, chiffres, prix, noms de magasins, adjectifs de comparaison). Prédisez le thème.
\item \textbf{Première écoute (globale) :} Identifiez le thème général, les interlocuteurs, le lieu. Ne paniquez pas si vous ne comprenez pas tout.
\item \textbf{Deuxième écoute (détaillée) :} Concentrez-vous sur les informations ciblées par les questions : \textbf{chiffres} (prix, pourcentages, quantités), \textbf{noms propres} (magasins, personnes), \textbf{connecteurs logiques} (\eng{however, but, because, so, in the end}).
\item \textbf{Après l'écoute :} Vérifiez vos réponses. Pour les questions \eng{True/False}, citez le passage exact du script (mots entendus) qui justifie votre choix.
\end{enumerate}
\end{methode}

\courssub{Texte support : script d'écoute « The Best Buy »}

\begin{textebox}[Listening Script — Main Text (to be read aloud twice by the teacher)]
“Good morning everyone. Today, let's talk about smart shopping. Last Saturday, Mr. Nguema wanted to buy a new gas cooker for his restaurant in Buea. He visited three shops in the main market.

At \textbf{Kitchen Plus}, the cooker cost 120,000 FCFA. It was a famous brand, very modern, with a two-year guarantee. But Mr. Nguema thought it was too expensive for his budget.

At \textbf{Cheap Deals}, the price was only 85,000 FCFA. The seller said: ‘This is a real bargain, sir!’ However, when Mr. Nguema checked the quality, he saw the metal was thin and there was no guarantee at all. He knew that ‘cheap things are expensive’ — it would break quickly.

Finally, at \textbf{Home Essentials}, the cooker was 95,000 FCFA. It was not the cheapest, but the quality was good, it had a one-year guarantee, and the shopkeeper offered a 5\% discount for cash payment. Mr. Nguema compared the three options: price, quality, guarantee and service. He chose Home Essentials because it offered the \textbf{best value for money}. Remember: a smart consumer always compares before buying.”
\end{textebox}

\courssub{Compréhension : questions de types variés}

\begin{exercice}[QCM : comprendre une annonce de prix]
Votre professeur lit (ou vous lisez) l'annonce suivante deux fois. Écoutez, puis choisissez la bonne réponse.

\begin{textebox}[Listening script — read aloud twice]
“Good morning, shoppers! This weekend at City Market, all school bags are reduced by twenty per cent. A standard bag now costs eight thousand five hundred francs instead of ten thousand. Sports shoes are down from fifteen thousand to twelve thousand francs. And don't forget: every customer who spends more than twenty thousand francs receives a free water bottle. Offer valid until Sunday evening.”
\end{textebox}

\begin{enumerate}
\item The announcement is about: \quad a) a job offer \quad b) a sales promotion \quad c) a new shop opening \quad d) a transport strike
\item School bags are reduced by: \quad a) 10 \% \quad b) 15 \% \quad c) 20 \% \quad d) 25 \%
\item A standard bag now costs: \quad a) 8,000 FCFA \quad b) 8,500 FCFA \quad c) 10,000 FCFA \quad d) 12,000 FCFA
\item Sports shoes used to cost: \quad a) 12,000 FCFA \quad b) 13,000 FCFA \quad c) 15,000 FCFA \quad d) 20,000 FCFA
\item To receive a free water bottle, you must spend more than: \quad a) 8,500 FCFA \quad b) 10,000 FCFA \quad c) 15,000 FCFA \quad d) 20,000 FCFA
\end{enumerate}
\end{exercice}

\begin{exercice}[Vrai ou faux, à justifier]
Écoutez le script ci-dessous lu deux fois, puis dites si les affirmations sont \eng{true} ou \eng{false}. Justifiez chaque réponse par une courte citation du texte en anglais.

\begin{textebox}[Listening script — read aloud twice]
“Madam Etoundi compared three shops before buying a bag of rice. At Shop A, a fifty-kilogram bag costs twenty-eight thousand francs. At Shop B, the same bag costs twenty-six thousand five hundred francs, but the rice is of lower quality. At Shop C, the price is twenty-seven thousand francs and the shopkeeper gives a free two-kilogram packet of salt. In the end, Madam Etoundi chose Shop C because it offered the best value for money.”
\end{textebox}

\begin{enumerate}
\item Madam Etoundi visited only two shops.
\item Shop A is the most expensive.
\item Shop B sells high-quality rice.
\item Shop C offers a free gift.
\item Madam Etoundi bought the cheapest bag of rice.
\end{enumerate}
\end{exercice}

\begin{exercice}[Vocabulaire : le bon mot]
Complétez chaque phrase avec le mot qui convient, choisi dans la liste : \eng{bargain — discount — receipt — quality — affordable — expensive}.

\begin{enumerate}
\item This phone is too \ldots\ldots\ldots for me; I only have 40,000 FCFA.
\item Always keep your \ldots\ldots\ldots in case you want to exchange the product.
\item The shop offers a ten per cent \ldots\ldots\ldots to all students.
\item Don't look at the price only; check the \ldots\ldots\ldots of the material.
\item At Mokolo market, you can \ldots\ldots\ldots with the sellers to get a lower price.
\item These shoes are good and \ldots\ldots\ldots: only 6,000 FCFA.
\end{enumerate}
\end{exercice}

\courssub{Grammaire : comparatifs et superlatifs}

\begin{propriete}[Rappel : formes comparatives et superlatives]
Pour comparer prix, qualité, services, on utilise les degrés de comparaison.
\begin{center}
\begin{tabular}{|l|l|l|l|}
\hline
\rowcolor{popblueL} \textbf{Type d'adjectif} & \textbf{Comparatif de supériorité} & \textbf{Superlatif} & \textbf{Exemple (prix/qualité)} \\
\hline
Court (1 syllabe) & \eng{cheap $\rightarrow$ cheaper than} & \eng{the cheapest} & \eng{Shop B is cheaper than Shop A.} \\
\hline
Court terminé par \emph{-e} & \eng{large $\rightarrow$ larger than} & \eng{the largest} & \eng{This shop has the largest choice.} \\
\hline
Court CVC (double consonne) & \eng{big $\rightarrow$ bigger than} & \eng{the biggest} & \eng{The biggest discount is 50\%.} \\
\hline
Terminé par \emph{-y} & \eng{busy $\rightarrow$ busier than} & \eng{the busiest} & \eng{Market days are busier than weekdays.} \\
\hline
Long ($\ge$ 2 syllabes) & \eng{expensive $\rightarrow$ more expensive than} & \eng{the most expensive} & \eng{This brand is more expensive.} \\
\hline
\multicolumn{4}{|l|}{\textbf{Irréguliers importants :}} \\
\hline
\eng{good / well} & \eng{better than} & \eng{the best} & \eng{This rice is better quality.} \\
\eng{bad / badly} & \eng{worse than} & \eng{the worst} & \eng{Shop B has the worst service.} \\
\eng{far} & \eng{farther / further than} & \eng{the farthest / furthest} & \eng{Market 3 is the furthest.} \\
\eng{much / many} & \eng{more than} & \eng{the most} & \eng{More customers came today.} \\
\eng{little} & \eng{less than} & \eng{the least} & \eng{The least expensive option.} \\
\hline
\end{tabular}
\end{center}
\textbf{Comparatif d'égalité :} \eng{as + adjectif + as} (ex. \eng{Shop C is as good as Shop A}).
\textbf{Comparatif d'infériorité :} \eng{less + adjectif + than} (ex. \eng{less expensive than}) ou \eng{not as + adjectif + as}.
\end{propriete}

\begin{exercice}[Comparatifs et superlatifs à compléter]
Complétez les phrases avec la forme correcte de l'adjectif entre parenthèses (comparatif ou superlatif).

\begin{enumerate}
\item Shop B is \ldots\ldots\ldots (cheap) than Shop A.
\item This is the \ldots\ldots\ldots (good) quality rice in the market.
\item The new bus is \ldots\ldots\ldots (comfortable) than the old one.
\item Amadou Travel is the \ldots\ldots\ldots (expensive) company on the Douala–Yaoundé route.
\item My phone has a \ldots\ldots\ldots (long) guarantee than yours.
\item This market is the \ldots\ldots\ldots (far) from our school, but prices are the \ldots\ldots\ldots (low).
\end{enumerate}
\end{exercice}

\courssub{Production guidée : conseiller le meilleur achat}

\begin{exercice}[Traduction vers l'anglais]
Traduisez les phrases suivantes en anglais, en utilisant le vocabulaire et la grammaire de la leçon.

\begin{enumerate}
\item Ce magasin est moins cher que l'autre, mais la qualité est meilleure.
\item Le vendeur m'a fait une remise de quinze pour cent.
\item Compare toujours les prix avant d'acheter.
\item Cette agence de voyage offre le meilleur rapport qualité-prix.
\item J'ai marchandé et j'ai payé le sac douze mille francs au lieu de quinze mille.
\end{enumerate}
\end{exercice}

\begin{exercice}[Transformation : comparer trois boutiques]
Transformez les phrases selon le modèle : \eng{Shop A is cheap. $\rightarrow$ Shop B is cheaper than Shop A. $\rightarrow$ Shop C is the cheapest of the three.}

\begin{enumerate}
\item Shop A is expensive. (Shop B / Shop C)
\item The first phone is good. (the second phone / the third phone)
\item Bus company X is fast. (company Y / company Z)
\item Market 1 is crowded. (Market 2 / Market 3)
\item This radio is bad. (that radio / the new radio)
\end{enumerate}
\end{exercice}

\begin{exercice}[Texte à trous : au marché]
Complétez le dialogue avec les mots de la liste : \eng{how much — cheaper — discount — quality — receipt — afford — bargain — price}.

\begin{textebox}[At the market]
Customer: Good morning, madam. \ldots\ldots\ldots is this bag of onions?

Seller: It is five thousand francs, my dear.

Customer: Five thousand! That is too high. At the other stand, the \ldots\ldots\ldots is four thousand five hundred.

Seller: But look at the \ldots\ldots\ldots of my onions! They are fresh and big.

Customer: Can you give me a small \ldots\ldots\ldots? I am a student; I cannot \ldots\ldots\ldots to pay so much.

Seller: All right, because I like you, take it for four thousand five hundred. But don't tell anyone!

Customer: Thank you! You really know how to \ldots\ldots\ldots. Can I have a \ldots\ldots\ldots, please?

Seller: Of course. Here you are. Come back next week — my tomatoes will be \ldots\ldots\ldots than today!
\end{textebox}
\end{exercice}

\begin{exercice}[Appariement : mots et définitions]
Associez chaque mot anglais (colonne de gauche) à sa définition anglaise (colonne de droite).

\begin{center}
\begin{tabularx}{\linewidth}{|X|X|}
\hline
\rowcolor{popblueL} \textbf{Word} & \textbf{Definition} \\
\hline
1. a bargain & a) a paper that proves you paid for something \\
\hline
2. a discount & b) to discuss the price in order to pay less \\
\hline
3. a receipt & c) a product bought at a very good price \\
\hline
4. a guarantee & d) a reduction on the normal price \\
\hline
5. to bargain & e) how good or bad something is \\
\hline
6. quality & f) a written promise to repair or replace a product \\
\hline
\end{tabularx}
\end{center}
\end{exercice}

\courssub{Je me prépare au Bac}

\begin{exercice}[Compréhension de texte — type Baccalauréat]
Lisez le texte suivant, puis répondez aux questions.

\begin{textebox}[Choosing a bus company in Cameroon]
Every Friday, hundreds of students and traders travel between Douala and Yaoundé. Before buying a ticket, wise passengers compare the three main bus companies on the road.

Express Line is the most famous. A ticket costs 5,000 FCFA, and the buses are new and air-conditioned. However, the company is often fully booked, so passengers must buy their tickets two days in advance.

Comfort Voyages is slightly cheaper: 4,500 FCFA per ticket. The buses leave every hour, which is very convenient, but they are older and the journey takes thirty minutes longer because the driver stops in every town to pick up passengers.

Finally, Budget Trans is the cheapest option at only 3,800 FCFA. Students love it, but the buses are slow, there is no air-conditioning, and luggage is charged extra — 500 FCFA per bag.

In the end, the best buy depends on what matters most to you: speed, comfort or price. As the saying goes, “cheap things are expensive” — a very cheap ticket can cost you time and comfort.
\end{textebox}

\textbf{A. Comprehension questions}

\begin{enumerate}
\item Say whether the following statements are \eng{true} or \eng{false}. Justify each answer with a short quotation from the text.
\begin{enumerate}
\item Express Line buses have air-conditioning.
\item You can always buy an Express Line ticket on the day of travel.
\item Comfort Voyages buses stop in many towns.
\item Budget Trans carries your luggage for free.
\end{enumerate}
\item Answer these questions in your own words.
\begin{enumerate}
\item Which company is the cheapest, and how much is a ticket?
\item Give one advantage and one disadvantage of Comfort Voyages.
\item Explain in one sentence the meaning of the saying “cheap things are expensive”.
\end{enumerate}
\end{enumerate}

\textbf{B. Language work}

\begin{enumerate}
\item Find in the text the comparative or superlative form of: \eng{famous}, \eng{cheap}, \eng{long}, \eng{good}.
\item Complete: Budget Trans is \ldots\ldots\ldots (cheap) company, but it is also the \ldots\ldots\ldots (slow).
\end{enumerate}
\end{exercice}

\begin{exercice}[Traduction — type Baccalauréat]
Traduisez le passage suivant en anglais.

\begin{textebox}[À traduire]
Avant d'acheter un nouveau téléphone, mon frère a comparé les prix dans trois boutiques de Yaoundé. La première boutique était la plus chère, mais elle offrait une garantie de deux ans. La deuxième était moins chère, mais le téléphone était de mauvaise qualité. Finalement, il a choisi la troisième boutique, qui proposait le meilleur rapport qualité-prix et une remise de dix pour cent pour les étudiants.
\end{textebox}
\end{exercice}

\begin{exercice}[Composition — type Baccalauréat]
\textbf{Topic:} Your friend wants to buy a smartphone and hesitates between two models: Phone A (45,000 FCFA, one-year guarantee, average camera) and Phone B (60,000 FCFA, two-year guarantee, excellent camera, free earphones).

Write a paragraph of about \textbf{80 words} advising your friend on the best buy. Compare the \textbf{price}, the \textbf{quality} and the \textbf{guarantee} of the two phones, and give your final advice. Use at least \textbf{two comparatives} and \textbf{one superlative}.
\end{exercice}

\begin{exemplebox}[Modèle de rédaction (environ 85 mots)]
“Hi! If you want my advice, go for Phone B. True, Phone A is \textbf{cheaper than} Phone B (45,000 vs 60,000 FCFA), but Phone B offers \textbf{better quality} overall. Its camera is \textbf{far superior to} Phone A's, and the \textbf{two-year guarantee} gives you more security than the one-year one. Plus, you get free earphones! Considering the camera, the guarantee and the gift, Phone B is \textbf{the best value for money}. It's a smarter investment for the long term.”
\end{exemplebox}

\begin{aretenir}[Check-list pour la composition]
\begin{itemize}
\item \textbf{Structure :} Choix clair $\rightarrow$ Comparaison Prix / Qualité / Garantie $\rightarrow$ Conseil final.
\item \textbf{Grammaire :} Au moins 2 comparatifs (\eng{cheaper than, better than, more expensive than, superior to}) et 1 superlatif (\eng{the best, the most reliable}).
\item \textbf{Vocabulaire :} \eng{value for money, guarantee, discount, affordable, bargain, invest, worth it}.
\item \textbf{Connecteurs :} \eng{However, moreover, therefore, in conclusion, as a result}.
\end{itemize}
\end{aretenir}

\begin{savaistu}[Le saviez-vous ? — Le commerce au Cameroun]
Au Cameroun, le secteur informel représente une part énorme de l'économie. Les marchés comme \textbf{Mokolo} (Yaoundé), \textbf{Sandaga} (Douala) ou \textbf{Mvog-Ada} sont des lieux où la négociation (\eng{bargaining/haggling}) est la norme. Contrairement aux supermarchés (Carrefour, U Market, Mahima) où les prix sont fixes (\eng{fixed prices}), au marché “\eng{the price is not written on the tag}”. L'expression \eng{“to strike a bargain”} signifie “conclure une bonne affaire”. Le \textbf{e-commerce} (Jumia, Glotelho, Kaymu) change la donne : on compare les prix en ligne (\eng{online price comparison}) avant d'aller en boutique.
\end{savaistu}


\newpage

\coursec{Corrigés des exercices}

\begin{corrige}[Exercice 1 — QCM : comprendre une annonce de prix]
\begin{enumerate}
\item \textbf{b) a sales promotion.} Mots-clés entendus : “reduced by twenty per cent”, “down from”, “offer valid until Sunday”. Il s'agit d'une promotion commerciale, pas d'une offre d'emploi ni d'une ouverture de magasin.
\item \textbf{c) 20 \%.} Justification : “all school bags are reduced by \textbf{twenty per cent}”.
\item \textbf{b) 8,500 FCFA.} Justification : “A standard bag now costs \textbf{eight thousand five hundred francs} instead of ten thousand.”
\item \textbf{c) 15,000 FCFA.} Justification : “Sports shoes are down \textbf{from fifteen thousand} to twelve thousand francs.” (\eng{used to cost} = coûtaient avant.)
\item \textbf{d) 20,000 FCFA.} Justification : “every customer who spends \textbf{more than twenty thousand francs} receives a free water bottle.”
\end{enumerate}
\textbf{Point de vigilance :} ne pas confondre le prix \emph{avant} (\eng{instead of}, \eng{down from}) et le prix \emph{après} réduction (\eng{now costs}).
\end{corrige}

\begin{corrige}[Exercice 2 — Vrai ou Faux, à justifier]
\begin{enumerate}
\item \textbf{False.} Justification : “Madam Etoundi compared \textbf{three shops} before buying a bag of rice.”
\item \textbf{True.} Justification : “At Shop A, a fifty-kilogram bag costs \textbf{twenty-eight thousand francs}” — contre 26,500 (Shop B) et 27,000 (Shop C). Shop A est donc le plus cher.
\item \textbf{False.} Justification : “…but the rice is of \textbf{lower quality}.” (Shop B vend du riz de qualité inférieure.)
\item \textbf{True.} Justification : “…the shopkeeper gives a \textbf{free two-kilogram packet of salt}.”
\item \textbf{False.} Justification : “In the end, Madam Etoundi \textbf{chose Shop C} because it offered the \textbf{best value for money}.” Le sac le moins cher était au Shop B (26,500 FCFA), qu'elle n'a pas choisi.
\end{enumerate}
\textbf{Point de vigilance :} \eng{the best value for money} (meilleur rapport qualité-prix) ne signifie pas \eng{the cheapest} (le moins cher).
\end{corrige}

\begin{corrige}[Exercice 3 — Vocabulaire : le bon mot]
\begin{enumerate}
\item \textbf{expensive} — “This phone is too \textbf{expensive} for me; I only have 40,000 FCFA.” (trop cher pour mon budget)
\item \textbf{receipt} — “Always keep your \textbf{receipt} in case you want to exchange the product.” (le reçu, preuve d'achat ; le \emph{p} ne se prononce pas)
\item \textbf{discount} — “The shop offers a ten per cent \textbf{discount} to all students.” (une remise)
\item \textbf{quality} — “Don't look at the price only; check the \textbf{quality} of the material.” (nom non comptable, sans article)
\item \textbf{bargain} — “At Mokolo market, you can \textbf{bargain} with the sellers to get a lower price.” (verbe : marchander ; ici après \eng{can}, base verbale)
\item \textbf{affordable} — “These shoes are good and \textbf{affordable}: only 6,000 FCFA.” (abordables, connotation positive)
\end{enumerate}
\end{corrige}

\begin{corrige}[Exercice 4 — Comparatifs et superlatifs à compléter]
\begin{enumerate}
\item Shop B is \textbf{cheaper} than Shop A. (adjectif court : \eng{cheap} + \emph{-er})
\item This is the \textbf{best} quality rice in the market. (irrégulier : \eng{good} $\rightarrow$ \eng{better} $\rightarrow$ \eng{the best})
\item The new bus is \textbf{more comfortable} than the old one. (adjectif long de 3 syllabes : \eng{more} + adjectif)
\item Amadou Travel is the \textbf{most expensive} company on the Douala–Yaoundé route. (superlatif d'adjectif long)
\item My phone has a \textbf{longer} guarantee than yours. (adjectif court : \eng{long} + \emph{-er})
\item This market is the \textbf{furthest} (ou \textbf{farthest}) from our school, but prices are the \textbf{lowest}. (irrégulier : \eng{far} $\rightarrow$ \eng{furthest/farthest} ; adjectif court : \eng{low} $\rightarrow$ \eng{the lowest})
\end{enumerate}
\textbf{Points de vigilance :} ne jamais écrire \eng{*more cheaper} ni \eng{*the most cheapest} (double marque interdite) ; \eng{good} ne donne jamais \eng{*gooder} ni \eng{*the goodest}.
\end{corrige}

\begin{corrige}[Exercice 5 — Traduction vers l'anglais]
\begin{enumerate}
\item “This shop is \textbf{cheaper than} the other one, but the quality is \textbf{better}.” (Variantes acceptées : \eng{less expensive than} ; \eng{the quality is better there}.)
\item “The seller \textbf{gave me a fifteen per cent discount}.” (Variante : \eng{offered me a 15\% discount}.) Attention : on dit \eng{give somebody a discount}, pas \eng{*make a discount to me}.
\item “\textbf{Always compare prices before buying}.” (Variante : \eng{before you buy}.) Impératif : base verbale sans sujet.
\item “This travel agency \textbf{offers the best value for money}.” (Superlatif irrégulier \eng{the best} ; expression figée sans article devant \eng{value}.)
\item “I \textbf{bargained} and \textbf{paid twelve thousand francs for the bag instead of fifteen thousand}.” (Variantes : \eng{haggled} ; \eng{paid 12,000 FCFA instead of 15,000}.) Attention à la préposition : \eng{pay + montant + \textbf{for} + objet}.
\end{enumerate}
\end{corrige}

\begin{corrige}[Exercice 6 — Transformation : comparer trois boutiques]
\begin{enumerate}
\item Shop B is \textbf{more expensive than} Shop A. / Shop C is \textbf{the most expensive of the three}.
\item The second phone is \textbf{better than} the first phone. / The third phone is \textbf{the best of the three}.
\item Company Y is \textbf{faster than} company X. / Company Z is \textbf{the fastest of the three}.
\item Market 2 is \textbf{more crowded than} Market 1. / Market 3 is \textbf{the most crowded of the three}.
\item That radio is \textbf{worse than} this radio. / The new radio is \textbf{the worst of the three}.
\end{enumerate}
\textbf{Points de vigilance :} le comparatif exige \eng{than} (jamais \eng{*that} ni \eng{*as}) ; le superlatif exige l'article \eng{the} et peut être suivi de \eng{of the three} ou \eng{in the market}.
\end{corrige}

\begin{corrige}[Exercice 7 — Texte à trous : au marché]
\begin{enumerate}
\item \textbf{How much} — “\textbf{How much} is this bag of onions?” (question de prix ; \eng{how much} + verbe pour un prix)
\item \textbf{price} — “At the other stand, the \textbf{price} is four thousand five hundred.”
\item \textbf{quality} — “But look at the \textbf{quality} of my onions!”
\item \textbf{discount} — “Can you give me a small \textbf{discount}?”
\item \textbf{afford} — “I am a student; I cannot \textbf{afford} to pay so much.” (structure : \eng{can/cannot afford to do})
\item \textbf{bargain} — “You really know how to \textbf{bargain}.” (verbe après \eng{how to})
\item \textbf{receipt} — “Can I have a \textbf{receipt}, please?”
\item \textbf{cheaper} — “My tomatoes will be \textbf{cheaper} than today!” (comparatif + \eng{than})
\end{enumerate}
\end{corrige}

\begin{corrige}[Exercice 8 — Appariement : mots et définitions]
\begin{center}
\begin{tabular}{|l|l|l|}
\hline
\rowcolor{popblueL} \textbf{Word} & \textbf{Réponse} & \textbf{Definition} \\
\hline
1. a bargain & \textbf{c} & a product bought at a very good price \\
2. a discount & \textbf{d} & a reduction on the normal price \\
3. a receipt & \textbf{a} & a paper that proves you paid for something \\
4. a guarantee & \textbf{f} & a written promise to repair or replace a product \\
5. to bargain & \textbf{b} & to discuss the price in order to pay less \\
6. quality & \textbf{e} & how good or bad something is \\
\hline
\end{tabular}
\end{center}
Donc : \textbf{1–c, 2–d, 3–a, 4–f, 5–b, 6–e.}
\end{corrige}

\begin{corrige}[Exercice 9 — Compréhension de texte, type Baccalauréat]
\textbf{Barème indicatif (sur 10 points) :} Partie A.1 : 4 $\times$ 1 pt (0,5 pt pour la réponse + 0,5 pt pour la citation justificative) ; A.2 : 3 $\times$ 1 pt ; B : 2 $\times$ 1,5 pt.

\textbf{A. Comprehension questions}
\begin{enumerate}
\item Vrai ou faux :
\begin{enumerate}
\item \textbf{True.} “…the buses are new and \textbf{air-conditioned}.”
\item \textbf{False.} “…the company is often \textbf{fully booked}, so passengers \textbf{must buy their tickets two days in advance}.”
\item \textbf{True.} “…the driver \textbf{stops in every town} to pick up passengers.”
\item \textbf{False.} “…luggage is \textbf{charged extra} — 500 FCFA per bag.”
\end{enumerate}
\item Réponses personnelles :
\begin{enumerate}
\item “\textbf{Budget Trans} is the cheapest company; a ticket costs \textbf{3,800 FCFA}.”
\item “One advantage of Comfort Voyages is that \textbf{buses leave every hour}, which is convenient. One disadvantage is that the \textbf{buses are older} and the \textbf{journey takes thirty minutes longer}.”
\item “The saying means that \textbf{a very cheap product can cost you more later}, for example in repairs, lost time or discomfort.”
\end{enumerate}
\end{enumerate}

\textbf{B. Language work}
\begin{enumerate}
\item \eng{famous} $\rightarrow$ \textbf{the most famous} ; \eng{cheap} $\rightarrow$ \textbf{cheaper} / \textbf{the cheapest} ; \eng{long} $\rightarrow$ \textbf{longer} ; \eng{good} $\rightarrow$ \textbf{the best} (dans “the best buy”).
\item Budget Trans is \textbf{the cheapest} company, but it is also \textbf{the slowest}.
\end{enumerate}
\textbf{Points de vigilance :} pour les questions Vrai/Faux, la citation doit être \emph{exacte} et tirée du texte ; aux questions ouvertes, reformulez avec vos propres mots sans recopier tout le paragraphe.
\end{corrige}

\begin{corrige}[Exercice 10 — Traduction, type Baccalauréat]
\textbf{Barème indicatif (sur 5 points) :} 1 pt par segment correct (5 segments), moins 0,25 à 0,5 pt par faute de grammaire ou de vocabulaire.

\textbf{Traduction de référence :}

“\textbf{Before buying} a new phone, my brother \textbf{compared prices} in three shops in Yaoundé. The first shop was \textbf{the most expensive}, but it \textbf{offered a two-year guarantee}. The second one was \textbf{cheaper}, but the phone was \textbf{of poor quality}. Finally, he \textbf{chose} the third shop, which \textbf{offered the best value for money} and \textbf{a ten per cent discount for students}.”

\textbf{Points de vigilance :}
\begin{itemize}
\item \eng{Before buying} (gérondif après \eng{before}), pas \eng{*before to buy}.
\item Superlatifs : \eng{the most expensive}, \eng{the best} — avec l'article \eng{the}.
\item \eng{of poor quality} ou \eng{of low quality} ; éviter \eng{*bad quality} dans un texte soigné.
\item Prétérits réguliers et irréguliers : \eng{compared}, \eng{offered}, \eng{chose} (irrégulier : \eng{choose} $\rightarrow$ \eng{chose} $\rightarrow$ \eng{chosen}).
\item \eng{a ten per cent discount} : pas de « s » à \eng{per cent} ici.
\end{itemize}
\end{corrige}

\begin{corrige}[Exercice 11 — Composition, type Baccalauréat]
\textbf{Barème indicatif (sur 10 points) :} respect de la consigne et de la longueur (2 pts) ; comparaison prix / qualité / garantie (3 pts) ; au moins deux comparatifs et un superlatif corrects (2 pts) ; vocabulaire du thème et connecteurs (2 pts) ; exactitude grammaticale et orthographe (1 pt).

\textbf{Proposition de rédaction modèle (environ 85 mots) :}

“If you want my advice, choose Phone B. Phone A is \textbf{cheaper than} Phone B — 45,000 FCFA instead of 60,000 — but Phone B is \textbf{better} in almost every way. Its camera is \textbf{far better than} Phone A's, and its two-year guarantee is \textbf{longer} and \textbf{more reassuring} than the one-year guarantee of Phone A. Moreover, you receive free earphones. In my opinion, Phone B is \textbf{the best value for money} and \textbf{the smartest choice} for a student who wants a phone that will last.”

\textbf{Points de vigilance :}
\begin{itemize}
\item Annoncer clairement son choix dès la première phrase (\eng{I advise you to choose…}, \eng{go for…}).
\item Comparer les trois critères demandés : prix (\eng{cheaper than}, \eng{more expensive than}), qualité (\eng{better camera}, \eng{higher quality}), garantie (\eng{a longer guarantee}).
\item Employer au moins un superlatif : \eng{the best value for money}, \eng{the most reliable}.
\item Éviter les fautes classiques : \eng{*more cheaper}, \eng{*the most best}, \eng{*advices} (nom non comptable : \eng{advice}).
\end{itemize}
\end{corrige}

\newpage

\coursec{Fiche bilan}

\begin{popfigure}
\begin{tikzpicture}[mindmap, grow cyclic, every node/.style={concept, align=center, font=\small},
root concept/.append style={concept color=popblue, font=\bfseries},
level 1/.append style={level distance=3.2cm, sibling angle=51},
level 2/.append style={level distance=1.9cm}]
\node[root concept, align=center]{Comparing Prices and Quality\\ Best Buys}
  child[concept color=poporange]{ node[align=center]{Key\\ Vocabulary} }
  child[concept color=popgreen]{ node[align=center]{Comparatives\\ Superlatives} }
  child[concept color=poppurple]{ node[align=center]{Listening\\ Strategies} }
  child[concept color=poppink]{ node[align=center]{Numbers\\ and Prices} }
  child[concept color=popgold]{ node[align=center]{Comprehension\\ Questions} }
  child[concept color=popblue!70]{ node[align=center]{Advising\\ the Best Buy} }
  child[concept color=popmuted]{ node[align=center]{Common\\ Traps} };
\end{tikzpicture}
\legende{Carte mentale de la leçon : écouter pour comparer prix et qualité.}
\end{popfigure}

\begin{bilan}
\textbf{1. Lexique clé (à connaître par cœur)}
\begin{itemize}
\item \eng{a bargain} : une bonne affaire ; \eng{a best buy} : le meilleur achat
\item \eng{cheap / expensive} : bon marché / cher ; \eng{affordable} : abordable
\item \eng{value for money} : le rapport qualité-prix ; \eng{good value} : avantageux
\item \eng{quality / durability} : qualité / durabilité ; \eng{reliable} : fiable
\item \eng{a discount / a reduction} : une remise ; \eng{on sale} : en solde
\item \eng{a receipt} : un reçu ; \eng{a refund} : un remboursement
\item \eng{a customer / a consumer} : un client / un consommateur
\item \eng{to compare prices} : comparer les prix ; \eng{to shop around} : comparer les offres
\item \eng{to save money} : économiser ; \eng{to waste money} : gaspiller de l'argent
\item \eng{community services} : services collectifs (eau, électricité, transport)
\end{itemize}

\textbf{2. Grammaire : comparatifs et superlatifs}
\begin{center}
\begin{tabularx}{\linewidth}{|l|X|X|}
\hline
\rowcolor{popblueL} Type & Comparatif & Superlatif \\
\hline
Adjectifs courts & \eng{cheaper than} & \eng{the cheapest} \\
\hline
Adjectifs longs & \eng{more expensive than} & \eng{the most expensive} \\
\hline
Irréguliers & \eng{good $\to$ better} ; \eng{bad $\to$ worse} & \eng{the best} ; \eng{the worst} \\
\hline
Égalité & \eng{as cheap as} & \eng{not as expensive as} \\
\hline
\end{tabularx}
\end{center}
Rappel : adjectif court terminé par une consonne + \emph{y} $\to$ \eng{happy, happier, the happiest}. Après un comparatif, toujours \eng{than} ; devant un superlatif, toujours \eng{the}.

\textbf{3. Méthodes express}
\begin{itemize}
\item Avant l'écoute : lire les questions, souligner les mots-clés, anticiper le thème.
\item Pendant l'écoute : noter les \cle{chiffres}, les \cle{prix}, les \cle{noms propres} et les comparatifs.
\item Pour un vrai/faux : repérer la phrase exacte du script qui justifie la réponse.
\item Pour conseiller un achat : \eng{I would recommend... because it is cheaper and more reliable.}
\end{itemize}
\end{bilan}

\begin{aretenir}[Checklist avant le Bac]
\begin{itemize}
\item[$\square$] Je connais le vocabulaire de la consommation : \eng{bargain, discount, receipt, refund, value for money}.
\item[$\square$] Je forme correctement les comparatifs : \eng{cheaper than}, \eng{more expensive than}.
\item[$\square$] Je forme correctement les superlatifs : \eng{the cheapest}, \eng{the most reliable}.
\item[$\square$] Je connais les formes irrégulières : \eng{good / better / the best} et \eng{bad / worse / the worst}.
\item[$\square$] J'utilise \eng{than} après un comparatif et jamais \emph{that} ou \emph{as}.
\item[$\square$] Je sais exprimer l'égalité : \eng{as good as}, \eng{not as expensive as}.
\item[$\square$] Avant l'écoute, je lis les questions et je souligne les mots-clés.
\item[$\square$] Pendant l'écoute, je note les chiffres, les prix et les noms propres.
\item[$\square$] Je justifie mes réponses vrai/faux par une information précise du texte.
\item[$\square$] Je distingue \eng{cheap} (bon marché) et \eng{affordable} (abordable) sans faux sens.
\item[$\square$] Je sais rédiger un conseil d'achat en deux ou trois phrases avec un comparatif.
\item[$\square$] Je prononce correctement les prix : \eng{5,000 CFA francs}, \eng{twenty-five thousand}.
\end{itemize}
\end{aretenir}

\findecours

\end{document}
